%Paper: alg-geom/9506023
%From: behrend@mpim-bonn.mpg.de (Kai Behrend)
%Date: Tue, 27 Jun 1995 16:32:00 +0200
%Date (revised): Thu, 26 Oct 1995 03:30:52 +0100

\documentstyle[amssymb,11pt]{amsart}
\title{Stacks of Stable Maps and Gromov-Witten Invariants}
\dedicatory{To Professor Friedrich Hirzebruch}
\date{October 18, 1995}
\author{K. Behrend}
\address{University of British Columbia\\
Mathematics Department\\
121--1984 Mathematics Road\\
Vancouver, B.C.\ Canada V6T 1Z2}
\email{behrend@@math.ubc.ca}
\author{Yu. Manin}
\address{Max-Planck-Institut f\"ur Mathematik\\
Gottfried-Claren-Str.\ 26\\
D--53225 Bonn}
\email{manin@@mpim-bonn.mpg.de}


% This is prepictex.tex  Version 1.1  9/10/87

% To use the PiCTeX macros under LaTeX, you first need to \input this
% file, then the main corpus of PiCTeX macros (pictex.tex), and then
% the file postpictex.tex. Do not \input the file latexpicobjs.tex.

\catcode`@=11 \catcode`!=11

% First of all, see if  \fiverm  is defined. If so do nothing;
% if not, let  \fiverm  take on meaning of LaTeX's  \fivrm.
\expandafter\ifx\csname fiverm\endcsname\relax
  \let\fiverm\fivrm
\fi

% Save meanings of LaTeX keywords that duplicate PiCTeX keywords
\let\!latexendpicture=\endpicture
\let\!latexframe=\frame
\let\!latexlinethickness=\linethickness
\let\!latexmultiput=\multiput
\let\!latexput=\put

% Redefine the LaTeX \@picture macro
\def\@picture(#1,#2)(#3,#4){%
  \@picht #2\unitlength
  \setbox\@picbox\hbox to #1\unitlength\bgroup
  \let\endpicture=\!latexendpicture
  \let\frame=\!latexframe
  \let\linethickness=\!latexlinethickness
  \let\multiput=\!latexmultiput
  \let\put=\!latexput
  \hskip -#3\unitlength \lower #4\unitlength \hbox\bgroup}

\catcode`@=12 \catcode`!=12

\newcommand{\fiverm}{}

% This is PiCTeX, Version 1.1   9/21/87

% CAVEAT: The PiCTeX manual often has a more lucid explanation
%   of any given topic than you will find in the internal documentation
%   of the macros.

% PiCTeX's commands can be classified into two groups: (1) public (or
%   external), and (2) private (or internal). The public macros are
%   discussed at length in the manual. The only discussion of the private
%   macros is the internal documentation. The private macros all have
%   names beginning with an exclamation point (!) of category code 11.
%   Since in normal usage "!" has category code 12, these macros can't
%   be accessed or modified by the general user.

% The macros are organized into thematically related groups. For example,
%   the macros dealing with dots & dashes are all in the DASHPATTERN group.
%   The table below shows which macros are in which groups. The table
%   covers all public macros, and many (but not all) of PiCTeX's upper level
%   private macros. Following the table, the various groups are listed
%   in the order in which they appear in the table.

% *********************** TABLE OF GROUPS OF MACROS **********************

% HACKS:  Utility macros
%    \PiC
%    \PiCTeX
%    \placevalueinpts
%    \!!loop
%    \!cfor
%    \!copylist
%    \!ecfor
%    \!etfor
%    \!getnext
%    \!getnextvalueof
%    \!ifempty
%    \!ifnextchar
%    \!leftappend
%    \!listaddon
%    \!loop
%    \!lop
%    \!mlap
%    \!not
%    \!removept
%    \!rightappend
%    \!tfor
%    \!vmlap
%    \!wlet

% ALLOCATION:  Allocates registers

% AREAS: Deals with plot areas
%    \axis
%    \grid
%    \invisibleaxes
%    \normalgraphs
%    \plotheading
%    \setplotarea
%    \visibleaxes

% ARROWS:  Draws arrows
%    \arrow
%    \betweenarrows

% BARS:  Draws bars
%    \putbar
%    \setbars

% BOXES:  Draws rectangles
%    \frame
%    \putrectangle
%    \rectangle
%    \shaderectangleson
%    \shaderectanglesoff

% CURVES:  Upper level plot commands
%    \hshade
%    \plot
%    \sethistograms
%    \setlinear
%    \setquadratic
%    \vshade

% DASHPATTERNS:  Sets up dash patterns
%    \findlength
%    \setdashes
%    \setdashesnear
%    \setdashpattern
%    \setdots
%    \setdotsnear
%    \setsolid
%    \!dashingoff
%    \!dashingon

% DIVISION:  Does long division of dimension registers
%    \Divide
%    \!divide

% ELLIPSES:  Draws ellipses and circles
%    \circulararc
%    \ellipticalarc

% RULES:  Draws rules, i.e., horizontal & vertical lines
%    \putrule
%    \!putdashedhline
%    \!putdashedvline
%    \!puthline
%    \!putsolidhline
%    \!putsolidvline
%    \!putvline

% LINEAR ARC:  Draws straight lines -- solid and dashed
%    \inboundscheckoff
%    \inboundscheckon
%    \!advancedashing
%    \!drawlinearsegment
%    \!initinboundscheck
%    \!linearsolid
%    \!lineardashed
%    \!ljoin
%    \!plotifinbounds
%    \!start

% LOGTEN:  Log_10 function
%    \!logten

% PICTURES:  Basic setups for PiCtures; \put commands
%    \accountingoff
%    \accountingon
%    \beginpicture
%    \endpicture
%    \endpicturesave
%    \lines
%    \multiput
%    \put
%    \setcoordinatemode
%    \setcoordinatesystem
%    \setdimensionmode
%    \stack
%    \Lines
%    \Xdistance
%    \Ydistance
%    \!dimenput
%    \!ifcoordmode
%    \!ifdimenmode
%    \!setcoordmode
%    \!setdimenmode
%    \!setputobject

% PLOTTING:  Things to do with plotting
%    \dontsavelinesandcurves
%    \replot
%    \savelinesandcurves
%    \setplotsymbol
%    \writesavefile
%    \!plot

% PYTHAGORAS:  Euclidean distance function
%    \placehypotenuse
%    \!Pythag

% QUADRATIC ARC:  Draws a quadratic arc
%    \!qjoin

% ROTATIONS:  Handles rotations
%    \startrotation
%    \stoprotation
%    \!rotateaboutpivot
%    \!rotateonly

% SHADING:  Handles shading
%    \setshadegrid
%    \setshadesymbol
%    \!lshade
%    \!qshade
%    \!starthshade
%    \!startvshade

% TICKS:  Draws ticks on graphs
%    \gridlines
%    \loggedticks
%    \nogridlines
%    \ticksin
%    \ticksout
%    \unloggesticks

% ***************** END OF TABLE OF GROUPS OF MACROS ********************


\catcode`!=11 %  ***** THIS MUST NEVER BE OMITTED
% *******************************
% *** HACKS  (Utility macros) ***
% *******************************

% ** User commands
% **   \PiC{P\kern-.12em\lower.5ex\hbox{I}\kern-.075emC}
% **   \PiCTeX{\PiC\kern-.11em\TeX}
% **   \placevalueinpts of <DIMENSION REGISTER> in {CONTROL SEQUENCE}

% ** Internal commands
% **   \!ifnextchar{CHARACTER}{TRUE ACTION}{FALSE ACTION}
% **   \!tfor NAME := LIST \do {BODY}
% **   \!etfor NAME:= LIST \do {BODY}
% **   \!cfor NAME := LIST \do {BODY}
% **   \!ecfor NAME:= LIST \do {BODY}
% **   \!ifempty{MACRO}{TRUE ACTION}{FALSE ACTION}
% **   \!getnext\\ITEMfrom\LIST
% **   \!getnextvalueof\DIMEN\from\LIST
% **   \!copylist\LISTMACRO_A\to\LISTMACRO_B
% **   \!wlet\CONTROL_SEQUENCE_A=\CONTROL_SEQUENCE_B
% **   \!listaddon ITEM LIST
% **   \!rightappendITEM\withCS\to\LISTMACRO
% **   \!leftappendITEM\withCS\to\LISTMACRO
% **   \!lop\LISTMACRO\to\ITEM
% **   \!loop ... repeat
% **   \!!loop ... repeat
% **   \!removept{DIMENSION REGISTER}{CONTROL SEQUENCE}
% **   \!mlap{...}
% **   \!vmlap{...}
% **   \!not{TEK if-CONDITION}

% ** First, here are the the PiCTeX logo, and the syllable PiC:
\def\PiC{P\kern-.12em\lower.5ex\hbox{I}\kern-.075emC}
\def\PiCTeX{\PiC\kern-.11em\TeX}

% ** The following macro expands to parameter #2 or parameter #3 according to
% ** whether the next non-blank character following the macro is or is not #1.
% ** Blanks following the macro are gobbled.
\def\!ifnextchar#1#2#3{%
  \let\!testchar=#1%
  \def\!first{#2}%
  \def\!second{#3}%
  \futurelet\!nextchar\!testnext}
\def\!testnext{%
  \ifx \!nextchar \!spacetoken
    \let\!next=\!skipspacetestagain
  \else
    \ifx \!nextchar \!testchar
      \let\!next=\!first
    \else
      \let\!next=\!second
    \fi
  \fi
  \!next}
\def\\{\!skipspacetestagain}
  \expandafter\def\\ {\futurelet\!nextchar\!testnext}
\def\\{\let\!spacetoken= } \\  %  ** set \spacetoken to a space token


% ** Borrow the "tfor" macro from Latex:
% **   \!tfor NAME := LIST \do {BODY}
% **   if, before expansion, LIST = T1 ... Tn,  where each  Ti  is a token
% **   or  {...},  then executes  BODY  n  times, with  NAME = Ti  on the
% **   i-th iteration.  Works for  n=0.
\def\!tfor#1:=#2\do#3{%
  \edef\!fortemp{#2}%
  \ifx\!fortemp\!empty
    \else
    \!tforloop#2\!nil\!nil\!!#1{#3}%
  \fi}
\def\!tforloop#1#2\!!#3#4{%
  \def#3{#1}%
  \ifx #3\!nnil
    \let\!nextwhile=\!fornoop
  \else
    #4\relax
    \let\!nextwhile=\!tforloop
  \fi
  \!nextwhile#2\!!#3{#4}}


% **   \!etfor NAME:= LIST\do {BODY}
% **   This is like \!cfor, but LIST is any balanced token list whose complete
% **     expansion has the form  T1 ... Tn
\def\!etfor#1:=#2\do#3{%
  \def\!!tfor{\!tfor#1:=}%
  \edef\!!!tfor{#2}%
  \expandafter\!!tfor\!!!tfor\do{#3}}


% **   modify the Latex \tfor (token-for) loop to a \cfor (comma-for) loop.
% **   \!cfor NAME := LIST \do {BODY}
% **     if, before expansion, LIST = a1,a2,...an, then executes  BODY n times,
% **     with  NAME = ai  on the i-th iteration.  Works for  n=0.
\def\!cfor#1:=#2\do#3{%
  \edef\!fortemp{#2}%
  \ifx\!fortemp\!empty
  \else
    \!cforloop#2,\!nil,\!nil\!!#1{#3}%
  \fi}
\def\!cforloop#1,#2\!!#3#4{%
  \def#3{#1}%
  \ifx #3\!nnil
    \let\!nextwhile=\!fornoop
  \else
    #4\relax
    \let\!nextwhile=\!cforloop
  \fi
  \!nextwhile#2\!!#3{#4}}


% **   \!ecfor NAME:= LIST\do {BODY}
% **   This is like \!cfor, but LIST is any balanced token list whose complete
% **     expansion has the form  a1,a2,...,an.
\def\!ecfor#1:=#2\do#3{%
  \def\!!cfor{\!cfor#1:=}%
  \edef\!!!cfor{#2}%
  \expandafter\!!cfor\!!!cfor\do{#3}}


\def\!empty{}
\def\!nnil{\!nil}
\def\!fornoop#1\!!#2#3{}


% **  \!ifempty{ARG}{TRUE ACTION}{FALSE ACTION}
\def\!ifempty#1#2#3{%
  \edef\!emptyarg{#1}%
  \ifx\!emptyarg\!empty
    #2%
  \else
    #3%
  \fi}

% **  \!getnext\\ITEMfrom\LIST
% **    \LIST has the form \\{item1}\\{item2}\\{item3}...\\{itemk}
% **    This routine sets \ITEM to item1, and cycles \LIST to
% **    \\{item2}\\{item3}...\\{itemk}\\{item1}
\def\!getnext#1\from#2{%
  \expandafter\!gnext#2\!#1#2}%
\def\!gnext\\#1#2\!#3#4{%
  \def#3{#1}%
  \def#4{#2\\{#1}}%
  \ignorespaces}


% ** \!getnextvalueof\DIMEN\from\LIST
% **   Similar to !getnext.
% **   \LIST has the form \\{dimen1}\\{dimen2}\\{dimen3} ...
% **   \DIMEN is a dimension register
% **   Works also for counts
%
\def\!getnextvalueof#1\from#2{%
  \expandafter\!gnextv#2\!#1#2}%
\def\!gnextv\\#1#2\!#3#4{%
  #3=#1%
  \def#4{#2\\{#1}}%
  \ignorespaces}


% ** \!copylist\LISTMACROA\to\LISTMACROB
% **   makes the replacement text of LISTMACRO B identical to that of
% **   list macro A.
\def\!copylist#1\to#2{%
  \expandafter\!!copylist#1\!#2}
\def\!!copylist#1\!#2{%
  \def#2{#1}\ignorespaces}


% **  \!wlet\CSA=\CSB
% **  lets control sequence \CSB = control sequence \CSA, and writes a
% **    message to that effect in the log file using plain TEK's \wlog
\def\!wlet#1=#2{%
  \let#1=#2
  \wlog{\string#1=\string#2}}

% ** \!listaddon ITEM LIST
% ** LIST <-- LIST \\ ITEM
\def\!listaddon#1#2{%
  \expandafter\!!listaddon#2\!{#1}#2}
\def\!!listaddon#1\!#2#3{%
  \def#3{#1\\#2}}

% **  \!rightappendITEM\to\LISTMACRO
% **    \LISTMACRO --> \LISTMACRO\\{ITEM}
%\def\!rightappend#1\to#2{\expandafter\!!rightappend#2\!{#1}#2}
%\def\!!rightappend#1\!#2#3{\def#3{#1\\{#2}}}


% **  \!rightappendITEM\withCS\to\LISTMACRO
% **    \LISTMACRO --> \LISTMACRO||CS||{ITEM}
\def\!rightappend#1\withCS#2\to#3{\expandafter\!!rightappend#3\!#2{#1}#3}
\def\!!rightappend#1\!#2#3#4{\def#4{#1#2{#3}}}


% **  \!leftappendITEM\withCS\to\LISTMACRO
% **    \LISTMACRO --> CS||{ITEM}||\LISTMACRO
\def\!leftappend#1\withCS#2\to#3{\expandafter\!!leftappend#3\!#2{#1}#3}
\def\!!leftappend#1\!#2#3#4{\def#4{#2{#3}#1}}


% **  \!lop\LISTMACRO\to\ITEM
% **    \\{item1}\\{item2}\\{item3} ... --> \\{item2}\\{item3} ...
% **    item1 --> \ITEM
\def\!lop#1\to#2{\expandafter\!!lop#1\!#1#2}
\def\!!lop\\#1#2\!#3#4{\def#4{#1}\def#3{#2}}


% **  \!placeNUMBER\of\LISTMACRO\in\ITEM
% **    the NUMBERth item of \LISTMACRO --> replacement text of \ITEM
%\def\!place#1\of#2\in#3{\def#3{\outofrange}%
%{\count0=#1\def\\##1{\advance\count0-1 \ifnum\count0=0 \gdef#3{##1}\fi}#2}}


% **  Following code converts a commalist to a list macro, with all items
% **    fully expanded.
%\!ecfor\item:=\commalist\do{\expandafter\!rightappend\item\to\list}


% ** \!loop ... repeat
% ** This is exactly like TEX's \loop ... repeat.  It can be used in nesting
% ** two loops, without puting the inner one inside a group.
\def\!loop#1\repeat{\def\!body{#1}\!iterate}
\def\!iterate{\!body\let\!next=\!iterate\else\let\!next=\relax\fi\!next}

% ** \!!loop ... repeat
% ** This is exactly like TEX's \loop ... repeat.  It can be used in nesting
% ** two loops, without puting the inner one inside a group.
\def\!!loop#1\repeat{\def\!!body{#1}\!!iterate}
\def\!!iterate{\!!body\let\!!next=\!!iterate\else\let\!!next=\relax\fi\!!next}
%  (\multiput uses \!!loop)

% ** \!removept{DIMENREG}{\CS}
% ** Defines the control sequence CS to be the value (in points) in the
% ** dimension register DIMENREG (but without the "pt" TEK usually adds)
% ** E.g., after  \dimen0=12.3pt \!removept\dimen0\A, \A expands to 12.3
\def\!removept#1#2{\edef#2{\expandafter\!!removePT\the#1}}
{\catcode`p=12 \catcode`t=12 \gdef\!!removePT#1pt{#1}}

% ** \pladevalueinpts of <DIMENSION REGISTER> in {CONTROL SEQUENCE}
\def\placevalueinpts of <#1> in #2 {%
  \!removept{#1}{#2}}

% ** \!mlap{...}  \!vmlap{...}
% ** Center  ...  in a box of width 0.
\def\!mlap#1{\hbox to 0pt{\hss#1\hss}}
\def\!vmlap#1{\vbox to 0pt{\vss#1\vss}}

% ** \!not{TEK if-CONDITION}
% ** By a TEK if-CONDITION is meant something like
% **     \ifnum\N<0,   or   \ifdim\A>\B
% ** \!not produces an if-condition which is false if the original condition
% ** is true, and true if the original condition is false.
\def\!not#1{%
  #1\relax
    \!switchfalse
  \else
    \!switchtrue
  \fi
  \if!switch
  \ignorespaces}


% *******************
% *** ALLOCATIONS ***
% *******************

% This section allocates all the registers PiCTeX uses. Following
% each allocation is a string of the form  ....N.D...L......... ;
% the various letters show which sections of PiCTeX make explicit
% reference to that register, according to the following code:

%   H Hacks
%   A Areas
%   W arroWs
%   B Bars
%   X boXes
%   C Curves
%   D Dashpattterns
%   V diVision
%   E Ellipses
%   U rUles
%   L Linear arc
%   G loGten
%   P Pictures
%   O plOtting
%   Y pYthagoras
%   Q Quadratic arc
%   R Rotations
%   S Shading
%   T Ticks

% Turn off messages from TeX's allocation macros
\let\!!!wlog=\wlog              % "\wlog" is defined in plain TeX
\def\wlog#1{}

\newdimen\headingtoplotskip     %.A.................
\newdimen\linethickness         %.A..X....U........T
\newdimen\longticklength        %.A................T
\newdimen\plotsymbolspacing     %......D...L....Q...
\newdimen\shortticklength       %.A................T
\newdimen\stackleading          %.A..........P......
\newdimen\tickstovaluesleading  %.A................T
\newdimen\totalarclength        %......D...L....Q...
\newdimen\valuestolabelleading  %.A.................

\newbox\!boxA                   %.AW...............T
\newbox\!boxB                   %..W................
\newbox\!picbox                 %............P......
\newbox\!plotsymbol             %..........L..O.....
\newbox\!putobject              %............PO...S.
\newbox\!shadesymbol            %.................S.

\newcount\!countA               %.A....D..UL....Q.ST
\newcount\!countB               %......D..U.....Q.ST
\newcount\!countC               %...............Q..T
\newcount\!countD               %...................
\newcount\!countE               %.............O....T
\newcount\!countF               %.............O....T
\newcount\!countG               %..................T
\newcount\!fiftypt              %.........U.........
\newcount\!intervalno           %..........L....Q...
\newcount\!npoints              %..........L........
\newcount\!nsegments            %.........U.........
\newcount\!ntemp                %............P......
\newcount\!parity               %.................S.
\newcount\!scalefactor          %..................T
\newcount\!tfs                  %.......V...........
\newcount\!tickcase             %..................T

\newdimen\!Xleft                %............P......
\newdimen\!Xright               %............P......
\newdimen\!Xsave                %.A................T
\newdimen\!Ybot                 %............P......
\newdimen\!Ysave                %.A................T
\newdimen\!Ytop                 %............P......
\newdimen\!angle                %........E..........
\newdimen\!arclength            %..W......UL....Q...
\newdimen\!areabloc             %.A........L........
\newdimen\!arealloc             %.A........L........
\newdimen\!arearloc             %.A........L........
\newdimen\!areatloc             %.A........L........
\newdimen\!bshrinkage           %.................S.
\newdimen\!checkbot             %..........L........
\newdimen\!checkleft            %..........L........
\newdimen\!checkright           %..........L........
\newdimen\!checktop             %..........L........
\newdimen\!dimenA               %.AW.X.DVEUL..OYQRST
\newdimen\!dimenB               %....X.DVEU...O.QRS.
\newdimen\!dimenC               %..W.X.DVEU......RS.
\newdimen\!dimenD               %..W.X.DVEU....Y.RS.
\newdimen\!dimenE               %..W........G..YQ.S.
\newdimen\!dimenF               %...........G..YQ.S.
\newdimen\!dimenG               %...........G..YQ.S.
\newdimen\!dimenH               %...........G..Y..S.
\newdimen\!dimenI               %...BX.........Y....
\newdimen\!distacross           %..........L....Q...
\newdimen\!downlength           %..........L........
\newdimen\!dp                   %.A..X.......P....S.
\newdimen\!dshade               %.................S.
\newdimen\!dxpos                %..W......U..P....S.
\newdimen\!dxprime              %...............Q...
\newdimen\!dypos                %..WB.....U..P......
\newdimen\!dyprime              %...............Q...
\newdimen\!ht                   %.A..X.......P....S.
\newdimen\!leaderlength         %......D..U.........
\newdimen\!lshrinkage           %.................S.
\newdimen\!midarclength         %...............Q...
\newdimen\!offset               %.A................T
\newdimen\!plotheadingoffset    %.A.................
\newdimen\!plotsymbolxshift     %..........L..O.....
\newdimen\!plotsymbolyshift     %..........L..O.....
\newdimen\!plotxorigin          %..........L..O.....
\newdimen\!plotyorigin          %..........L..O.....
\newdimen\!rootten              %...........G.......
\newdimen\!rshrinkage           %.................S.
\newdimen\!shadesymbolxshift    %.................S.
\newdimen\!shadesymbolyshift    %.................S.
\newdimen\!tenAa                %...........G.......
\newdimen\!tenAc                %...........G.......
\newdimen\!tenAe                %...........G.......
\newdimen\!tshrinkage           %.................S.
\newdimen\!uplength             %..........L........
\newdimen\!wd                   %....X.......P....S.
\newdimen\!wmax                 %...............Q...
\newdimen\!wmin                 %...............Q...
\newdimen\!xB                   %...............Q...
\newdimen\!xC                   %...............Q...
\newdimen\!xE                   %..W.....E.L....Q.S.
\newdimen\!xM                   %..W.....E......Q.S.
\newdimen\!xS                   %..W.....E.L....Q.S.
\newdimen\!xaxislength          %.A................T
\newdimen\!xdiff                %..........L........
\newdimen\!xleft                %............P......
\newdimen\!xloc                 %..WB.....U.......S.
\newdimen\!xorigin              %.A........L.P....S.
\newdimen\!xpivot               %................R..
\newdimen\!xpos                 %..........L.P..Q.ST
\newdimen\!xprime               %...............Q...
\newdimen\!xright               %............P......
\newdimen\!xshade               %.................S.
\newdimen\!xshift               %..W.........PO...S.
\newdimen\!xtemp                %............P......
\newdimen\!xunit                %.AWBX...EUL.P..QRS.
\newdimen\!xxE                  %........E..........
\newdimen\!xxM                  %........E..........
\newdimen\!xxS                  %........E..........
\newdimen\!xxloc                %..WB....EU.........
\newdimen\!yB                   %...............Q...
\newdimen\!yC                   %...............Q...
\newdimen\!yE                   %..W.....E.L....Q...
\newdimen\!yM                   %..W.....E......Q...
\newdimen\!yS                   %..W.....E.L....Q...
\newdimen\!yaxislength          %.A................T
\newdimen\!ybot                 %............P......
\newdimen\!ydiff                %..........L........
\newdimen\!yloc                 %..WB.....U.......S.
\newdimen\!yorigin              %.A........L.P....S.
\newdimen\!ypivot               %................R..
\newdimen\!ypos                 %..........L.P..Q.ST
\newdimen\!yprime               %...............Q...
\newdimen\!yshade               %.................S.
\newdimen\!yshift               %..W.........PO...S.
\newdimen\!ytemp                %............P......
\newdimen\!ytop                 %............P......
\newdimen\!yunit                %.AWBX...EUL.P..QRS.
\newdimen\!yyE                  %........E..........
\newdimen\!yyM                  %........E..........
\newdimen\!yyS                  %........E..........
\newdimen\!yyloc                %..WB....EU.........
\newdimen\!zpt                  %.AWBX.DVEULGP.YQ.ST

\newif\if!axisvisible           %.A.................
\newif\if!gridlinestoo          %..................T
\newif\if!keepPO                %...................
\newif\if!placeaxislabel        %.A.................
\newif\if!switch                %H..................
\newif\if!xswitch               %.A................T

\newtoks\!axisLaBeL             %.A.................
\newtoks\!keywordtoks           %.A.................

\newwrite\!replotfile           %.............O.....

\newhelp\!keywordhelp{The keyword mentioned in the error message in unknown.
Replace NEW KEYWORD in the indicated response by the keyword that
should have been specified.}    %.A.................

% The following commands assign alternate names to some of the
% above registers.  "\!wlet"  is defined in  Hacks.
\!wlet\!!origin=\!xM                   %.A................T
\!wlet\!!unit=\!uplength               %.A................T
\!wlet\!Lresiduallength=\!dimenG       %.........U.........
\!wlet\!Rresiduallength=\!dimenF       %.........U.........
\!wlet\!axisLength=\!distacross        %.A................T
\!wlet\!axisend=\!ydiff                %.A................T
\!wlet\!axisstart=\!xdiff              %.A................T
\!wlet\!axisxlevel=\!arclength         %.A................T
\!wlet\!axisylevel=\!downlength        %.A................T
\!wlet\!beta=\!dimenE                  %...............Q...
\!wlet\!gamma=\!dimenF                 %...............Q...
\!wlet\!shadexorigin=\!plotxorigin     %.................S.
\!wlet\!shadeyorigin=\!plotyorigin     %.................S.
\!wlet\!ticklength=\!xS                %..................T
\!wlet\!ticklocation=\!xE              %..................T
\!wlet\!ticklocationincr=\!yE          %..................T
\!wlet\!tickwidth=\!yS                 %..................T
\!wlet\!totalleaderlength=\!dimenE     %.........U.........
\!wlet\!xone=\!xprime                  %....X..............
\!wlet\!xtwo=\!dxprime                 %....X..............
\!wlet\!ySsave=\!yM                    %...................
\!wlet\!ybB=\!yB                       %.................S.
\!wlet\!ybC=\!yC                       %.................S.
\!wlet\!ybE=\!yE                       %.................S.
\!wlet\!ybM=\!yM                       %.................S.
\!wlet\!ybS=\!yS                       %.................S.
\!wlet\!ybpos=\!yyloc                  %.................S.
\!wlet\!yone=\!yprime                  %....X..............
\!wlet\!ytB=\!xB                       %.................S.
\!wlet\!ytC=\!xC                       %.................S.
\!wlet\!ytE=\!downlength               %.................S.
\!wlet\!ytM=\!arclength                %.................S.
\!wlet\!ytS=\!distacross               %.................S.
\!wlet\!ytpos=\!xxloc                  %.................S.
\!wlet\!ytwo=\!dyprime                 %....X..............


% Initial values for registers
\!zpt=0pt                              % static
\!xunit=1pt
\!yunit=1pt
\!arearloc=\!xunit
\!areatloc=\!yunit
\!dshade=5pt
\!leaderlength=24in
\!tfs=256                              % static
\!wmax=5.3pt                           % static
\!wmin=2.7pt                           % static
\!xaxislength=\!xunit
\!xpivot=\!zpt
\!yaxislength=\!yunit
\!ypivot=\!zpt
\plotsymbolspacing=.4pt
  \!dimenA=50pt \!fiftypt=\!dimenA     % static

\!rootten=3.162278pt                   % static
\!tenAa=8.690286pt                     % static  (A5)
\!tenAc=2.773839pt                     % static  (A3)
\!tenAe=2.543275pt                     % static  (A1)

% Initial values for control sequences
\def\!cosrotationangle{1}      %................R..
\def\!sinrotationangle{0}      %................R..
\def\!xpivotcoord{0}           %................R..
\def\!xref{0}                  %............P......
\def\!xshadesave{0}            %.................S.
\def\!ypivotcoord{0}           %................R..
\def\!yref{0}                  %............P......
\def\!yshadesave{0}            %.................S.
\def\!zero{0}                  %..................T

% Reset TeX to report allocations
\let\wlog=\!!!wlog
%  *************************************
%  ***  AREAS: Deals with plot areas ***
%  *************************************
%
%  ** User commands
%  **   \setplotarea x from LEFT XCOORD to RIGTH XCOORD, y from BOTTOM YCOORD
%  **      to TOP YCOORD
%  **   \axis BOTTOM-LEFT-TOP-RIGHT  [SHIFTEDTO xy=COORD] [VISIBLE-INVISIBLE]
%  **      [LABEL {label}] [TICKS] /
%  **   \visibleaxes
%  **   \invisibleaxes
%  **   \plotheading {HEADING}
%  **   \grid {# of columns} {# of rows}
%  **   \normalgraphs

%  **  \normalgraphs
%  **    Sets defaults for graph setup. See Subsection 3.4 of manual.
\def\normalgraphs{%
  \longticklength=.4\baselineskip
  \shortticklength=.25\baselineskip
  \tickstovaluesleading=.25\baselineskip
  \valuestolabelleading=.8\baselineskip
  \linethickness=.4pt
  \stackleading=.17\baselineskip
  \headingtoplotskip=1.5\baselineskip
  \visibleaxes
  \ticksout
  \nogridlines
  \unloggedticks}
%
% **  \setplotarea x from LEFT XCOORD to RIGTH XCOORD, y from BOTTOM YCOORD
% **    to TOP YCOORD
% **  Reserves space in PICBOX for a rectangular box with the indicated
% **   coordinates.  Must be specified before calls to  \axis,
% **   \grid, \plotheading.
% **  See Subsection 3.1 of the manual.
\def\setplotarea x from #1 to #2, y from #3 to #4 {%
  \!arealloc=\!M{#1}\!xunit \advance \!arealloc -\!xorigin
  \!areabloc=\!M{#3}\!yunit \advance \!areabloc -\!yorigin
  \!arearloc=\!M{#2}\!xunit \advance \!arearloc -\!xorigin
  \!areatloc=\!M{#4}\!yunit \advance \!areatloc -\!yorigin
  \!initinboundscheck
  \!xaxislength=\!arearloc  \advance\!xaxislength -\!arealloc
  \!yaxislength=\!areatloc  \advance\!yaxislength -\!areabloc
  \!plotheadingoffset=\!zpt
  \!dimenput {{\setbox0=\hbox{}\wd0=\!xaxislength\ht0=\!yaxislength\box0}}
     [bl] (\!arealloc,\!areabloc)}
%
% ** \visibleaxes, \invisibleaxes
% ** Switches for setting visibility of subsequent axes.
% ** See Subsection 3.2 of the manual.
\def\visibleaxes{%
  \def\!axisvisibility{\!axisvisibletrue}}
\def\invisibleaxes{%
  \def\!axisvisibility{\!axisvisiblefalse}}
%
% ** The next few macros enable the user to fix up an erroneous keyword
% **   in the \axis command.
%  \newhelp is in ALLOCATIONS
%  \newhelp\!keywordhelp{The keyword mentioned in the error message in unknown.
%  Replace NEW KEYWORD in the indicated response by the keyword that
%  should have been specified.}

\def\!fixkeyword#1{%
  \errhelp=\!keywordhelp
  \errmessage{Unrecognized keyword `#1': \the\!keywordtoks{NEW KEYWORD}'}}

%  \newtoks\!keywordtoks    In ALLOCATIONS.
\!keywordtoks={enter `i\fixkeyword}

\def\fixkeyword#1{%
  \!nextkeyword#1 }

% ** \axis BOTTOM-LEFT-TOP-RIGHT  [SHIFTEDTO xy=COORD] [VISIBLE-INVISIBLE]
% **   [LABEL {label}] [TICKS] /
% ** Exactly one of the keywords BOTTOM, LEFT, TOP, RIGHT must be
% ** specified. Axis is drawn along the indicated edge of the current
% ** plot area, shifted if the SHIFTEDTO option is used, visible or
% ** invisible according the selected option, with an optional LABEL,
% ** and optional TICKS (see ticks.tex for the options avialabel with
% ** TICKS). The TICKS option must be the last one specified. The \axis
% ** MUST be terminated with a / followed by a space.
% ** See Subsection 3.2 of the manual for more information.

% ** The various options of the \axis command are processed by the
% ** \!nextkeyword macro defined below. For example,
% ** `\!nextkeyword shiftedto ' expands to `\!axisshiftedto'.
\def\axis {%
  \def\!nextkeyword##1 {%
    \expandafter\ifx\csname !axis##1\endcsname \relax
      \def\!next{\!fixkeyword{##1}}%
    \else
      \def\!next{\csname !axis##1\endcsname}%
    \fi
    \!next}%
  \!offset=\!zpt
  \!axisvisibility
  \!placeaxislabelfalse
  \!nextkeyword}

% ** This and the various macros that follow handle the keyword
% ** specifications on the \axis command
% ** See Subsection 3.2 of the manual.
\def\!axisbottom{%
  \!axisylevel=\!areabloc
  \def\!tickxsign{0}%
  \def\!tickysign{-}%
  \def\!axissetup{\!axisxsetup}%
  \def\!axislabeltbrl{t}%
  \!nextkeyword}

\def\!axistop{%
  \!axisylevel=\!areatloc
  \def\!tickxsign{0}%
  \def\!tickysign{+}%
  \def\!axissetup{\!axisxsetup}%
  \def\!axislabeltbrl{b}%
  \!nextkeyword}

\def\!axisleft{%
  \!axisxlevel=\!arealloc
  \def\!tickxsign{-}%
  \def\!tickysign{0}%
  \def\!axissetup{\!axisysetup}%
  \def\!axislabeltbrl{r}%
  \!nextkeyword}

\def\!axisright{%
  \!axisxlevel=\!arearloc
  \def\!tickxsign{+}%
  \def\!tickysign{0}%
  \def\!axissetup{\!axisysetup}%
  \def\!axislabeltbrl{l}%
  \!nextkeyword}

\def\!axisshiftedto#1=#2 {%
  \if 0\!tickxsign
    \!axisylevel=\!M{#2}\!yunit
    \advance\!axisylevel -\!yorigin
  \else
    \!axisxlevel=\!M{#2}\!xunit
    \advance\!axisxlevel -\!xorigin
  \fi
  \!nextkeyword}

\def\!axisvisible{%
  \!axisvisibletrue
  \!nextkeyword}

\def\!axisinvisible{%
  \!axisvisiblefalse
  \!nextkeyword}

\def\!axislabel#1 {%
  \!axisLaBeL={#1}%
  \!placeaxislabeltrue
  \!nextkeyword}

\expandafter\def\csname !axis/\endcsname{%
  \!axissetup % This could done already by "ticks"; if so, now \relax
  \if!placeaxislabel
    \!placeaxislabel
  \fi
  \if +\!tickysign %                 ** (A "top" axis)
    \!dimenA=\!axisylevel
    \advance\!dimenA \!offset %      ** dimA = top of the axis structure
    \advance\!dimenA -\!areatloc %   ** dimA = excess over the plot area
    \ifdim \!dimenA>\!plotheadingoffset
      \!plotheadingoffset=\!dimenA % ** Greatest excess over the plot area
    \fi
  \fi}

% ** \grid {c} {r}
% ** Partitions the plot area into c columns and r rows; see Subsection 3.3
% ** of the manual.
% ** (Other grid patterns can be drawn with the TICKS option of the \axis
% ** command.
\def\grid #1 #2 {%
  \!countA=#1\advance\!countA 1
  \axis bottom invisible ticks length <\!zpt> andacross quantity {\!countA} /
  \!countA=#2\advance\!countA 1
  \axis left   invisible ticks length <\!zpt> andacross quantity {\!countA} / }

% ** \plotheading{HEADING}
% ** Places HEADING centered above the top of the plotarea (and above
% ** any top axis ticks marks, tick labels, and axis label); see
% ** Subsection 3.3 of the manual.
\def\plotheading#1 {%
  \advance\!plotheadingoffset \headingtoplotskip
  \!dimenput {#1} [B] <.5\!xaxislength,\!plotheadingoffset>
    (\!arealloc,\!areatloc)}

% ** From here on, the routines are internal.
\def\!axisxsetup{%
  \!axisxlevel=\!arealloc
  \!axisstart=\!arealloc
  \!axisend=\!arearloc
  \!axisLength=\!xaxislength
  \!!origin=\!xorigin
  \!!unit=\!xunit
  \!xswitchtrue
  \if!axisvisible
    \!makeaxis
  \fi}

\def\!axisysetup{%
  \!axisylevel=\!areabloc
  \!axisstart=\!areabloc
  \!axisend=\!areatloc
  \!axisLength=\!yaxislength
  \!!origin=\!yorigin
  \!!unit=\!yunit
  \!xswitchfalse
  \if!axisvisible
    \!makeaxis
  \fi}

\def\!makeaxis{%
  \setbox\!boxA=\hbox{% (Make a pseudo-y[x] tick for an x[y]-axis)
    \beginpicture
      \!setdimenmode
      \setcoordinatesystem point at {\!zpt} {\!zpt}
      \putrule from {\!zpt} {\!zpt} to
        {\!tickysign\!tickysign\!axisLength}
        {\!tickxsign\!tickxsign\!axisLength}
    \endpicturesave <\!Xsave,\!Ysave>}%
    \wd\!boxA=\!zpt
    \!placetick\!axisstart}

\def\!placeaxislabel{%
  \advance\!offset \valuestolabelleading
  \if!xswitch
    \!dimenput {\the\!axisLaBeL} [\!axislabeltbrl]
      <.5\!axisLength,\!tickysign\!offset> (\!axisxlevel,\!axisylevel)
    \advance\!offset \!dp  % ** advance offset by the "tallness"
    \advance\!offset \!ht  % ** of the label
  \else
    \!dimenput {\the\!axisLaBeL} [\!axislabeltbrl]
      <\!tickxsign\!offset,.5\!axisLength> (\!axisxlevel,\!axisylevel)
  \fi
  \!axisLaBeL={}}


% *******************************
% *** ARROWS  (Draws arrows)  ***
% *******************************
%
% ** User commands
% **  \arrow <ARROW HEAD LENGTH> [MID FRACTION, BASE FRACTION]
% **    [<XSHIFT,YSHIFT>] from XFROM YFROM to XTO YTO
% **  \betweenarrows {TEXT} [orientation & shift] from XFROM YFROM to XTO YTO

% ** \arrow <ARROW HEAD LENGTH> [MID FRACTION, BASE FRACTION]
% **    [<XSHIFT,YSHIFT>] from XFROM YFROM to XTO YTO
% ** Draws an arrow from (XFROM,YFROM) to (XTO,YTO).  The arrow head
% ** is constructed two quadratic arcs, which extend back a distance
% ** ARROW HEAD LENGTH (a dimension) on both sides of the arrow shaft.
% ** All the way back the arcs are a distance BASE FRACTION*ARROW HEAD
% ** LENGTH apart, while half-way back they are a distance MID FRACTION*
% ** ARROW HEAD LENGTH apart. <XSHIFT,YSHIFT> is optional, and has
% ** its usual interpreation. See Subsection 5.4 of the manual.

\def\arrow <#1> [#2,#3]{%
  \!ifnextchar<{\!arrow{#1}{#2}{#3}}{\!arrow{#1}{#2}{#3}<\!zpt,\!zpt> }}

\def\!arrow#1#2#3<#4,#5> from #6 #7 to #8 #9 {%
%
% ** convert to dimensions
  \!xloc=\!M{#8}\!xunit
  \!yloc=\!M{#9}\!yunit
  \!dxpos=\!xloc  \!dimenA=\!M{#6}\!xunit  \advance \!dxpos -\!dimenA
  \!dypos=\!yloc  \!dimenA=\!M{#7}\!yunit  \advance \!dypos -\!dimenA
  \let\!MAH=\!M%                         ** save current c/d mode
  \!setdimenmode%                        ** go into dimension mode
%
  \!xshift=#4\relax  \!yshift=#5\relax%  ** pick up shift
  \!reverserotateonly\!xshift\!yshift%   ** back rotate shift
  \advance\!xshift\!xloc  \advance\!yshift\!yloc
%
% **  draw shaft of arrow
  \!xS=-\!dxpos  \advance\!xS\!xshift
  \!yS=-\!dypos  \advance\!yS\!yshift
  \!start (\!xS,\!yS)
  \!ljoin (\!xshift,\!yshift)
%
% ** find 32*cosine and 32*sine of angle of rotation
  \!Pythag\!dxpos\!dypos\!arclength
  \!divide\!dxpos\!arclength\!dxpos
  \!dxpos=32\!dxpos  \!removept\!dxpos\!!cos
  \!divide\!dypos\!arclength\!dypos
  \!dypos=32\!dypos  \!removept\!dypos\!!sin
%
% ** construct arrowhead
  \!halfhead{#1}{#2}{#3}%                ** draw half of arrow head
  \!halfhead{#1}{-#2}{-#3}%              ** draw other half
%
  \let\!M=\!MAH%                         ** restore old c/d mode
  \ignorespaces}
%
% ** draw half of arrow head
  \def\!halfhead#1#2#3{%
    \!dimenC=-#1%
    \divide \!dimenC 2 %                 ** half way back
    \!dimenD=#2\!dimenC%                 ** half the mid width
    \!rotate(\!dimenC,\!dimenD)by(\!!cos,\!!sin)to(\!xM,\!yM)
    \!dimenC=-#1%                        ** all the way back
    \!dimenD=#3\!dimenC
    \!dimenD=.5\!dimenD%                 ** half the full width
    \!rotate(\!dimenC,\!dimenD)by(\!!cos,\!!sin)to(\!xE,\!yE)
    \!start (\!xshift,\!yshift)
    \advance\!xM\!xshift  \advance\!yM\!yshift
    \advance\!xE\!xshift  \advance\!yE\!yshift
    \!qjoin (\!xM,\!yM) (\!xE,\!yE)
    \ignorespaces}


% ** \betweenarrows {TEXT} [orientation & shift] from XFROM YFROM to XTO YTO
% **   Makes things like <--- text --->, using arrow heads from TeX's fonts.
% **   See Subsection 5.4 of the manual.
\def\betweenarrows #1#2 from #3 #4 to #5 #6 {%
  \!xloc=\!M{#3}\!xunit  \!xxloc=\!M{#5}\!xunit%
  \!yloc=\!M{#4}\!yunit  \!yyloc=\!M{#6}\!yunit%
  \!dxpos=\!xxloc  \advance\!dxpos by -\!xloc
  \!dypos=\!yyloc  \advance\!dypos by -\!yloc
  \advance\!xloc .5\!dxpos
  \advance\!yloc .5\!dypos
%
  \let\!MBA=\!M%           ** save current coord\dimen mode
  \!setdimenmode%          ** express locations in dimens
  \ifdim\!dypos=\!zpt
    \ifdim\!dxpos<\!zpt \!dxpos=-\!dxpos \fi
    \put {\!lrarrows{\!dxpos}{#1}}#2{} at {\!xloc} {\!yloc}
  \else
    \ifdim\!dxpos=\!zpt
      \ifdim\!dypos<\!zpt \!dypos=-\!zpt \fi
      \put {\!udarrows{\!dypos}{#1}}#2{} at {\!xloc} {\!yloc}
    \fi
  \fi
  \let\!M=\!MBA%           ** restore previous c/d mode
  \ignorespaces}

% ** Subroutine for left-right between arrows
\def\!lrarrows#1#2{% #1=width, #2=text
  {\setbox\!boxA=\hbox{$\mkern-2mu\mathord-\mkern-2mu$}%
   \setbox\!boxB=\hbox{$\leftarrow$}\!dimenE=\ht\!boxB
   \setbox\!boxB=\hbox{}\ht\!boxB=2\!dimenE
   \hbox to #1{$\mathord\leftarrow\mkern-6mu
     \cleaders\copy\!boxA\hfil
     \mkern-6mu\mathord-$%
     \kern.4em $\vcenter{\box\!boxB}$$\vcenter{\hbox{#2}}$\kern.4em
     $\mathord-\mkern-6mu
     \cleaders\copy\!boxA\hfil
     \mkern-6mu\mathord\rightarrow$}}}

% ** Subroutine for up-down between arrows
\def\!udarrows#1#2{% #1=width, #2=text
  {\setbox\!boxB=\hbox{#2}%
   \setbox\!boxA=\hbox to \wd\!boxB{\hss$\vert$\hss}%
   \!dimenE=\ht\!boxA \advance\!dimenE \dp\!boxA \divide\!dimenE 2
   \vbox to #1{\offinterlineskip
      \vskip .05556\!dimenE
      \hbox to \wd\!boxB{\hss$\mkern.4mu\uparrow$\hss}\vskip-\!dimenE
      \cleaders\copy\!boxA\vfil
      \vskip-\!dimenE\copy\!boxA
      \vskip\!dimenE\copy\!boxB\vskip.4em
      \copy\!boxA\vskip-\!dimenE
      \cleaders\copy\!boxA\vfil
      \vskip-\!dimenE \hbox to \wd\!boxB{\hss$\mkern.4mu\downarrow$\hss}
      \vskip .05556\!dimenE}}}


% ***************************
% *** BARS  (Draws bars)  ***
% ***************************
%
% ** User commands:
% ** \putbar [<XSHIFT,YSHIFT>] breadth <BREADTH> from XSTART YSTART
% **   to XEND YEND
% ** \setbars [<XSHIFT,YSHIFT>] breadth <BREADTH> baseline at XY = COORD
% **   [baselabels ([B_ORIENTATION_x,B_ORIENTATION_y] <B_XSHIFT,B_YSHIFT>)]
% **   [endlabels  ([E_ORIENTATION_x,E_ORIENTATION_y] <E_XSHIFT,E_YSHIFT>)]


% ** \putbar [<XSHIFT,YSHIFT>] breadth <BREADTH> from XSTART YSTART
% **   to XEND YEND
% ** Either XSTART=XEND or YSTART=YEND. Draws a rectangle between
% **   (XSTART,YSTART) & (XEND,YEND). The "depth" of the rectangle
% **   is determined by those two plot positions; its other
% **   dimension "breadth" is specified by the dimension BREADTH.
% ** See Subsection 4.2 of the manual.
\def\putbar#1breadth <#2> from #3 #4 to #5 #6 {%
  \!xloc=\!M{#3}\!xunit  \!xxloc=\!M{#5}\!xunit%
  \!yloc=\!M{#4}\!yunit  \!yyloc=\!M{#6}\!yunit%
  \!dypos=\!yyloc  \advance\!dypos by -\!yloc
  \!dimenI=#2
%
  \ifdim \!dimenI=\!zpt %            ** If 0 breadth
    \putrule#1from {#3} {#4} to {#5} {#6} % ** Then draw line
  \else %                            ** Else, put in a rectangle
    \let\!MBar=\!M%                  ** save current c/d mode
    \!setdimenmode %                 ** go into dimension mode
    \divide\!dimenI 2
    \ifdim \!dypos=\!zpt
      \advance \!yloc -\!dimenI %    ** Equal y coordinates
      \advance \!yyloc \!dimenI
    \else
      \advance \!xloc -\!dimenI %    ** Equal x coordinates
      \advance \!xxloc \!dimenI
    \fi
    \putrectangle#1corners at {\!xloc} {\!yloc} and {\!xxloc} {\!yyloc}
    \let\!M=\!MBar %                 ** restore c/d mode
  \fi
  \ignorespaces}


% ** \setbars [<XSHIFT,YSHIFT>] breadth <BREADTH> baseline at XY = COORD
% **   [baselabels ([B_ORIENTATION_x,B_ORIENTATION_y] <B_XSHIFT,B_YSHIFT>)]
% **   [endlabels  ([E_ORIENTATION_x,E_ORIENTATION_y] <E_XSHIFT,E_YSHIFT>)]
% ** This command puts PiCTeX into the bar graph drawing mode described
% **   in Subsection 4.4 of the manual.
\def\setbars#1breadth <#2> baseline at #3 = #4 {%
  \edef\!barshift{#1}%
  \edef\!barbreadth{#2}%
  \edef\!barorientation{#3}%
  \edef\!barbaseline{#4}%
  \def\!bardobaselabel{\!bardoendlabel}%
  \def\!bardoendlabel{\!barfinish}%
  \let\!drawcurve=\!barcurve
  \!setbars}
\def\!setbars{%
  \futurelet\!nextchar\!!setbars}
\def\!!setbars{%
  \if b\!nextchar
    \def\!!!setbars{\!setbarsbget}%
  \else
    \if e\!nextchar
      \def\!!!setbars{\!setbarseget}%
    \else
      \def\!!!setbars{\relax}%
    \fi
  \fi
  \!!!setbars}
\def\!setbarsbget baselabels (#1) {%
  \def\!barbaselabelorientation{#1}%
  \def\!bardobaselabel{\!!bardobaselabel}%
  \!setbars}
\def\!setbarseget endlabels (#1) {%
  \edef\!barendlabelorientation{#1}%
  \def\!bardoendlabel{\!!bardoendlabel}%
  \!setbars}

% ** \!barcurve
% ** Draws a bargraph with preset values of barshift, barbreadth,
% ** barorientation (x or y) and barbaseline (coordinate)
\def\!barcurve #1 #2 {%
  \if y\!barorientation
    \def\!basexarg{#1}%
    \def\!baseyarg{\!barbaseline}%
  \else
    \def\!basexarg{\!barbaseline}%
    \def\!baseyarg{#2}%
  \fi
  \expandafter\putbar\!barshift breadth <\!barbreadth> from {\!basexarg}
    {\!baseyarg} to {#1} {#2}
  \def\!endxarg{#1}%
  \def\!endyarg{#2}%
  \!bardobaselabel}

\def\!!bardobaselabel "#1" {%
  \put {#1}\!barbaselabelorientation{} at {\!basexarg} {\!baseyarg}
  \!bardoendlabel}

\def\!!bardoendlabel "#1" {%
  \put {#1}\!barendlabelorientation{} at {\!endxarg} {\!endyarg}
  \!barfinish}

\def\!barfinish{%
  \!ifnextchar/{\!finish}{\!barcurve}}


% ********************************
% *** BOXES (Draws rectangles) ***
% ********************************
%
% ** User commands:
% **   \putrectangle [<XSHIFT,YSHIFT>] corners at  XCOORD1 YCOORD1
% **     and  XCOORD2 YCOORD2
% **   \shaderectangleson
% **   \shaderectanglesoff
% **   \frame [<SEPARATION>] {TEXT}
% **   \rectangle <WIDTH> <HEIGHT>
%
%
% **  \putrectangle [<XSHIFT,YSHIFT>] corners at XCOORD1 YCOORD1
% **    and  XCOORD2 YCOORD2
% **  Draws a rectangle with corners at (X1,Y1), (X2,Y1), (X1,Y2), (X2,Y2)
% **  Lines have thickness \linethickness, and overlap at the corners.
% **  The optional field  <XSHIFT,YSHIFT>  functions as with a \put command.
% **  See Subsection 4.2 of the manual.
\def\putrectangle{%
  \!ifnextchar<{\!putrectangle}{\!putrectangle<\!zpt,\!zpt> }}
\def\!putrectangle<#1,#2> corners at #3 #4 and #5 #6 {%
%
% ** get locations
  \!xone=\!M{#3}\!xunit  \!xtwo=\!M{#5}\!xunit%
  \!yone=\!M{#4}\!yunit  \!ytwo=\!M{#6}\!yunit%
  \ifdim \!xtwo<\!xone
    \!dimenI=\!xone  \!xone=\!xtwo  \!xtwo=\!dimenI
  \fi
  \ifdim \!ytwo<\!yone
    \!dimenI=\!yone  \!yone=\!ytwo  \!ytwo=\!dimenI
  \fi
  \!dimenI=#1\relax  \advance\!xone\!dimenI  \advance\!xtwo\!dimenI
  \!dimenI=#2\relax  \advance\!yone\!dimenI  \advance\!ytwo\!dimenI
  \let\!MRect=\!M%                  ** save current coord/dimen mode
  \!setdimenmode
%
% ** shade rectangle if appropriate
  \!shaderectangle
%
% ** draw horizontal edges
  \!dimenI=.5\linethickness
  \advance \!xone  -\!dimenI%       ** adjust x-location to overlap corners
  \advance \!xtwo   \!dimenI%       ** ditto
  \putrule from {\!xone} {\!yone} to {\!xtwo} {\!yone}
  \putrule from {\!xone} {\!ytwo} to {\!xtwo} {\!ytwo}
%
% ** draw vertical edges
  \advance \!xone   \!dimenI%       ** restore original x-values
  \advance \!xtwo  -\!dimenI%
  \advance \!yone  -\!dimenI%       ** adjust y-location to overlap corners
  \advance \!ytwo   \!dimenI%       ** ditto
  \putrule from {\!xone} {\!yone} to {\!xone} {\!ytwo}
  \putrule from {\!xtwo} {\!yone} to {\!xtwo} {\!ytwo}
%
  \let\!M=\!MRect%                  ** restore coord/dimen mode
  \ignorespaces}

% ** \shaderectangleson
% **   Subsequent rectangles will be shaded according to
% **   the current shading pattern.  Affects \putrectangle, \putbar,
% **   \frame, \sethistograms, and \setbars. See Subsection 7.5 of the manual.
\def\shaderectangleson{%
  \def\!shaderectangle{\!!shaderectangle}%
  \ignorespaces}
% ** \shaderectanglesoff
% **    Suppresses  \shaderectangleson.  The default.
\def\shaderectanglesoff{%
  \def\!shaderectangle{}%
  \ignorespaces}

\shaderectanglesoff

% ** The following internal routine shades the current rectangle, when
% **   \!shaderectangle = \!!shaderectangle .
\def\!!shaderectangle{%
  \!dimenA=\!xtwo  \advance \!dimenA -\!xone
  \!dimenB=\!ytwo  \advance \!dimenB -\!yone
  \ifdim \!dimenA<\!dimenB
    \!startvshade (\!xone,\!yone,\!ytwo)
    \!lshade      (\!xtwo,\!yone,\!ytwo)
  \else
    \!starthshade (\!yone,\!xone,\!xtwo)
    \!lshade      (\!ytwo,\!xone,\!xtwo)
  \fi
  \ignorespaces}

% ** \frame [<SEPARATION>] {TEXT}
% ** Draws a frame of thickness linethickness about the box enclosing
% **   TEXT; the frame is separated from the box by a distance of
% **   SEPARATION.  The result is an hbox with the same baseline as TEXT.
% **   If <SEPARATION> is omitted, you get the effect of <0pt>.
% ** See Subsection 4.2 of the manual.
\def\frame{%
  \!ifnextchar<{\!frame}{\!frame<\!zpt> }}
\long\def\!frame<#1> #2{%
  \beginpicture
    \setcoordinatesystem units <1pt,1pt> point at 0 0
    \put {#2} [Bl] at 0 0
    \!dimenA=#1\relax
    \!dimenB=\!wd \advance \!dimenB \!dimenA
    \!dimenC=\!ht \advance \!dimenC \!dimenA
    \!dimenD=\!dp \advance \!dimenD \!dimenA
    \let\!MFr=\!M
    \!setdimenmode
    \putrectangle corners at {-\!dimenA} {-\!dimenD} and {\!dimenB} {\!dimenC}
    \!setcoordmode
    \let\!M=\!MFr
  \endpicture
  \ignorespaces}

% ** \rectangle <WIDTH> <HEIGHT>
% ** Constructs a rectangle of width WIDTH and heigth HEIGHT.
% ** See Subsection 4.2 of the manual.
\def\rectangle <#1> <#2> {%
  \setbox0=\hbox{}\wd0=#1\ht0=#2\frame {\box0}}


% *********************************************
% ***  CURVES  (Upper level \plot commands) ***
% *********************************************
%
% ** User commands
% **   \plot  DATA  /
% **   \plot  "FILE NAME"
% **   \setquadratic
% **   \setlinear
% **   \sethistograms
% **   \vshade  ...
% **   \hshade  ...

% \plot: multi-purpose command. Draws histograms, bar graphs, piecewise-linear
% or piecewise quadratic curves, depending on the setting of \!drawcurve.
% See Subsections 4.3-4.5, 5.1, 5.2 of the manual.
\def\plot{%
  \!ifnextchar"{\!plotfromfile}{\!drawcurve}}
\def\!plotfromfile"#1"{%
  \expandafter\!drawcurve \input #1 /}

% Command to set piecewise quadratic mode
% See Subsections 5.1, 7.3, and 7.4 of the manual.
\def\setquadratic{%
  \let\!drawcurve=\!qcurve
  \let\!!Shade=\!!qShade
  \let\!!!Shade=\!!!qShade}

% Command to set piecewise linear mode
% See Subsections 5.1, 7.3, and 7.4 of the manual.
\def\setlinear{%
  \let\!drawcurve=\!lcurve
  \let\!!Shade=\!!lShade
  \let\!!!Shade=\!!!lShade}

% Command to set histogram mode
% See Subsection 4.3 of the manual.
\def\sethistograms{%
  \let\!drawcurve=\!hcurve}

% Commands to cycle through list of coordinates in piecewise quadratic
% interpolation mode
\def\!qcurve #1 #2 {%
  \!start (#1,#2)
  \!Qjoin}
\def\!Qjoin#1 #2 #3 #4 {%
  \!qjoin (#1,#2) (#3,#4)             % \!qjoin  is defined in QUADRATIC
  \!ifnextchar/{\!finish}{\!Qjoin}}

% Commands to cycle through list of coordinates in piecewise linear
% interpolation mode
\def\!lcurve #1 #2 {%
  \!start (#1,#2)
  \!Ljoin}
\def\!Ljoin#1 #2 {%
  \!ljoin (#1,#2)                    % \!ljoin  is defined in LINEAR
  \!ifnextchar/{\!finish}{\!Ljoin}}

\def\!finish/{\ignorespaces}

% Command to cycle through list of coordinates in histogram mode
\def\!hcurve #1 #2 {%
  \edef\!hxS{#1}%
  \edef\!hyS{#2}%
  \!hjoin}
\def\!hjoin#1 #2 {%
  \putrectangle corners at {\!hxS} {\!hyS} and {#1} {#2}
  \edef\!hxS{#1}%
  \!ifnextchar/{\!finish}{\!hjoin}}


% \vshade: See Subsection 7.3 of the manual.
\def\vshade #1 #2 #3 {%
  \!startvshade (#1,#2,#3)
  \!Shadewhat}

% \hshade: See Subsection 7.4 of the manual.
\def\hshade #1 #2 #3 {%
  \!starthshade (#1,#2,#3)
  \!Shadewhat}

% Commands to cycle through coordinates and optional "edge effect"
% fields while shading.
\def\!Shadewhat{%
  \futurelet\!nextchar\!Shade}
\def\!Shade{%
  \if <\!nextchar
    \def\!nextShade{\!!Shade}%
  \else
    \if /\!nextchar
      \def\!nextShade{\!finish}%
    \else
      \def\!nextShade{\!!!Shade}%
    \fi
  \fi
  \!nextShade}
\def\!!lShade<#1> #2 #3 #4 {%
  \!lshade <#1> (#2,#3,#4)                 % \!lshade is defined in SHADING
  \!Shadewhat}
\def\!!!lShade#1 #2 #3 {%
  \!lshade (#1,#2,#3)
  \!Shadewhat}
\def\!!qShade<#1> #2 #3 #4 #5 #6 #7 {%
  \!qshade <#1> (#2,#3,#4) (#5,#6,#7)      % \!qshade is defined in SHADING
  \!Shadewhat}
\def\!!!qShade#1 #2 #3 #4 #5 #6 {%
  \!qshade (#1,#2,#3) (#4,#5,#6)
  \!Shadewhat}

% ** Set default interpolation mode
\setlinear


%  ********************************************
%  *** DASHPATTERNS (Sets up dash patterns) ***
%  ********************************************

%  **  User commands:
%  **    \setdashpattern <DIMEN1,DIMEN2,DIMEN3,...>
%  **    \setdots <INTRADOT_DISTANCE>
%  **    \setdotsnear <INTRADOT_DISTANCE> for <ARC LENGTH>
%  **    \setdashes <DASH/SKIP_DISTANCE>
%  **    \setdashesnear <DASH/SKIP_DISTANCE> for <ARC LENGTH>
%  **    \setsolid
%  **    \findlength {CURVE CMDS}

%  **  Internal commands:
%  **    \!dashingon
%  **    \!dashingoff

%  **  Dash patterns are specified by a balanced token list whose complete
%  **    expansion has the form: DIMEN1,DIMEN2,DIMEN3,DIMEN4,... ; this
%%produces
%  **    an arc of length DIMEN1, a skip of length DIMEN2, an arc of length
%  **    DIMEN3, a skip of length DIMEN4, ... .  Any number of DIMEN values may
%  **    be given. The pattern is repeated as many times (perhaps fractional)
%  **    as necessary to draw the curve.
%  **  A dash pattern remains in effect until it is overridden by a call to
%  **    \setdashpattern, or to \setdots, \setdotsnear ... , \setdashes,
%  **    \setdashesnear ... , or \setsolid.
%  **  Solid lines are the default.


%  **  \def\setdashpattern <DIMEN1,DIMEN2,DIMEN3,...>
%  **  The following routine converts a balanced list of tokens whose
%  **  complete expansion has the form  DIMEN1,DIMEN2, ... , DIMENk  into
%  **  three list macros that are used in drawing dashed rules and curves:
%  **    !Flist:   \!Rule{DIMEN1}\!Skip{DIMEN2}\!Rule{DIMEN3}\!Skip{DIMEN4} ...
%  **    !Blist:   ...\!Skip{DIMEN4}\!Rule{DIMEN3}\!Skip{DIMEN2}\!Rule{DIMEN1}
%  **    !UDlist:  \\{DIMEN1}\\{DIMEN2}\\{DIMEN3}\\{DIMEN4} ...;
%  **  calculates \!leaderlength := DIMEN1 + ... + DIMENk; and
%  **  sets the curve drawing routines to dash mode.
%  **  Those lists are used by the curve drawing routines.
%  **  Dimenj ... may be given as an explicit dimension (e.g., 5pt), or
%  **  as an expression involving a dimension register (e.g., -2.5\dimen0).
%  **  See Subsection 6.2 of the manual
\def\setdashpattern <#1>{%
  \def\!Flist{}\def\!Blist{}\def\!UDlist{}%
  \!countA=0
  \!ecfor\!item:=#1\do{%
    \!dimenA=\!item\relax
    \expandafter\!rightappend\the\!dimenA\withCS{\\}\to\!UDlist%
    \advance\!countA  1
    \ifodd\!countA
      \expandafter\!rightappend\the\!dimenA\withCS{\!Rule}\to\!Flist%
      \expandafter\!leftappend\the\!dimenA\withCS{\!Rule}\to\!Blist%
    \else
      \expandafter\!rightappend\the\!dimenA\withCS{\!Skip}\to\!Flist%
      \expandafter\!leftappend\the\!dimenA\withCS{\!Skip}\to\!Blist%
    \fi}%
  \!leaderlength=\!zpt
  \def\!Rule##1{\advance\!leaderlength  ##1}%
  \def\!Skip##1{\advance\!leaderlength  ##1}%
  \!Flist%
  \ifdim\!leaderlength>\!zpt
  \else
    \def\!Flist{\!Skip{24in}}\def\!Blist{\!Skip{24in}}\ignorespaces
    \def\!UDlist{\\{\!zpt}\\{24in}}\ignorespaces
    \!leaderlength=24in
  \fi
  \!dashingon}


%  **  \!dashingon  -- puts the curve drawing routines into dash mode
%  **  \!dashingoff -- puts the curve drawing routines into solid mode
%  **  These are internal commands, invoked by \setdashpattern and \setsolid
\def\!dashingon{%
  \def\!advancedashing{\!!advancedashing}%
  \def\!drawlinearsegment{\!lineardashed}%
  \def\!puthline{\!putdashedhline}%
  \def\!putvline{\!putdashedvline}%
%  \def\!putsline{\!putdashedsline}%
  \ignorespaces}%
\def\!dashingoff{%
  \def\!advancedashing{\relax}%
  \def\!drawlinearsegment{\!linearsolid}%
  \def\!puthline{\!putsolidhline}%
  \def\!putvline{\!putsolidvline}%
%  \def\!putsline{\!putsolidsline}%
  \ignorespaces}


%  **  \setdots <LENGTH>  --  sets up a dot/skip pattern where dot (actually
%  **    the current plotsymbol) is plunked down once for every LENGTH
%  **    traveled along the curve.  LENGTH defaults to 5pt.
%  **    See Subsection 6.1 of the manual.
\def\setdots{%
  \!ifnextchar<{\!setdots}{\!setdots<5pt>}}
\def\!setdots<#1>{%
  \!dimenB=#1\advance\!dimenB -\plotsymbolspacing
  \ifdim\!dimenB<\!zpt
    \!dimenB=\!zpt
  \fi
\setdashpattern <\plotsymbolspacing,\!dimenB>}

% ** \setdotsnear <LENGTH> for <ARC LENGTH>
% ** sets up a dot pattern where the dots are approximately LENGTH apart,
% ** the total length of the pattern is ARC LENGTH, and the pattern
% ** begins and ends with a dot. See Subsection 6.3 of the manual.
\def\setdotsnear <#1> for <#2>{%
  \!dimenB=#2\relax  \advance\!dimenB -.05pt
  \!dimenC=#1\relax  \!countA=\!dimenC
  \!dimenD=\!dimenB  \advance\!dimenD .5\!dimenC  \!countB=\!dimenD
  \divide \!countB  \!countA
  \ifnum 1>\!countB
    \!countB=1
  \fi
  \divide\!dimenB  \!countB
  \setdots <\!dimenB>}

%  **  \setdashes <LENGTH>  --  sets up a dash/skip pattern where the dash
%  **    and the skip are each of length LENGTH (the dash is formed by
%  **    plunking down the current plotsymbol over an arc of length LENGTH
%  **    and so may actually be longer than LENGTH.  LENGTH defaults to 5pt.
%  **    See Subsection 6.1 of the manual.
\def\setdashes{%
  \!ifnextchar<{\!setdashes}{\!setdashes<5pt>}}
\def\!setdashes<#1>{\setdashpattern <#1,#1>}

% ** \setdashesnear ...
% ** Like \setdotsnear; the pattern begins and ends with a dash.
% ** See Subsection 6.3 of the manual.
\def\setdashesnear <#1> for <#2>{%
  \!dimenB=#2\relax
  \!dimenC=#1\relax  \!countA=\!dimenC
  \!dimenD=\!dimenB  \advance\!dimenD .5\!dimenC  \!countB=\!dimenD
  \divide \!countB  \!countA
  \ifodd \!countB
  \else
    \advance \!countB  1
  \fi
  \divide\!dimenB  \!countB
  \setdashes <\!dimenB>}

%  **  \setsolid  --  puts the curve drawing routines in "solid line" mode,
%  **    the default mode.  See Subsection 6.1 of the manual.
\def\setsolid{%
  \def\!Flist{\!Rule{24in}}\def\!Blist{\!Rule{24in}}%
  \def\!UDlist{\\{24in}\\{\!zpt}}%
  \!dashingoff}
\setsolid

%  **  \findlength {CURVE CMDS}
%  **  PiCTeX executes the \start, \ljoin, and \qjoin cmds comprising
%  **  CURVE CMDS without plotting anything, but stashes the length
%  **  of the phantom curve away in \totalarclength.
%  **  See Subsection 6.3 of the manual.
\def\findlength#1{%
  \begingroup
    \setdashpattern <0pt, \maxdimen>
    \setplotsymbol ({})
    \dontsavelinesandcurves
    #1%
  \endgroup
  \ignorespaces}


% *************************************************************
% *** DIVISION  (Does long division of dimension registers) ***
% *************************************************************

% ** User command:
% **   \Divide {DIVIDEND} by {DIVISOR} forming {RESULT}

% ** Internal command
% **   \!divide{DIVIDEND}{DIVISOR}{RESULT}

% **  \!divide DIVIDEND [by] DIVISOR [to get] ANSWER
% **  Divides the dimension DIVIDEND by the dimension DIVISOR, placing the
% **  quotient in the dimension register ANSWER.  Values are understood to
% **  be in points.  E.g.  12.5pt/1.4pt=8.92857pt.
% **  Quotient is accurate to 1/65536pt=2**[-16]pt
% **  |DIVISOR| should be < 2048pt (about 28 inches).
\def\!divide#1#2#3{%
  \!dimenB=#1%                      **  dimB  holds current remainder (r)
  \!dimenC=#2%                      **  dimC  holds divisor (d)
  \!dimenD=\!dimenB%                **  dimD  holds quotient q=r/d for this
  \divide \!dimenD \!dimenC%        **    step, in units of scaled pts
  \!dimenA=\!dimenD%                **  dimA  eventually holds answer (a)
  \multiply\!dimenD \!dimenC%       **  r <-- r - dq
  \advance\!dimenB -\!dimenD%       **  First step complete. Have integer part
%                                   **  of a, and corresponding remainder.
  \!dimenD=\!dimenC%                **  Temporarily use dimD to hold |d|
    \ifdim\!dimenD<\!zpt \!dimenD=-\!dimenD
  \fi
  \ifdim\!dimenD<64pt%              **  Branch on the magnitude of |d|
    \!divstep[\!tfs]\!divstep[\!tfs]%
  \else
    \!!divide
  \fi
  #3=\!dimenA\ignorespaces}

% **  The following code handles divisors  d  with
% **    (1)  .88in =  64pt <= d <  256pt =  3.54in
% **    (2) 3.54in = 256pt <= d < 2048pt = 28.34in
% **  Anything bigger than that may result in an overflow condition.
% **  For our purposes, we should never even see case (2).
\def\!!divide{%
  \ifdim\!dimenD<256pt
    \!divstep[64]\!divstep[32]\!divstep[32]%
  \else
    \!divstep[8]\!divstep[8]\!divstep[8]\!divstep[8]\!divstep[8]%
    \!dimenA=2\!dimenA
  \fi}


% **  The following macro does the real long division work.
\def\!divstep[#1]{%                 **  #1 = "B"
  \!dimenB=#1\!dimenB%              **  r <-- B*r
  \!dimenD=\!dimenB%                **  dimD  holds quotient q=r/d for this
    \divide \!dimenD by \!dimenC%   **    step, in units of scaled pts
  \!dimenA=#1\!dimenA%              **  a <-- B*a + q
    \advance\!dimenA by \!dimenD%
  \multiply\!dimenD by \!dimenC%    **  r <-- r - dq
    \advance\!dimenB by -\!dimenD}

% **  \Divide:  See Subsection 9.3 of the manual.
\def\Divide <#1> by <#2> forming <#3> {%
  \!divide{#1}{#2}{#3}}


% *********************************************
% *** ELLIPSES (Draws ellipses and circles) ***
% *********************************************

% ** User commands
% **   \ellipticalarc  axes ratio A:B  DEGREES degrees from XSTART YSTART
% **      center at XCENTER YCENTER
% **   \circulararc DEGREES degrees from XSTART YSTART
% **      center at XCENTER YCENTER

% ** Internal command
% **   \!sinandcos{32*ANGLE in radians}{32*SIN}{32*COS}


% **   \ellipticalarc  axes ratio A:B  DEGREES degrees from XSTART YSTART
% **      center at XCENTER YCENTER
% **    Draws a elliptical arc starting at the coordinate point
%%(XSTART,YSTART).
% **    The center of the ellipse of which the arc is a segment is at
% **      (XCENTER,YCENTER).
% **    The arc extends through an angle of DEGREES degrees (may be + or -).
% **    A:B is the ratio of the length of the xaxis to the length of
% **      the yaxis of the ellipse
% **    Sqrt{[(XSTART-XCENTER)/A]**2 + [(YSTART-YCENTER)/B]**2}
% **      must be < 512pt (about 7in).
% **    Doesn't modify the dimensions (ht, dp, wd) of the PiCture under
% **      construction.

% ** \circulararc  --  See Subsection 5.3 of the manual.
\def\circulararc{%
  \ellipticalarc axes ratio 1:1 }

% ** \ellipticalarc  --  See Subsection 5.3 of the manual.
\def\ellipticalarc axes ratio #1:#2 #3 degrees from #4 #5 center at #6 #7 {%
  \!angle=#3pt\relax%                    ** get angle
  \ifdim\!angle>\!zpt
    \def\!sign{}%                        ** counterclockwise
  \else
    \def\!sign{-}\!angle=-\!angle%       ** clockwise
  \fi
  \!xxloc=\!M{#6}\!xunit%                ** convert CENTER to dimension
  \!yyloc=\!M{#7}\!yunit
  \!xxS=\!M{#4}\!xunit%                  ** get STARTing point on rim of
%%ellipse
  \!yyS=\!M{#5}\!yunit
  \advance\!xxS -\!xxloc%                ** make center of ellipse (0,0)
  \advance\!yyS -\!yyloc
  \!divide\!xxS{#1pt}\!xxS %             ** scale point on ellipse to point on
  \!divide\!yyS{#2pt}\!yyS %                 corresponding circle
%
  \let\!MC=\!M%                          ** save current c/d mode
  \!setdimenmode%                        ** go into dimension mode
%
  \!xS=#1\!xxS  \advance\!xS\!xxloc
  \!yS=#2\!yyS  \advance\!yS\!yyloc
  \!start (\!xS,\!yS)%
  \!loop\ifdim\!angle>14.9999pt%         ** draw in major portion of ellipse
    \!rotate(\!xxS,\!yyS)by(\!cos,\!sign\!sin)to(\!xxM,\!yyM)
    \!rotate(\!xxM,\!yyM)by(\!cos,\!sign\!sin)to(\!xxE,\!yyE)
    \!xM=#1\!xxM  \advance\!xM\!xxloc  \!yM=#2\!yyM  \advance\!yM\!yyloc
    \!xE=#1\!xxE  \advance\!xE\!xxloc  \!yE=#2\!yyE  \advance\!yE\!yyloc
    \!qjoin (\!xM,\!yM) (\!xE,\!yE)
    \!xxS=\!xxE  \!yyS=\!yyE
    \advance \!angle -15pt
  \repeat
  \ifdim\!angle>\!zpt%                   ** complete remaining arc, if any
    \!angle=100.53096\!angle%            ** convert angle to radians, divide
    \divide \!angle 360 %                **   by 2, and multiply by 32
    \!sinandcos\!angle\!!sin\!!cos%      ** get 32*sin & 32*cos
    \!rotate(\!xxS,\!yyS)by(\!!cos,\!sign\!!sin)to(\!xxM,\!yyM)
    \!rotate(\!xxM,\!yyM)by(\!!cos,\!sign\!!sin)to(\!xxE,\!yyE)
    \!xM=#1\!xxM  \advance\!xM\!xxloc  \!yM=#2\!yyM  \advance\!yM\!yyloc
    \!xE=#1\!xxE  \advance\!xE\!xxloc  \!yE=#2\!yyE  \advance\!yE\!yyloc
    \!qjoin (\!xM,\!yM) (\!xE,\!yE)
  \fi
%
  \let\!M=\!MC%                          ** restore c/d mode
  \ignorespaces}%                        **   if appropriate


%  ** \!rotate(XREG,YREG)by(32cos,32sin)to(XXREG,YYREG)
%  ** rotates (XREG,YREG) by angle with specfied scaled cos & sin to
%  ** (XXREG,YYREG).  Uses \!dimenA & \!dimenB as scratch registers.
\def\!rotate(#1,#2)by(#3,#4)to(#5,#6){%
  \!dimenA=#3#1\advance \!dimenA -#4#2%   ** Rcos(x+t)=Rcosx*cost - Rsinx*sint
  \!dimenB=#3#2\advance \!dimenB  #4#1%   ** Rsin(x+t)=Rsinx*cost + Rcosx*sint
  \divide \!dimenA 32  \divide \!dimenB 32
  #5=\!dimenA  #6=\!dimenB
  \ignorespaces}
\def\!sin{4.17684}%                       ** 32*sin(pi/24) (pi/24=7.5deg)
\def\!cos{31.72624}%                      ** 32*cos(pi/24)


%  ** \!sinandcos{32*ANGLE in radians}{\SINCS}{\COSCS}
%  **   Computes the 32*sine and 32*cosine of a small ANGLE expressed in
%  **   radians/32 and puts these values in the replacement texts of
%  **   \SINCS and \COSCS
\def\!sinandcos#1#2#3{%
 \!dimenD=#1%                **  angle is expressed in radians/32: 1pt =
%%1/32rad
 \!dimenA=\!dimenD%          **  dimA will eventually contain 32sin(angle)in
%%pts
 \!dimenB=32pt%              **  dimB will eventually contain 32cos(angle)in
%%pts
 \!removept\!dimenD\!value%  **  get value of 32*angle, without "pt"
 \!dimenC=\!dimenD%          **  holds 32*angle**i/i! in pts
 \!dimenC=\!value\!dimenC \divide\!dimenC by 64 %   ** now 32*angle**2/2
 \advance\!dimenB by -\!dimenC%                     ** 32-32*angle**2/2
 \!dimenC=\!value\!dimenC \divide\!dimenC by 96 %   ** now 32*angle**3/3!
 \advance\!dimenA by -\!dimenC%                     ** now
%%32*(angle-angle**3/6)
 \!dimenC=\!value\!dimenC \divide\!dimenC by 128 %  ** now 32*angle**4/4!
 \advance\!dimenB by \!dimenC%
 \!removept\!dimenA#2%                              ** set 32*sin(angle)
 \!removept\!dimenB#3%                              ** set 32*cos(angle)
 \ignorespaces}


% *****************************************************************
% ***  RULES  (Draws rules, i.e., horizontal & vertical lines)  ***
% *****************************************************************

% **  User command:
% **    \putrule [<XDIMEN,YDIMEN>] from  XCOORD1 YCOORD1
% **      to  XCOORD2 YCOORD2

% **  Internal commands:
% **    \!puthline [<XDIMEN,YDIMEN>]    (h = horizontal)
% **      Set by dashpat to either: \!putsolidhline  or \!putdashedhline
% **    \!putvline [<XDIMEN,YDIMEN>]    (v = vertical)
% **      Either:  \!putsolidvline  or  \!putdashedvline


% **  \putrule [<XDIMEN,YDIMEN>] from XCOORD1 YCOORD1
% **    to XCOORD2 YCOORD2
% **  Draws a rule -- dashed or solid depending on the current dash pattern --
% **    from (X1,Y1) to (X2,Y2).  Uses TEK's  \hrule & \vrule & \leaders
% **    constructions to handle horizontal & vertical lines efficiently both
% **    in terms of execution time and space in the DVI file.
% **  See Subsection 4.1 of the manual.
\def\putrule#1from #2 #3 to #4 #5 {%
  \!xloc=\!M{#2}\!xunit  \!xxloc=\!M{#4}\!xunit%
  \!yloc=\!M{#3}\!yunit  \!yyloc=\!M{#5}\!yunit%
  \!dxpos=\!xxloc  \advance\!dxpos by -\!xloc
  \!dypos=\!yyloc  \advance\!dypos by -\!yloc
%
  \ifdim\!dypos=\!zpt
    \def\!!Line{\!puthline{#1}}\ignorespaces
  \else
    \ifdim\!dxpos=\!zpt
      \def\!!Line{\!putvline{#1}}\ignorespaces
    \else
       \def\!!Line{}
    \fi
  \fi
  \let\!ML=\!M%           ** save current coord\dimen mode
  \!setdimenmode%         ** express locations in dimens
  \!!Line%
  \let\!M=\!ML%           ** restore previous c/d mode
  \ignorespaces}


% **  \!putsolidhline [<XDIMEN,YDIMEN>]
% **  Place horizontal solid line
\def\!putsolidhline#1{%
  \ifdim\!dxpos>\!zpt
    \put{\!hline\!dxpos}#1[l] at {\!xloc} {\!yloc}
  \else
    \put{\!hline{-\!dxpos}}#1[l] at {\!xxloc} {\!yyloc}
  \fi
  \ignorespaces}

% **  \!putsolidvline [shifted <XDIMEN,YDIMEN>]
% **  Place vertical solid line
\def\!putsolidvline#1{%
  \ifdim\!dypos>\!zpt
    \put{\!vline\!dypos}#1[b] at {\!xloc} {\!yloc}
  \else
    \put{\!vline{-\!dypos}}#1[b] at {\!xxloc} {\!yyloc}
  \fi
  \ignorespaces}

\def\!hline#1{\hbox to #1{\leaders \hrule height\linethickness\hfill}}
\def\!vline#1{\vbox to #1{\leaders \vrule width\linethickness\vfill}}


% **  \!putdashedhline [<XDIMEN,YDIMEN>]
% **  Place dashed horizontal line
\def\!putdashedhline#1{%
  \ifdim\!dxpos>\!zpt
    \!DLsetup\!Flist\!dxpos
    \put{\hbox to \!totalleaderlength{\!hleaders}\!hpartialpattern\!Rtrunc}
      #1[l] at {\!xloc} {\!yloc}
  \else
    \!DLsetup\!Blist{-\!dxpos}
    \put{\!hpartialpattern\!Ltrunc\hbox to \!totalleaderlength{\!hleaders}}
      #1[r] at {\!xloc} {\!yloc}
  \fi
  \ignorespaces}

% **  \!putdashedhline [<XDIMEN,YDIMEN>]
% **  Place dashed vertical line
\def\!putdashedvline#1{%
  \!dypos=-\!dypos%            ** vertical leaders go from top to bottom
  \ifdim\!dypos>\!zpt
    \!DLsetup\!Flist\!dypos
    \put{\vbox{\vbox to \!totalleaderlength{\!vleaders}
      \!vpartialpattern\!Rtrunc}}#1[t] at {\!xloc} {\!yloc}
  \else
    \!DLsetup\!Blist{-\!dypos}
    \put{\vbox{\!vpartialpattern\!Ltrunc
      \vbox to \!totalleaderlength{\!vleaders}}}#1[b] at {\!xloc} {\!yloc}
  \fi
  \ignorespaces}


% **  The rest of the macros in this section are subroutines used by
% **  \!putdashedhline and \!putdashedvline.
\def\!DLsetup#1#2{%            ** Dashed-Line set up
  \let\!RSlist=#1%             ** set !Rule-Skip list
  \!countB=#2%                 ** convert rule length to integer (number of
%%sps)
  \!countA=\!leaderlength%     ** ditto, leaderlength
  \divide\!countB by \!countA% ** number of complete leader units
  \!totalleaderlength=\!countB\!leaderlength
  \!Rresiduallength=#2%
  \advance \!Rresiduallength by -\!totalleaderlength%  \** excess length
  \!Lresiduallength=\!leaderlength
  \advance \!Lresiduallength by -\!Rresiduallength
  \ignorespaces}

\def\!hleaders{%
  \def\!Rule##1{\vrule height\linethickness width##1}%
  \def\!Skip##1{\hskip##1}%
  \leaders\hbox{\!RSlist}\hfill}

\def\!hpartialpattern#1{%
  \!dimenA=\!zpt \!dimenB=\!zpt
  \def\!Rule##1{#1{##1}\vrule height\linethickness width\!dimenD}%
  \def\!Skip##1{#1{##1}\hskip\!dimenD}%
  \!RSlist}

\def\!vleaders{%
  \def\!Rule##1{\hrule width\linethickness height##1}%
  \def\!Skip##1{\vskip##1}%
  \leaders\vbox{\!RSlist}\vfill}

\def\!vpartialpattern#1{%
  \!dimenA=\!zpt \!dimenB=\!zpt
  \def\!Rule##1{#1{##1}\hrule width\linethickness height\!dimenD}%
  \def\!Skip##1{#1{##1}\vskip\!dimenD}%
  \!RSlist}

\def\!Rtrunc#1{\!trunc{#1}>\!Rresiduallength}
\def\!Ltrunc#1{\!trunc{#1}<\!Lresiduallength}

\def\!trunc#1#2#3{%
  \!dimenA=\!dimenB
  \advance\!dimenB by #1%
  \!dimenD=\!dimenB  \ifdim\!dimenD#2#3\!dimenD=#3\fi
  \!dimenC=\!dimenA  \ifdim\!dimenC#2#3\!dimenC=#3\fi
  \advance \!dimenD by -\!dimenC}


%  ****************************************************************
%  ***  LINEAR ARC  (Draws straight lines -- solid and dashed)  ***
%  ****************************************************************

%  **  User commands
%  **    \inboundscheckoff
%  **    \inboundscheckon

%  **  Internal commands
%  **    \!start (XCOORD,YCOORD)
%  **    \!ljoin (XCOORD,YCOORD)
%  **    \!drawlinearsegment  --  set by \dashpat to either
%  **      \!linearsolid  or  \!lineardashed
%  **    \!advancedashing     --  set by \dashpat to either
%  **       \relax  or  \!!advancedashing
%  **    \!plotifinbounds     --  set by \inboundscheck off/on to either
%  **       \!plot  or  \!!plotifinbounds
%  **    \!initinboundscheck  --  set by \inboundscheck off/on to either
%  **       \relax  or  \!!initinboundscheck


%  \plotsymbolspacing  ** distance between consecutive plot positions
%  \!xS                ** starting x
%  \!yS                ** starting y
%  \!xE                ** ending   x
%  \!yE                ** ending   y
%  \!xdiff             ** x_end - x_start
%  \!ydiff             ** y_end - y_start
%  \!distacross        ** how far along curve next point to be plotted is
%  \!arclength         ** approximate length of arc for current interval
%  \!downlength        ** remaining length for "pen" to be down
%  \!uplength          ** length for "pen" to be down
%  \!intervalno        ** counts segments to curve
%  \totalarclength     ** cumulative distance along curve
%  \!npoints           ** approximately  (arc length / plotsymbolspacing)

%  **  Calls -- \!Pythag, \!divide, \!plot


%  **  \!start (XCOORD,YCOORD)
%  **  Sets initial point for linearly (or quadratically) interpolated curve
\def\!start (#1,#2){%
  \!plotxorigin=\!xorigin  \advance \!plotxorigin by \!plotsymbolxshift
  \!plotyorigin=\!yorigin  \advance \!plotyorigin by \!plotsymbolyshift
  \!xS=\!M{#1}\!xunit \!yS=\!M{#2}\!yunit
  \!rotateaboutpivot\!xS\!yS
  \!copylist\!UDlist\to\!!UDlist% **\!UDlist has the form
%%\\{dimen1}\\{dimen2}..
%                                 ** Routine will draw dashed line with pen
%                                 ** down for dimen1, up for dimen2, ...
  \!getnextvalueof\!downlength\from\!!UDlist
  \!distacross=\!zpt%             ** 1st point goes at start of curve
  \!intervalno=0 %                ** initialize interval counter
  \global\totalarclength=\!zpt%   ** initialize distance traveled along curve
  \ignorespaces}


%  **  \!ljoin (XCOORD,YCOORD)
%  **  Draws a straight line starting at the last point specified
%  **    by the most recent \!start, \!ljoin, or \!qjoin, and
%  **    ending at (XCOORD,YCOORD).
\def\!ljoin (#1,#2){%
  \advance\!intervalno by 1
  \!xE=\!M{#1}\!xunit \!yE=\!M{#2}\!yunit
  \!rotateaboutpivot\!xE\!yE
  \!xdiff=\!xE \advance \!xdiff by -\!xS%**  xdiff = xE - xS
  \!ydiff=\!yE \advance \!ydiff by -\!yS%**  ydiff = yE - yS
  \!Pythag\!xdiff\!ydiff\!arclength%     **  arclength =
%%sqrt(xdiff**2+ydiff**2)
  \global\advance \totalarclength by \!arclength%
  \!drawlinearsegment%   ** set by dashpat to \!linearsolid or \!lineardashed
  \!xS=\!xE \!yS=\!yE%   ** shift ending points to starting points
  \ignorespaces}


% **  The following routine is used to draw a "solid" line between (xS,yS)
% **  and (xE,yE).  Points are spaced nearly every  \plotsymbolspacing length
% **  along the line.
\def\!linearsolid{%
  \!npoints=\!arclength
  \!countA=\plotsymbolspacing
  \divide\!npoints by \!countA%      ** now #pts =. arclength/plotsymbolspacing
  \ifnum \!npoints<1
    \!npoints=1
  \fi
  \divide\!xdiff by \!npoints
  \divide\!ydiff by \!npoints
  \!xpos=\!xS \!ypos=\!yS
%
  \loop\ifnum\!npoints>-1
    \!plotifinbounds
    \advance \!xpos by \!xdiff
    \advance \!ypos by \!ydiff
    \advance \!npoints by -1
  \repeat
  \ignorespaces}


% ** The following routine is used to draw a dashed line between (xS,yS)
% ** and (xE,yE). The dash pattern continues from the previous segment.
\def\!lineardashed{%
% **
  \ifdim\!distacross>\!arclength
    \advance \!distacross by -\!arclength  %nothing to plot in this interval
%
  \else
%
    \loop\ifdim\!distacross<\!arclength
%     ** plot point, interpolating linearly in x and y
      \!divide\!distacross\!arclength\!dimenA%  ** dimA = across/arclength
      \!removept\!dimenA\!t%  ** \!t holds value in dimA, without the "pt"
      \!xpos=\!t\!xdiff \advance \!xpos by \!xS
      \!ypos=\!t\!ydiff \advance \!ypos by \!yS
      \!plotifinbounds
      \advance\!distacross by \plotsymbolspacing
      \!advancedashing
    \repeat
%
    \advance \!distacross by -\!arclength%    ** prepare for next interval
  \fi
  \ignorespaces}


\def\!!advancedashing{%
  \advance\!downlength by -\plotsymbolspacing
  \ifdim \!downlength>\!zpt
  \else
    \advance\!distacross by \!downlength
    \!getnextvalueof\!uplength\from\!!UDlist
    \advance\!distacross by \!uplength
    \!getnextvalueof\!downlength\from\!!UDlist
  \fi}


% ** \inboundscheckoff & \inboundscheckon: See Subsection 5.5 of the manual.
\def\inboundscheckoff{%
  \def\!plotifinbounds{\!plot(\!xpos,\!ypos)}%
  \def\!initinboundscheck{\relax}\ignorespaces}
\def\inboundscheckon{%
  \def\!plotifinbounds{\!!plotifinbounds}%
  \def\!initinboundscheck{\!!initinboundscheck}%
  \!initinboundscheck\ignorespaces}
\inboundscheckoff

% ** The following code plots the current point only if it falls in the
% ** current plotarea.  It doesn't matter if the coordinate system has
% ** changed since the plotarea was set up.  However, shifts of the plot
% ** are ignored (how the plotsymbol stands relative to its plot position is
% ** unknown anyway).
\def\!!plotifinbounds{%
  \ifdim \!xpos<\!checkleft
  \else
    \ifdim \!xpos>\!checkright
    \else
      \ifdim \!ypos<\!checkbot
      \else
         \ifdim \!ypos>\!checktop
         \else
           \!plot(\!xpos,\!ypos)
         \fi
      \fi
    \fi
  \fi}


\def\!!initinboundscheck{%
  \!checkleft=\!arealloc     \advance\!checkleft by \!xorigin
  \!checkright=\!arearloc    \advance\!checkright by \!xorigin
  \!checkbot=\!areabloc      \advance\!checkbot by \!yorigin
  \!checktop=\!areatloc      \advance\!checktop by \!yorigin}


% *********************************
% *** LOGTEN  (Log_10 function) ***
% *********************************
%
% ** \!logten{X}
% ** Calculates log_10 of X.  X and LOG10(X) are in fixed point notation.
% **  X must be positive; it may have an optional `+' sign; any number
% **  of digits may be specified for X.  The absolute error in LOG10(X) is
% **  less than .0001 (probably < .00006).  That's about as good as you
% **  hope for, since TEX only operates to 5 figures after the decimal
% **  point anyway.

%  \!rootten=3.162278pt       **** These are values are set in ALLOCATIONS
%  \!tenAe=2.543275pt  (=A5)
%  \!tenAc=2.773839pt  (=A3)
%  \!tenAa=8.690286pt  (=A1)

\def\!logten#1#2{%
  \expandafter\!!logten#1\!nil
  \!removept\!dimenF#2%
  \ignorespaces}

\def\!!logten#1#2\!nil{%
  \if -#1%
    \!dimenF=\!zpt
    \def\!next{\ignorespaces}%
  \else
    \if +#1%
      \def\!next{\!!logten#2\!nil}%
    \else
      \if .#1%
        \def\!next{\!!logten0.#2\!nil}%
      \else
        \def\!next{\!!!logten#1#2..\!nil}%
      \fi
    \fi
  \fi
  \!next}

\def\!!!logten#1#2.#3.#4\!nil{%
  \!dimenF=1pt %                 ** DimF holds log10 original argument
  \if 0#1%
    \!!logshift#3pt %            ** Argument < 1
  \else %                        ** Argument >= 1
    \!logshift#2/%               ** Shift decimal pt as many places
    \!dimenE=#1.#2#3pt %         **   as there are figures in #2
  \fi %                          ** Now dimE holds revised X want log10 of
  \ifdim \!dimenE<\!rootten%          ** Transform X to XX between sqrt(10)
    \multiply \!dimenE 10 %           **   and 10*sqrt(10)
    \advance  \!dimenF -1pt
  \fi
  \!dimenG=\!dimenE%                  ** dimG <- (XX + 10)
    \advance\!dimenG 10pt
  \advance\!dimenE -10pt %            ** dimE <- (XX - 10)
  \multiply\!dimenE 10 %              ** dimE = 10*(XX-10)
  \!divide\!dimenE\!dimenG\!dimenE%   ** Now dimE=10t==10*(XX-10)/(XX+10)
  \!removept\!dimenE\!t%              ** !t=10t, with "pt" removed
  \!dimenG=\!t\!dimenE%               ** dimG=100t**2
  \!removept\!dimenG\!tt%             ** !tt=100t**2, with "pt" removed
  \!dimenH=\!tt\!tenAe%               ** dimH=10*a5*(10t)**2 /100
    \divide\!dimenH 100
  \advance\!dimenH \!tenAc%           ** ditto + 10*a3
  \!dimenH=\!tt\!dimenH%              ** ditto * (10t)**2 /100
    \divide\!dimenH 100
  \advance\!dimenH \!tenAa%           ** ditto + 10*a1
  \!dimenH=\!t\!dimenH%               ** ditto * 10t / 100
    \divide\!dimenH 100 %             ** Now dimH = log10(XX) - 1
  \advance\!dimenF \!dimenH}%         ** dimF = log10(X)

\def\!logshift#1{%
  \if #1/%
    \def\!next{\ignorespaces}%
  \else
    \advance\!dimenF 1pt
    \def\!next{\!logshift}%
  \fi
  \!next}

 \def\!!logshift#1{%
   \advance\!dimenF -1pt
   \if 0#1%
     \def\!next{\!!logshift}%
   \else
     \if p#1%
       \!dimenF=1pt
       \def\!next{\!dimenE=1p}%
     \else
       \def\!next{\!dimenE=#1.}%
     \fi
   \fi
   \!next}


% ***********************************************************
% *** PICTURES (Basic setups for PiCtures; \put commands) ***
% ***********************************************************

% **  User Commands:
% **    \beginpicture
% **    \endpicture
% **    \endpicturesave <XREG,YREG>
% **    \setcoordinatesystem units <XUNIT,YUNIT> point at XREF YREF
% **    \put {OBJECT} [ORIENTATION] <XSHIFT,YSHIFT> at XCOORD YCOORD
% **    \multiput {OJBECT} [ORIENTATION] <XSHIFT,YSHIFT>) at
% **      XCOORD YCOORD
% **      *NUMBER_OF_TIMES DXCOORD DYCOORD  /
% **    \accountingon
% **    \accountingoff
% **    \stack [ORIENTATION] <LEADING> {LIST OF ITEMS}
% **    \lines [ORIENTATION] {LINES}
% **    \Lines [ORIENTATION] {LINES}
% **    \setdimensionmode
% **    \setcoordinatemode
% **    \Xdistance
% **    \Ydistance

% **  Internal commands:
% **    \!setputobject{OBJECT}{[ORIENTATION]<XSHIFT,YSHIFT>}
% **    \!dimenput{OBJECT}[ORIENTATION]<XSHIFT,YSHIFT>(XDIMEN,YDIMEN)
% **    \!setdimenmode
% **    \!setcoordmode
% **    \!ifdimenmode
% **    \!ifcoordmode


% **  \beginpicture
% **  \endpicture
% **  \endpicturesave <XREG,YREG>
% **    \beginpicture ... \endpicture  creates an hbox.  Objects are
% **    placed in this box using the \put command and the like (see below).
% **    The location of an object is specified in terms of coordinate system(s)
% **    established by \setcoordinatesystem.  Each coordinate system (there
% **    might be just one) specifies the length of 1 horizontal unit, the
%%length
% **    of 1 vertical unit, and the coordinates of a "reference point".  The
% **    reference points of various coordinate systems will be in the same
% **    physical location.  The macros keep track of the size of the objects
% **    and their locations. The resulting hbox is the smallest hbox which
% **    encloses all the objects, and whose TEK reference point is the point
% **    on the left edge of the box closest vertically to the PICTEX reference
% **    point. Using \endpicturesave, you can (globally) save the distance
%%TEK's
% **    reference point is to the right (respectively, up from) PICTEX's
% **    reference point in the dimension register \XREG (respectively \YREG).
% **    You can then \put the picture OBJECT into a larger picture so that its
% **    reference point is at (XCOORD,YCOORD) with the command
% **      \put {picture OBJECT} [Bl] <\XREG, \YREG> at  XCOORD YCOORD

% **  \beginpicture : See Subsection 1.1 of the manual.
\def\beginpicture{%
  \setbox\!picbox=\hbox\bgroup%
  \!xleft=\maxdimen
  \!xright=-\maxdimen
  \!ybot=\maxdimen
  \!ytop=-\maxdimen}

% **  \endpicture : See Subsection 1.1 of the manual.
\def\endpicture{%
  \ifdim\!xleft=\maxdimen%  ** check if nothing was put in picbox
    \!xleft=\!zpt \!xright=\!zpt \!ybot=\!zpt \!ytop=\!zpt
  \fi
  \global\!Xleft=\!xleft \global\!Xright=\!xright
  \global\!Ybot=\!ybot \global\!Ytop=\!ytop
  \egroup%
  \ht\!picbox=\!Ytop  \dp\!picbox=-\!Ybot
  \ifdim\!Ybot>\!zpt
  \else
    \ifdim\!Ytop<\!zpt
      \!Ybot=\!Ytop
    \else
      \!Ybot=\!zpt
    \fi
  \fi
  \hbox{\kern-\!Xleft\lower\!Ybot\box\!picbox\kern\!Xright}}

% **  \endpicturesave : See Subsection 8.4 of the manual.
\def\endpicturesave <#1,#2>{%
  \endpicture \global #1=\!Xleft \global #2=\!Ybot \ignorespaces}


% **   \setcoordinatesystem units <XUNIT,YUNIT>
% **     point at XREF YREF
% **   Each of `units <XUNIT,YUNIT>' and `point at XREF YREF'
% **     are optional.
% **   Unit lengths must be given in dimensions (e.g., <10pt,1in>).
% **     Default unit lengths are 1pt, 1pt, or previous unit lengths.
% **   Reference point is specified in current units (e.g., 3 5 ).
% **     Default reference point is 0 0 , or previous reference point.
% **   Unit lengths and reference points obey TEX's scoping rules.
% **   See Subsection 1.2 of the manual.
\def\setcoordinatesystem{%
  \!ifnextchar{u}{\!getlengths }
    {\!getlengths units <\!xunit,\!yunit>}}
\def\!getlengths units <#1,#2>{%
  \!xunit=#1\relax
  \!yunit=#2\relax
  \!ifcoordmode
    \let\!SCnext=\!SCccheckforRP
  \else
    \let\!SCnext=\!SCdcheckforRP
  \fi
  \!SCnext}
\def\!SCccheckforRP{%
  \!ifnextchar{p}{\!cgetreference }
    {\!cgetreference point at {\!xref} {\!yref} }}
\def\!cgetreference point at #1 #2 {%
  \edef\!xref{#1}\edef\!yref{#2}%
  \!xorigin=\!xref\!xunit  \!yorigin=\!yref\!yunit
  \!initinboundscheck % ** See linear.tex
  \ignorespaces}
\def\!SCdcheckforRP{%
  \!ifnextchar{p}{\!dgetreference}%
    {\ignorespaces}}
\def\!dgetreference point at #1 #2 {%
  \!xorigin=#1\relax  \!yorigin=#2\relax
  \ignorespaces}


%  ** \put {OBJECT} [XY] <XDIMEN,YDIMEN> at (XCOORD,YCOORD)
%  **   `[XY]' and `<XDIMEN,YDIMEN>' are optional.
%  **   First OBJECT is placed in an hbox (the "objectbox") and then a
%  **     "reference point" is assigned to the objectbox as follows:
%  **     [1] first, the reference point is taken to be the center of the box;
%  **     [2] next, centering is overridden by the specifications
%  **           X=l -- reference point along the left edge of the objectbox
%  **           X=r -- reference point along the right edge of the objectbox
%  **           Y=b -- reference point along the bottom edge of the objectbox
%  **           Y=B -- reference point along the Baseline of the objectbox
%  **           Y=t -- reference point along the top edge of the objectbox;
%  **     [3] finally the reference point is shifted left by XDIMEN, down
%  **           by YDIMEN  (both default to 0pt).
%  **   The objectbox is placed within PICBOX with its reference point at
%  **     (XCOORD,YCOORD).
%  **   If OBJECT is a saved box, say  box0, you have to write
%  **     \put{\box0}...   or  \put{\copy0}...
%  **   The objectbox is void after the put.
%  **   See Subsection 2.1 of the manual.
\long\def\put#1#2 at #3 #4 {%
  \!setputobject{#1}{#2}%
  \!xpos=\!M{#3}\!xunit  \!ypos=\!M{#4}\!yunit
  \!rotateaboutpivot\!xpos\!ypos%
  \advance\!xpos -\!xorigin  \advance\!xpos -\!xshift
  \advance\!ypos -\!yorigin  \advance\!ypos -\!yshift
  \kern\!xpos\raise\!ypos\box\!putobject\kern-\!xpos%
  \!doaccounting\ignorespaces}

%  **   \multiput etc.  Like  \put.  The objectbox is not voided until the
%  **     termininating /, and is placed repeatedly with:
%  **     XCOORD YCOORD -- the objectbox is put down with its reference point
%  **       at (XCOORD,YCOORD);
%  **     *N DXCOORD DYCOORD -- each of N times the current
%  **       (xcoord,ycoord) is incremented by (DXCOORD,DYCOORD), and the
%  **       objectbox is put down with its reference point at (xcoord,ycoord)
%  **       (This specification has to follow an XCOORD YCOORD pair)
%  **     See Subsection 2.2 of the manual.
\long\def\multiput #1#2 at {%
  \!setputobject{#1}{#2}%
  \!ifnextchar"{\!putfromfile}{\!multiput}}
\def\!putfromfile"#1"{%
  \expandafter\!multiput \input #1 /}
\def\!multiput{%
  \futurelet\!nextchar\!!multiput}
\def\!!multiput{%
  \if *\!nextchar
    \def\!nextput{\!alsoby}%
  \else
    \if /\!nextchar
      \def\!nextput{\!finishmultiput}%
    \else
      \def\!nextput{\!alsoat}%
    \fi
  \fi
  \!nextput}
\def\!finishmultiput/{%
  \setbox\!putobject=\hbox{}%
  \ignorespaces}

%  **   \!alsoat XCOORD YCOORD
%  **     The objectbox is put down with reference point at XCOORD,YCOORD
\def\!alsoat#1 #2 {%
  \!xpos=\!M{#1}\!xunit  \!ypos=\!M{#2}\!yunit
  \!rotateaboutpivot\!xpos\!ypos%
  \advance\!xpos -\!xorigin  \advance\!xpos -\!xshift
  \advance\!ypos -\!yorigin  \advance\!ypos -\!yshift
  \kern\!xpos\raise\!ypos\copy\!putobject\kern-\!xpos%
  \!doaccounting
  \!multiput}

% **   \!alsoby*N DXCOORD DYCOORD
% **     N times, the current (XCOORD,YCOORD) is advanced by (DXCOORD,DYCOORD),
% **     and the current (shifted, oriented) OBJECT is put down.
\def\!alsoby*#1 #2 #3 {%
  \!dxpos=\!M{#2}\!xunit \!dypos=\!M{#3}\!yunit
  \!rotateonly\!dxpos\!dypos
  \!ntemp=#1%
  \!!loop\ifnum\!ntemp>0
    \advance\!xpos by \!dxpos  \advance\!ypos by \!dypos
    \kern\!xpos\raise\!ypos\copy\!putobject\kern-\!xpos%
    \advance\!ntemp by -1
  \repeat
  \!doaccounting
  \!multiput}

% **  \accountingoff : Suspends PiCTeX's accounting of the aggregate
% **    size of the picture box.
% **  \accounting on : Reinstates accounting.
% **  See Subsection 8.2 of the manual.
\def\accountingon{\def\!doaccounting{\!!doaccounting}\ignorespaces}
\def\accountingoff{\def\!doaccounting{}\ignorespaces}
\accountingon
\def\!!doaccounting{%
  \!xtemp=\!xpos
  \!ytemp=\!ypos
  \ifdim\!xtemp<\!xleft
     \!xleft=\!xtemp
  \fi
  \advance\!xtemp by  \!wd
  \ifdim\!xright<\!xtemp
    \!xright=\!xtemp
  \fi
  \advance\!ytemp by -\!dp
  \ifdim\!ytemp<\!ybot
    \!ybot=\!ytemp
  \fi
  \advance\!ytemp by  \!dp
  \advance\!ytemp by  \!ht
  \ifdim\!ytemp>\!ytop
    \!ytop=\!ytemp
  \fi}

\long\def\!setputobject#1#2{%
  \setbox\!putobject=\hbox{#1}%
  \!ht=\ht\!putobject  \!dp=\dp\!putobject  \!wd=\wd\!putobject
  \wd\!putobject=\!zpt
  \!xshift=.5\!wd   \!yshift=.5\!ht   \advance\!yshift by -.5\!dp
  \edef\!putorientation{#2}%
  \expandafter\!SPOreadA\!putorientation[]\!nil%
  \expandafter\!SPOreadB\!putorientation<\!zpt,\!zpt>\!nil\ignorespaces}

\def\!SPOreadA#1[#2]#3\!nil{\!etfor\!orientation:=#2\do\!SPOreviseshift}

\def\!SPOreadB#1<#2,#3>#4\!nil{\advance\!xshift by -#2\advance\!yshift by -#3}

\def\!SPOreviseshift{%
  \if l\!orientation
    \!xshift=\!zpt
  \else
    \if r\!orientation
      \!xshift=\!wd
    \else
      \if b\!orientation
        \!yshift=-\!dp
      \else
        \if B\!orientation
          \!yshift=\!zpt
        \else
          \if t\!orientation
            \!yshift=\!ht
          \fi
        \fi
      \fi
    \fi
  \fi}


%  **  \!dimenput{OBJECT} <XDIMEN,YDIMEN> [XY] (XLOC,YLOC)
%  **    This is an internal put routine, similar to \put, except that
%  **    XLOC=distance right from reference point, YLOC=distance up from
%  **    reference point. XLOC and YLOC are dimensions, so this routine
%  **    is completely independent of the current coordinate system.
%  **    This routine does NOT do ROTATIONS.
\long\def\!dimenput#1#2(#3,#4){%
  \!setputobject{#1}{#2}%
  \!xpos=#3\advance\!xpos by -\!xshift
  \!ypos=#4\advance\!ypos by -\!yshift
  \kern\!xpos\raise\!ypos\box\!putobject\kern-\!xpos%
  \!doaccounting\ignorespaces}


%  ** The following macros permit the picture drawing routines to be used
%  ** either in the default "coordinate mode", or in "dimension mode".
%  **   In coordinate mode  \!M(1.5,\!xunit)    expands to  1.5\!xunit
%  **   In dimension  mode  \!M(1.5pt,\!xunit)  expands to  1.5pt
%  ** Dimension mode is useful in coding macros.
%  ** Any special purpose picture macro that sets dimension mode should
%  ** reset coordinate mode before completion.
%  ** See Subsection 9.2 of the manual.
\def\!setdimenmode{%
  \let\!M=\!M!!\ignorespaces}
\def\!setcoordmode{%
  \let\!M=\!M!\ignorespaces}
\def\!ifcoordmode{%
  \ifx \!M \!M!}
\def\!ifdimenmode{%
  \ifx \!M \!M!!}
\def\!M!#1#2{#1#2}
\def\!M!!#1#2{#1}
\!setcoordmode
\let\setdimensionmode=\!setdimenmode
\let\setcoordinatemode=\!setcoordmode

%  ** \Xdistance{XCOORD}, \Ydistance{YCOORD}  are the horizontal and
%  **   vertical distances from the origin (0,0) to the point
%  **   (XCOORD,YCOORD)  in the current coordinate system.
%  ** See Subsection 9.2 of the manual.
\def\Xdistance#1{%
  \!M{#1}\!xunit
  \ignorespaces}
\def\Ydistance#1{%
  \!M{#1}\!yunit
  \ignorespaces}

% ** The following macros -- \stack, \line, and \Lines -- are useful for
% **   annotating PiCtures. They can be used outside the \beginpicture ...
% **   \endpicture environment.

% ** \stack [POSITIONING] <LEADING> {VALUESLIST}
% ** Builds a vertical stack of the values in VALUESLIST. Values in
% ** VALUESLIST are separated by commas.  In the resulting stack, values are
% ** centered by default, and positioned flush left (right) if
% ** POSITIONING = l (r).  Values are separated vertically by LEADING,
% ** which defaults to \stackleading.
% ** See Subsection 2.3 of the manual.
\def\stack{%
  \!ifnextchar[{\!stack}{\!stack[c]}}
\def\!stack[#1]{%
  \let\!lglue=\hfill \let\!rglue=\hfill
  \expandafter\let\csname !#1glue\endcsname=\relax
  \!ifnextchar<{\!!stack}{\!!stack<\stackleading>}}
\def\!!stack<#1>#2{%
  \vbox{\def\!valueslist{}\!ecfor\!value:=#2\do{%
    \expandafter\!rightappend\!value\withCS{\\}\to\!valueslist}%
    \!lop\!valueslist\to\!value
    \let\\=\cr\lineskiplimit=\maxdimen\lineskip=#1%
    \baselineskip=-1000pt\halign{\!lglue##\!rglue\cr \!value\!valueslist\cr}}%
  \ignorespaces}

% ** \lines [POSITIONING] {LINES}
% ** Builds a vertical array of the lines in LINES. Each line in LINES
% ** is terminated by a \cr.  In the resulting array, lines are
% ** centered by default, and positioned flush left (right) if
% ** POSITIONING = l (r).  The lines in the array are subject to TeX's
% ** usual spacing rules: in particular the baselines are ordinarily an equal
% ** distance apart. The baseline of the array is the baseline of the
% ** the bottom line.
% ** See Subsection 2.3 of the manual.
\def\lines{%
  \!ifnextchar[{\!lines}{\!lines[c]}}
\def\!lines[#1]#2{%
  \let\!lglue=\hfill \let\!rglue=\hfill
  \expandafter\let\csname !#1glue\endcsname=\relax
  \vbox{\halign{\!lglue##\!rglue\cr #2\crcr}}%
  \ignorespaces}

% ** \Lines [POSITIONING] {LINES}
% ** Like \lines, but the baseline of the array is the baseline of the
% ** top line.  See Subsection 2.3 of the manual.
\def\Lines{%
  \!ifnextchar[{\!Lines}{\!Lines[c]}}
\def\!Lines[#1]#2{%
  \let\!lglue=\hfill \let\!rglue=\hfill
  \expandafter\let\csname !#1glue\endcsname=\relax
  \vtop{\halign{\!lglue##\!rglue\cr #2\crcr}}%
  \ignorespaces}


% *********************************************
% *** PLOTTING (Things to do with plotting) ***
% *********************************************

% **  User commands
% **    \setplotsymbol ({PLOTSYMBOL} [ORIENTATION] <XSHIFT,YSHIFT>)
% **    \savelinesandcurves on "FILE_NAME"
% **    \dontsavelinesandcurves
% **    \writesavefile {MESSAGE}
% **    \replot {FILE_NAME}

% **  Internal command
% **    \!plot(XDIMEN,YDIMEN)

% **  \setplotsymbol ({PLOTSYMBOL} [ ] < , >)
% **  Save PLOTSYMBOL away in an hbox for use with curve plotting routines
% **  See Subsection 5.2 of the manual.
\def\setplotsymbol(#1#2){%
  \!setputobject{#1}{#2}
  \setbox\!plotsymbol=\box\!putobject%
  \!plotsymbolxshift=\!xshift
  \!plotsymbolyshift=\!yshift
  \ignorespaces}

\setplotsymbol({\fiverm .})%       ** initialize plotsymbol


% **  \!plot is either \!!plot (when no lines and curves are being saved) or
% **                   \!!!plot (when   lines and curves are being saved)

% **  \!!plot(XDIMEN,YDIMEN)
% **  Places the current plotsymbol a horizontal distance=XDIMEN-xorigin
% **    and a vertical distance=YDIMEN-yorigin from the current
% **    reference point.
\def\!!plot(#1,#2){%
  \!dimenA=-\!plotxorigin \advance \!dimenA by #1%    ** over
  \!dimenB=-\!plotyorigin \advance \!dimenB by #2%    ** up
  \kern\!dimenA\raise\!dimenB\copy\!plotsymbol\kern-\!dimenA%
  \ignorespaces}

% **  \!!!plot(XDIMEN,YDIMEN)
% **  Like \!!plot, but also saves the plot location in units of
% **    scaled point, on file `replotfile'
\def\!!!plot(#1,#2){%
  \!dimenA=-\!plotxorigin \advance \!dimenA by #1%    ** over
  \!dimenB=-\!plotyorigin \advance \!dimenB by #2%    ** up
  \kern\!dimenA\raise\!dimenB\copy\!plotsymbol\kern-\!dimenA%
  \!countE=\!dimenA
  \!countF=\!dimenB
  \immediate\write\!replotfile{\the\!countE,\the\!countF.}%
  \ignorespaces}


% ** \savelinesandcurves on "FILE_NAME"
% **   Switch to save locations used for plotting lines and curves
% **   (No advantage in saving locations for solid lines; however
% **   replotting curve locations speeds things up by a factor of about 4.
% ** \dontsavelinesandcurves
% **   Terminates \savelinesandcurves. The default.
% ** See Subsection 5.6 of the manual.
\def\savelinesandcurves on "#1" {%
  \immediate\closeout\!replotfile
  \immediate\openout\!replotfile=#1%
  \let\!plot=\!!!plot}

\def\dontsavelinesandcurves {%
  \let\!plot=\!!plot}
\dontsavelinesandcurves

% ** \writesavefile {MESSAGE}
% ** The message is preceded by a "%", so that it won't interfere
% ** with replotting.
% ** See Subsection 5.6 of the manual.
{\catcode`\%=11\xdef\!Commentsignal{%}}
\def\writesavefile#1 {%
  \immediate\write\!replotfile{\!Commentsignal #1}%
  \ignorespaces}

% ** \replot "FILE_NAME"
% **   Replots the locations saved earlier under \savelinesandcurves
% **   on "FILE_NAME"
% ** See Subsection 5.6 of the manual.
\def\replot"#1" {%
  \expandafter\!replot\input #1 /}
\def\!replot#1,#2. {%
  \!dimenA=#1sp
  \kern\!dimenA\raise#2sp\copy\!plotsymbol\kern-\!dimenA
  \futurelet\!nextchar\!!replot}
\def\!!replot{%
  \if /\!nextchar
    \def\!next{\!finish}%
  \else
    \def\!next{\!replot}%
  \fi
  \!next}
% **************************************************
% ***  PYTHAGORAS  (Euclidean distance function) ***
% **************************************************

% ** User command:
% **   \placehypotenuse for <dimension1> and <dimension2> in <register>

% ** Internal command:
% **   \!Pythag{X}{Y}{Z}
% **     Input X,Y are dimensions, or dimension registers.
% **     Output Z == sqrt(X**2+Y**2) must be a dimension register.
% **     Assumes that |X|+|Y| < 2048pt (about 28in).

% ** Without loss of generality, suppose  x>0, y>0.  Put s = x+y,
% **   z = sqrt(x**2+y**2). Then  z = s*f,  where  f = sqrt(t**2 + (1-t)**2)
% **   = sqrt((1+tau**2)/2), where  t = x/s  and  tau = 2(t-1/2) .

% ** Uses the \!divide macro (which uses registers \!dimenA--\!dimenD.
% ** Uses the \!removept macro   (e.g., 123.45pt --> 123.45)
% ** Uses registers \!dimenE--\!dimenI.
\def\!Pythag#1#2#3{%
  \!dimenE=#1\relax
  \ifdim\!dimenE<\!zpt
    \!dimenE=-\!dimenE
  \fi%                                            ** dimE = |x|
  \!dimenF=#2\relax
  \ifdim\!dimenF<\!zpt
    \!dimenF=-\!dimenF
  \fi%                                            ** dimF = |y|
  \advance \!dimenF by \!dimenE%                  ** dimF = s = |x|+|y|
  \ifdim\!dimenF=\!zpt
    \!dimenG=\!zpt%                               ** dimG = z = sqrt(x**2+y**2)
  \else
    \!divide{8\!dimenE}\!dimenF\!dimenE%          ** now dimE = 8t = (8|x|)/s
    \advance\!dimenE by -4pt%                     ** 8tau = (8t-4)*2
      \!dimenE=2\!dimenE%                         **   (tau = 2*t - 1)
    \!removept\!dimenE\!!t%                       ** 8tau, without "pt"
    \!dimenE=\!!t\!dimenE%                        ** (8tau)**2, in pts
    \advance\!dimenE by 64pt%                     ** u = [64 + (8tau)**2]/2
    \divide \!dimenE by 2%                        **   [u = (8f)**2]
    \!dimenH=7pt%                                 ** initial guess g at sqrt(u)
    \!!Pythag\!!Pythag\!!Pythag%                  ** 3 iterations give sqrt(u)
    \!removept\!dimenH\!!t%                       ** 8f=sqrt(u), without "pt"
    \!dimenG=\!!t\!dimenF%                        ** z = (8f)*s/8
    \divide\!dimenG by 8
  \fi
  #3=\!dimenG
  \ignorespaces}

\def\!!Pythag{%                                   ** Newton-Raphson for sqrt
  \!divide\!dimenE\!dimenH\!dimenI%               ** v = u/g
  \advance\!dimenH by \!dimenI%                   ** g <-- (g + u/g)/2
    \divide\!dimenH by 2}

% **  \placehypotenuse for <XI> and <ETA> in <ZETA>
% **  See Subsection 9.3 of the manual.
\def\placehypotenuse for <#1> and <#2> in <#3> {%
  \!Pythag{#1}{#2}{#3}}


% **********************************************
% *** QUADRATIC ARC  (Draws a quadratic arc) ***
% **********************************************

% **  Internal command
% **    \!qjoin (XCOORD1,YCOORD1) (XCOORD2,YCOORD2)

% **  \!qjoin (XCOORD1,YCOORD1) (XCOORD2,YCOORD2)
% **  Draws an arc starting at the (last) point specified by the most recent
% **  \!qjoin, or \!ljoin, or \!start  and passing through (X_1,Y_1),
%%(X_2,Y_2).
% **  Uses quadratic interpolation in both  x  and  y:
% **    x(t), 0 <= t <= 1, interpolates  x_0, x_1, x_2  at  t=0, .5, 1
% **    y(t), 0 <= t <= 1, interpolates  y_0, y_1, y_2  at  t=0, .5, 1

\def\!qjoin (#1,#2) (#3,#4){%
  \advance\!intervalno by 1
  \!ifcoordmode
    \edef\!xmidpt{#1}\edef\!ymidpt{#2}%
  \else
    \!dimenA=#1\relax \edef\!xmidpt{\the\!dimenA}%
    \!dimenA=#2\relax \edef\!xmidpt{\the\!dimenA}%
  \fi
  \!xM=\!M{#1}\!xunit  \!yM=\!M{#2}\!yunit   \!rotateaboutpivot\!xM\!yM
  \!xE=\!M{#3}\!xunit  \!yE=\!M{#4}\!yunit   \!rotateaboutpivot\!xE\!yE
%
% ** Find coefficients for x(t)=a_x + b_x*t + c_x*t**2
  \!dimenA=\!xM  \advance \!dimenA by -\!xS%   ** dimA = I = xM - xS
  \!dimenB=\!xE  \advance \!dimenB by -\!xM%   ** dimB = II = xE-xM
  \!xB=3\!dimenA \advance \!xB by -\!dimenB%   ** b=3I-II
  \!xC=2\!dimenB \advance \!xC by -2\!dimenA%  ** c=2(II-I)
%
% ** Find coefficients for y(t)=y_x + b_y*t + c_y*t**2
  \!dimenA=\!yM  \advance \!dimenA by -\!yS%
  \!dimenB=\!yE  \advance \!dimenB by -\!yM%
  \!yB=3\!dimenA \advance \!yB by -\!dimenB%
  \!yC=2\!dimenB \advance \!yC by -2\!dimenA%
%
% ** Use Simpson's rule to calculate arc length over [0,1/2]:
% **   arc length = 1/2[1/6 f(0) + 4/6 f(1/4) + 1/6 f(1/2)]
% ** with f(t) = sqrt(x'(t)**2 + y'(t)**2).
  \!xprime=\!xB  \!yprime=\!yB%          ** x'(t) = b + 2ct
  \!dxprime=.5\!xC  \!dyprime=.5\!yC%    ** dt=1/4 ==> dx'(t) = c/2
  \!getf \!midarclength=\!dimenA
  \!getf \advance \!midarclength by 4\!dimenA
  \!getf \advance \!midarclength by \!dimenA
  \divide \!midarclength by 12
%
% ** Get arc length over [0,1].
  \!arclength=\!dimenA
  \!getf \advance \!arclength by 4\!dimenA
  \!getf \advance \!arclength by \!dimenA
  \divide \!arclength by 12%             ** Now have arc length over [1/2,1]
  \advance \!arclength by \!midarclength
  \global\advance \totalarclength by \!arclength
%
%
% ** Check to see if there's anything to plot in this interval
  \ifdim\!distacross>\!arclength
    \advance \!distacross by -\!arclength%   ** nothing
%
  \else
    \!initinverseinterp%  ** initialize for inverse interpolation on arc length
    \loop\ifdim\!distacross<\!arclength%     ** loop over points on arc
      \!inverseinterp%    ** find  t  such that arc length[0,t] = distacross,
%                         **   using inverse quadratic interpolation
%                         ** now evaluate x(t)=(c*t + b)*t + a
      \!xpos=\!t\!xC \advance\!xpos by \!xB
        \!xpos=\!t\!xpos \advance \!xpos by \!xS
%                                             ** evaluate y(t)
      \!ypos=\!t\!yC \advance\!ypos by \!yB
        \!ypos=\!t\!ypos \advance \!ypos by \!yS
      \!plotifinbounds%                       ** plot point if in bounds
      \advance\!distacross \plotsymbolspacing%** advance arc length for next pt
      \!advancedashing%                       ** see "linear"
    \repeat
%
    \advance \!distacross by -\!arclength%    ** prepare for next interval
  \fi
%
  \!xS=\!xE%              ** shift ending points to starting points
  \!yS=\!yE
  \ignorespaces}


% ** \!getf -- Calculates sqrt(x'(t)**2 + y'(t)**2) and advances
% **   x'(t) and y'(t)
\def\!getf{\!Pythag\!xprime\!yprime\!dimenA%
  \advance\!xprime by \!dxprime
  \advance\!yprime by \!dyprime}


% ** \!initinverseinterp -- initializes for inverse quadratic interpolation
% ** of arc length provided  1/3 < midarclength/arclength < 2/3; otherwise
% ** initializes for inverse linear interpolation.
\def\!initinverseinterp{%
  \ifdim\!arclength>\!zpt
    \!divide{8\!midarclength}\!arclength\!dimenE% ** dimE=8w=8r/s, where  r
%                                               **  = midarclength, s=arclength
% **  Test for  w  out of range:  w<1/3  or w>2/3
    \ifdim\!dimenE<\!wmin \!setinverselinear
    \else
      \ifdim\!dimenE>\!wmax \!setinverselinear
      \else%                                    ** w  in range: initialize
        \def\!inverseinterp{\!inversequad}\ignorespaces
%
% **     Calculate the coefficients  \!beta  and  \!gamma  of the quadratic
% **                    t = \!beta*v + \!gamma*v**2
% **     taking the values  t=0, 1/2, 1  at  v=0, w==r/s, 1  respectively:
% **        \!beta = (1/2 - w**2)/[w(1-w)]
% **        \!gamma = 1 - beta.
%
         \!removept\!dimenE\!Ew%           **  8w, without "pt"
         \!dimenF=-\!Ew\!dimenE%           **  -(8w)**2
         \advance\!dimenF by 32pt%         **  32 - (8w)**2
         \!dimenG=8pt
         \advance\!dimenG by -\!dimenE%    **  8 - 8w
         \!dimenG=\!Ew\!dimenG%            **  (8w)*(8-8w)
         \!divide\!dimenF\!dimenG\!beta%   **  beta = (32-(8w)**2)/(8w(8-8w))
%                                          **       = (1/2 - w**2)/(w(1-w))
         \!gamma=1pt
         \advance \!gamma by -\!beta%      **  gamma = 1-beta
      \fi%       ** end of the \ifdim\!dimenE>\!wmax
    \fi%         ** end of the \ifdim\!dimenE<\!wmin
  \fi%           ** end of the \ifdim\!arclength>\!zpt
  \ignorespaces}


% ** For 0 <= t <= 1, let AL(t) = arclength[0,t]/arclength[0,1]; note
% ** AL(0)=0, AL(1/2)=midarclength/arclength, AL(1)=1.  This routine
% ** calculates an approximation to AL^{-1}(distance across/arclength),
% ** using the assumption that AL^{-1} is quadratic.  Specifically,
% ** it finds  t  such that
% **    AL^{-1}(v) =. t = v*(\!beta + \!gamma*v)
% ** where  \!beta  and  \!gamma  are set by \!initinv, and where
% ** v=distance across/arclength
\def\!inversequad{%
  \!divide\!distacross\!arclength\!dimenG%   ** dimG = v = distacross/arclength
  \!removept\!dimenG\!v%                     ** v, without "pt"
  \!dimenG=\!v\!gamma%                       ** gamma*v
  \advance\!dimenG by \!beta%                ** beta + gamma*v
  \!dimenG=\!v\!dimenG%                      ** t = v*(beta + gamma*v)
  \!removept\!dimenG\!t}%                    ** t, without "pt"


% ** When  w <= 1/3  or  w >= 2/3, the following routine writes (using
% ** plain TEK's \wlog command) a warning message on the user's log file,
% ** and initializes for inverse linear interpolation on arc length.
\def\!setinverselinear{%
  \def\!inverseinterp{\!inverselinear}%
  \divide\!dimenE by 8 \!removept\!dimenE\!t
  \!countC=\!intervalno \multiply \!countC 2
  \!countB=\!countC     \advance \!countB -1
  \!countA=\!countB     \advance \!countA -1
  \wlog{\the\!countB th point (\!xmidpt,\!ymidpt) being plotted
    doesn't lie in the}%
  \wlog{ middle third of the arc between the \the\!countA th
    and \the\!countC th points:}%
  \wlog{ [arc length \the\!countA\space to \the\!countB]/[arc length
    \the \!countA\space to \the\!countC]=\!t.}%
  \ignorespaces}

% **  Inverse linear interpolation
\def\!inverselinear{%
  \!divide\!distacross\!arclength\!dimenG
  \!removept\!dimenG\!t}


% **************************************
% **  ROTATIONS  (Handles rotations) ***
% **************************************

% ** User commands
% **   \startrotation [by COS_OF_ANGLE SIN_OF_ANGLE] [about XPIVOT YPIVOT]
% **   \stoprotation

% **   \startrotation [by COS_OF_ANGLE SIN_OF_ANGLE] [about XPIVOT YPIVOT]
% ** Future (XCOORD,YCOORD)'s will be rotated about (XPIVOT,YPIVOT)
% ** by the angle with the give COS and SIN. Both fields are optional.
% ** [COS,SIN] defaults to previous value, or (1,0).
% ** (XPIVOT,YPIVOT) defaults to previous value, or (0,0)
% ** You can't change the coordinate system in the scope of a rotation.
% ** See Subsection 9.1 of the manual.
\def\startrotation{%
  \let\!rotateaboutpivot=\!!rotateaboutpivot
  \let\!rotateonly=\!!rotateonly
  \!ifnextchar{b}{\!getsincos }%
    {\!getsincos by {\!cosrotationangle} {\!sinrotationangle} }}
\def\!getsincos by #1 #2 {%
  \edef\!cosrotationangle{#1}%
  \edef\!sinrotationangle{#2}%
  \!ifcoordmode
    \let\!ROnext=\!ccheckforpivot
  \else
    \let\!ROnext=\!dcheckforpivot
  \fi
  \!ROnext}
\def\!ccheckforpivot{%
  \!ifnextchar{a}{\!cgetpivot}%
    {\!cgetpivot about {\!xpivotcoord} {\!ypivotcoord} }}
\def\!cgetpivot about #1 #2 {%
  \edef\!xpivotcoord{#1}%
  \edef\!ypivotcoord{#2}%
  \!xpivot=#1\!xunit  \!ypivot=#2\!yunit
  \ignorespaces}
\def\!dcheckforpivot{%
  \!ifnextchar{a}{\!dgetpivot}{\ignorespaces}}
\def\!dgetpivot about #1 #2 {%
  \!xpivot=#1\relax  \!ypivot=#2\relax
  \ignorespaces}


% ** Following terminates rotation.
% ** See Subsection 9.1 of the manual.
\def\stoprotation{%
  \let\!rotateaboutpivot=\!!!rotateaboutpivot
  \let\!rotateonly=\!!!rotateonly
  \ignorespaces}

% ** !!rotateaboutpivot{XREG}{YREG}
% ** XREG <-- xpvt + cos(angle)*(XREG-xpvt) - sin(angle)*(YREG-ypvt)
% ** YREG <-- ypvt + cos(angle)*(YREG-ypvt) + sin(angle)*(XREG-xpvt)
% ** XREG,YREG are dimension registers. Can't be \!dimenA to \!dimenD
\def\!!rotateaboutpivot#1#2{%
  \!dimenA=#1\relax  \advance\!dimenA -\!xpivot
  \!dimenB=#2\relax  \advance\!dimenB -\!ypivot
  \!dimenC=\!cosrotationangle\!dimenA
    \advance \!dimenC -\!sinrotationangle\!dimenB
  \!dimenD=\!cosrotationangle\!dimenB
    \advance \!dimenD  \!sinrotationangle\!dimenA
  \advance\!dimenC \!xpivot  \advance\!dimenD \!ypivot
  #1=\!dimenC  #2=\!dimenD
  \ignorespaces}

% ** \!!rotateonly{XREG}{YREG}
% ** Like \!!rotateaboutpivot, but with a pivot of  (0,0)
\def\!!rotateonly#1#2{%
  \!dimenA=#1\relax  \!dimenB=#2\relax
  \!dimenC=\!cosrotationangle\!dimenA
    \advance \!dimenC -\!rotsign\!sinrotationangle\!dimenB
  \!dimenD=\!cosrotationangle\!dimenB
    \advance \!dimenD  \!rotsign\!sinrotationangle\!dimenA
  #1=\!dimenC  #2=\!dimenD
  \ignorespaces}
\def\!rotsign{}
\def\!!!rotateaboutpivot#1#2{\relax}
\def\!!!rotateonly#1#2{\relax}
\stoprotation

\def\!reverserotateonly#1#2{%
  \def\!rotsign{-}%
  \!rotateonly{#1}{#2}%
  \def\!rotsign{}%
  \ignorespaces}


% **********************************
% *** SHADING  (Handles shading) ***
% **********************************

% **  User commands
% **    \setshadegrid [span <SPAN>] [point at XSHADE YSHADE]
% **    \setshadesymbol [<LS, RS, BS, TS>] ({SHADESYMBOL}
% **      <XDIMEN,YDIMEN> [ORIENTATION])

% **  Internal commands:
% **    \!startvshade  (xS,ybS,ytS)
% **    \!starthshade  (yS,xlS,xrS)
% **    \!lshade [<LS,RS,BS,TS>]
% **       ** when shading vertically:
% **       [the region from (xS,ybS,ytS) to] (xE,ybE,ytE)
% **       ** when shading horizontally:
% **       [the region from (yS,xlS,xrS) to] (yE,xlE,xrE)
% **    \!qshade [<LS,RS,BS,TS>]
% **       ** when shading vertically:
% **       [the region from (xS,ybS,ytS) to] (xM,ybM,ytM)  (xE,ybE,ytE)
% **       ** when shading horizontally:
% **       [the region from (yS,xlS,xrS) to] (yM,xlM,xrM)  (yE,xlE,xrE)
% **    \!lattice{ANCHOR}{SPAN}{LOCATION}{INDEX}{LATTICE LOCATION}
% **    \!override{NOMINAL DIMEN}{REPLACEMENT DIMEN}{DIMEN}


% **  The shading routine can operate either in a "vertical mode" or a
% **  "horizontal mode".  In vertical mode, the region to be shaded is
%%specified
% **  in the form
% **                 {(x,y): xl <= x <= xr  &  yb(x) <= y <= yt(x)}
% **  where  yb  and  yt  are functions of  x.  In horizontal mode, the region
% **  is specified in the form
% **                 {(x,y): yb <= y <= yt  &  xl(y) <= x <= xr(y)}.
% **  The functions  yb  and  yt  may be either both linear or both quadratic;
% **  similarly for  xl  and  xr.  A region with say, piecewise quadratic
%%bottom
% **  and top boundaries, can be shaded by consecutive (vertical) \!qshades,
% **  proceeding from left to right.  Similarly, a region with piecewise
% **  quadratic left and right boundaries can be shaded by consecutive
% **  (horizontal) \!qshades, proceeding from bottom to top.  More complex
% **  regions can be shaded by partitioning them into appropriate subregions,
% **  and shading those.

% **  Shading is accomplished by placing a user-selected shading symbol at
% **  those points of a regular grid which fall within the region to be
% **  shaded.  This region can be "shrunk" so that a largish shading symbol
% **  will not extend outside it.  Shrinking is accomplished by specifying
% **  shrinkages for the left, right, bottom, and top boundaries, in a manner
% **  discussed further below.

% **  \shades and \!joins MUST NOT be intermingled.  Finish drawing a curve
% **  before starting to shade a region, and finish shading a region before
% **  starting to draw a curve.


% **  \setshadegrid [span <SPAN>] [point at XSHADE YSHADE]
% **  The shading symbol is placed down on the points of a grid centered
% **  at the coordinate point (XSHADE,YSHADE).  The grid points are of the
% **  form (j*SPAN,k*SPAN), with  j+k  even.  SPAN is specified
% **  as a dimension.
% **  (XSHADE,YSHADE) defaults to previous (XSHADE,YSHADE) (or (0,0) if none)
% **  SPAN defaults to previous span (or 5pt if none)
% **  See Subsection 7.2 of the manual.
\def\setshadegrid{%
  \!ifnextchar{s}{\!getspan }
    {\!getspan span <\!dshade>}}
\def\!getspan span <#1>{%
  \!dshade=#1\relax
  \!ifcoordmode
    \let\!GRnext=\!GRccheckforAP
  \else
    \let\!GRnext=\!GRdcheckforAP
  \fi
  \!GRnext}
\def\!GRccheckforAP{%
  \!ifnextchar{p}{\!cgetanchor }
    {\!cgetanchor point at {\!xshadesave} {\!yshadesave} }}
\def\!cgetanchor point at #1 #2 {%
  \edef\!xshadesave{#1}\edef\!yshadesave{#2}%
  \!xshade=\!xshadesave\!xunit  \!yshade=\!yshadesave\!yunit
  \ignorespaces}
\def\!GRdcheckforAP{%
  \!ifnextchar{p}{\!dgetanchor}%
    {\ignorespaces}}
\def\!dgetanchor point at #1 #2 {%
  \!xshade=#1\relax  \!yshade=#2\relax
  \ignorespaces}

% **  \setshadesymbol  [<LS, RS, BS, TS>] ({SHADESYMBOL}
% **    <XDIMEN,YDIMEN> [ORIENTATION])
% **  Saves SHADESYMBOL away in an hbox for use with shading routines.
% **  A shade symbol will not be plotted if its plot position comes within
% **    distance LS of the left boundary,  RS of the right boundary,  TS of the
% **    top boundary,  BS of the bottom boundary.  These parameters have
% **    default values that should work in most cases (see below).
% **    To override a default value, specify the replacement value
% **    in the appropriate subfield of the shrinkages field.
% **    0pt may be coded as  "z" (without the quotes).  To accept a
% **    default value, leave the field empty.  Thus
% **      [,z,,5pt]  sets  LS=default, RS=0pt, BS=default, TS=5pt .
% **    Skipping the shrinkages field accepts all the defaults.
% **  See Subsection 7.1 of the manual.
\def\setshadesymbol{%
  \!ifnextchar<{\!setshadesymbol}{\!setshadesymbol<,,,> }}

\def\!setshadesymbol <#1,#2,#3,#4> (#5#6){%
% **  set the shadesymbol
  \!setputobject{#5}{#6}%
  \setbox\!shadesymbol=\box\!putobject%
  \!shadesymbolxshift=\!xshift \!shadesymbolyshift=\!yshift
%
% **  set the shrinkages
  \!dimenA=\!xshift \advance\!dimenA \!smidge% ** default LS = xshift - smidge
  \!override\!dimenA{#1}\!lshrinkage%
  \!dimenA=\!wd \advance \!dimenA -\!xshift%   ** default RS = width - xshift
    \advance\!dimenA \!smidge%                                  - smidge
    \!override\!dimenA{#2}\!rshrinkage
  \!dimenA=\!dp \advance \!dimenA \!yshift%    ** default BS = depth + yshift
    \advance\!dimenA \!smidge%                                  - smidge
    \!override\!dimenA{#3}\!bshrinkage
  \!dimenA=\!ht \advance \!dimenA -\!yshift%   ** default TS = height - yshift
    \advance\!dimenA \!smidge%                                  - smidge
    \!override\!dimenA{#4}\!tshrinkage
  \ignorespaces}
\def\!smidge{-.2pt}%

% ** \!override{NOMINAL DIMEN}{REPLACEMENT DIMEN}{DIMEN}
% ** Overrides the NOMINAL DIMEN by the REPLACEMENT DIMEN to produce DIMEN,
% ** according to the following rules:
% **   REPLACEMENT DIMEN empty: DIMEN <-- NOMINAL DIMEN
% **   REPLACEMENT DIMEN z:     DIMEN <-- 0pt
% **   otherwise:               DIMEN <-- REPLACEMENT DIMEN
% ** DIMEN must be a dimension register
\def\!override#1#2#3{%
  \edef\!!override{#2}%
  \ifx \!!override\empty
    #3=#1\relax
  \else
    \if z\!!override
      #3=\!zpt
    \else
      \ifx \!!override\!blankz
        #3=\!zpt
      \else
        #3=#2\relax
      \fi
    \fi
  \fi
  \ignorespaces}
\def\!blankz{ z}

\setshadesymbol ({\fiverm .})%       ** initialize plotsymbol
%                                    ** \fivesy ^^B  is a small cross


% ** \!startvshade [at] (xS,ybS,ytS)
% ** Initiates vertical shading mode
\def\!startvshade#1(#2,#3,#4){%
  \let\!!xunit=\!xunit%
  \let\!!yunit=\!yunit%
  \let\!!xshade=\!xshade%
  \let\!!yshade=\!yshade%
  \def\!getshrinkages{\!vgetshrinkages}%
  \let\!setshadelocation=\!vsetshadelocation%
  \!xS=\!M{#2}\!!xunit
  \!ybS=\!M{#3}\!!yunit
  \!ytS=\!M{#4}\!!yunit
  \!shadexorigin=\!xorigin  \advance \!shadexorigin \!shadesymbolxshift
  \!shadeyorigin=\!yorigin  \advance \!shadeyorigin \!shadesymbolyshift
  \ignorespaces}

% ** \!starthshade [at] (yS,xlS,xrS)
% ** Initiates horizontal shading mode
\def\!starthshade#1(#2,#3,#4){%
  \let\!!xunit=\!yunit%
  \let\!!yunit=\!xunit%
  \let\!!xshade=\!yshade%
  \let\!!yshade=\!xshade%
  \def\!getshrinkages{\!hgetshrinkages}%
  \let\!setshadelocation=\!hsetshadelocation%
  \!xS=\!M{#2}\!!xunit
  \!ybS=\!M{#3}\!!yunit
  \!ytS=\!M{#4}\!!yunit
  \!shadexorigin=\!xorigin  \advance \!shadexorigin \!shadesymbolxshift
  \!shadeyorigin=\!yorigin  \advance \!shadeyorigin \!shadesymbolyshift
  \ignorespaces}


% **  \!lattice{ANCHOR}{SPAN}{LOCATION}{INDEX}{LATTICE LOCATION}
% **  Consider the lattice with points  ANCHOR + j*SPAN. This routine
%%determines
% **  the index  k  of the smallest lattice point >= LOCATION, and sets
% **  LATTICE LOCATION = ANCHOR + k*SPAN.
% **  INDEX is assumed to be a count register, LATTICE LOCATION a dimen reg.
\def\!lattice#1#2#3#4#5{%
  \!dimenA=#1%                        ** dimA = ANCHOR
  \!dimenB=#2%                        ** dimB = SPAN  (assumed > 0pt)
  \!countB=\!dimenB%                  ** ctB  = SPAN, as a count
%
% ** Determine index of smallest lattice point >= LOCATION
  \!dimenC=#3%                        ** dimC = LOCATION
  \advance\!dimenC -\!dimenA%         ** now dimC = LOCATION-ANCHOR
  \!countA=\!dimenC%                  ** ctA = above, as a count
  \divide\!countA \!countB%           ** now ctA = desired index, if dimC <= 0
  \ifdim\!dimenC>\!zpt
    \!dimenD=\!countA\!dimenB%        ** (tentative k)*span
    \ifdim\!dimenD<\!dimenC%          ** if this is false, ctA = desired index
      \advance\!countA 1 %            ** if true, have to add 1
    \fi
  \fi
%
  \!dimenC=\!countA\!dimenB%          ** lattice location = anchor + ctA*span
    \advance\!dimenC \!dimenA
  #4=\!countA%                        ** the desired index
  #5=\!dimenC%                        ** corresponding lattice location
  \ignorespaces}


% ** \!qshade [with shrinkages] [[LS,RS,BS,TS]]
% ***** during vertical shading:
% **    [the region from (xS,ybS,ytS) to] (xM,ybM,ytM) [and] (xE,ybE,ytE)
% ** Shades the region {(x,y): xS <= x <= xE, yb(x) <= y <= yt(x)}, where
% **   yb is the quadratic thru (xS,ybS) & (xM,ybM) & (xE,ybE)
% **   yt is the quadratic thru (xS,ytS) & (xM,ybM) & (xE,ytE)
% ** xS,ybS,ytS are either given by \!startvshade or carried over
% **   as the ending values of the immediately preceding \!qshade.
% ** For the interpretation of LS, RS, BS, & TS, see \setshadesymbol. The
% **   values set there can be overridden, for the course of this \!qshade
% **   only, in the same manner as overrides are specified for
% **   \setshadesymbol.
% ***** during horizontal shading:
% **    [the region from (yS,xlS,xrS) to] (yM,xlM,xrM) [and] (yE,xlE,xrE)
\def\!qshade#1(#2,#3,#4)#5(#6,#7,#8){%
  \!xM=\!M{#2}\!!xunit
  \!ybM=\!M{#3}\!!yunit
  \!ytM=\!M{#4}\!!yunit
  \!xE=\!M{#6}\!!xunit
  \!ybE=\!M{#7}\!!yunit
  \!ytE=\!M{#8}\!!yunit
  \!getcoeffs\!xS\!ybS\!xM\!ybM\!xE\!ybE\!ybB\!ybC%**Get coefficients B & C for
  \!getcoeffs\!xS\!ytS\!xM\!ytM\!xE\!ytE\!ytB\!ytC%**y=y0 + B(x-X0) +
%%C(x-X0)**2
  \def\!getylimits{\!qgetylimits}%
  \!shade{#1}\ignorespaces}

% ** \!lshade ... (xE,ybE,ytE)
% ** This is like \!qshade, but the top and bottom boundaries are linear,
% ** rather than quadratic.
\def\!lshade#1(#2,#3,#4){%
  \!xE=\!M{#2}\!!xunit
  \!ybE=\!M{#3}\!!yunit
  \!ytE=\!M{#4}\!!yunit
  \!dimenE=\!xE  \advance \!dimenE -\!xS%   ** xE-xS
  \!dimenC=\!ytE \advance \!dimenC -\!ytS%  ** ytE-ytS
  \!divide\!dimenC\!dimenE\!ytB%            ** ytB = (ytE-ytS)/(xE-xS)
  \!dimenC=\!ybE \advance \!dimenC -\!ybS%  ** ybE-ybS
  \!divide\!dimenC\!dimenE\!ybB%            ** ybB = (ybE-ybS)/(xE-xS)
  \def\!getylimits{\!lgetylimits}%
  \!shade{#1}\ignorespaces}

% **  \!getcoeffs{X0}{Y0}{X1}{Y1}{X2}{Y2}{B}{C}
% **  Finds  B  and  C  such that the quadratic  y = Y0 + B(x-X0) + C(x-X0)**2
% **  passes through (X1,Y1) and (X2,Y2):  when X0=0=Y0, the formulas are:
% **                   B = S1 - X1*C,   C = (S2-S1)/X2
% **  with
% **                 S1 = Y1/X1,   S2 = (Y2-Y1)/(X2-X1).
\def\!getcoeffs#1#2#3#4#5#6#7#8{%
  \!dimenC=#4\advance \!dimenC -#2%            ** dimC=Y1-Y0
  \!dimenE=#3\advance \!dimenE -#1%            ** dimE=X1-X0
  \!divide\!dimenC\!dimenE\!dimenF%            ** dimF=S1
  \!dimenC=#6\advance \!dimenC -#4%            ** dimC=Y2-Y1
  \!dimenH=#5\advance \!dimenH -#3%            ** dimH=X2-X1
  \!divide\!dimenC\!dimenH\!dimenG%            ** dimG=S2
  \advance\!dimenG -\!dimenF%                  ** dimG=S2-S1
  \advance \!dimenH \!dimenE%                  ** dimH=X2-X0
  \!divide\!dimenG\!dimenH#8%                  ** C=(S2-S1)/(X2-X0)
  \!removept#8\!t%                             ** C, without "pt"
  #7=-\!t\!dimenE%                             ** -C*(X1-X0)
  \advance #7\!dimenF%                         ** B=S1-C*(X1-X0)
  \ignorespaces}


\def\!shade#1{%
% ** Get LS,RS,BS,TS for this panel
  \!getshrinkages#1<,,,>\!nil% %       ** now effective LS=dimE, RS=dimF,
%                                      **   BS=dimG, TS=dimH
  \advance \!dimenE \!xS%              ** now dimE=xS+LS
  \!lattice\!!xshade\!dshade\!dimenE%  ** set parity=index of left-mst
%%x-lattice
    \!parity\!xpos%                    **   point >= xS+LS, xpos=its location
  \!dimenF=-\!dimenF%                  ** set dimF=xE-RS
    \advance\!dimenF \!xE
%
  \!loop\!not{\ifdim\!xpos>\!dimenF}%  ** loop over x-lattice points <= xE-RS
    \!shadecolumn%
    \advance\!xpos \!dshade%           ** move over to next column
    \advance\!parity 1%                ** increase index of x-point
  \repeat
%
  \!xS=\!xE%                           ** shift ending values to starting
%%values
  \!ybS=\!ybE
  \!ytS=\!ytE
  \ignorespaces}


\def\!vgetshrinkages#1<#2,#3,#4,#5>#6\!nil{%
  \!override\!lshrinkage{#2}\!dimenE
  \!override\!rshrinkage{#3}\!dimenF
  \!override\!bshrinkage{#4}\!dimenG
  \!override\!tshrinkage{#5}\!dimenH
  \ignorespaces}
\def\!hgetshrinkages#1<#2,#3,#4,#5>#6\!nil{%
  \!override\!lshrinkage{#2}\!dimenG
  \!override\!rshrinkage{#3}\!dimenH
  \!override\!bshrinkage{#4}\!dimenE
  \!override\!tshrinkage{#5}\!dimenF
  \ignorespaces}


\def\!shadecolumn{%
  \!dxpos=\!xpos
  \advance\!dxpos -\!xS%            ** dx = x - xS
  \!removept\!dxpos\!dx%            ** ditto, without "pt"
  \!getylimits%                     ** get top and bottom y-values
  \advance\!ytpos -\!dimenH%        ** less TS
  \advance\!ybpos \!dimenG%         ** plus BS
  \!yloc=\!!yshade%                 ** get anchor point for this column
  \ifodd\!parity
     \advance\!yloc \!dshade
  \fi
  \!lattice\!yloc{2\!dshade}\!ybpos%
    \!countA\!ypos%                 ** ypos=smallest y point for this column
  \!dimenA=-\!shadexorigin \advance \!dimenA \!xpos%      ** over
  \loop\!not{\ifdim\!ypos>\!ytpos}% ** loop over ypos <= yt(t)
    \!setshadelocation%             ** vmode: xloc=xpos, yloc=ypos
%                                   ** hmode: xloc=ypos, yloc=xpos
    \!rotateaboutpivot\!xloc\!yloc%
    \!dimenA=-\!shadexorigin \advance \!dimenA \!xloc%    ** over
    \!dimenB=-\!shadeyorigin \advance \!dimenB \!yloc%    ** up
    \kern\!dimenA \raise\!dimenB\copy\!shadesymbol \kern-\!dimenA
    \advance\!ypos 2\!dshade
  \repeat
  \ignorespaces}

\def\!qgetylimits{%
  \!dimenA=\!dx\!ytC
  \advance\!dimenA \!ytB%         ** yt(t)=ytS + dx*(Bt + dx*Ct)
  \!ytpos=\!dx\!dimenA
  \advance\!ytpos \!ytS
  \!dimenA=\!dx\!ybC
  \advance\!dimenA \!ybB%         ** yb(t)=ybS + dx*(Bb + dx*Cb)
  \!ybpos=\!dx\!dimenA
  \advance\!ybpos \!ybS}

\def\!lgetylimits{%
  \!ytpos=\!dx\!ytB%              ** yt(t)=ytS + dx*Bt
  \advance\!ytpos \!ytS
  \!ybpos=\!dx\!ybB%              ** yb(t)=ybS + dx*Bb
  \advance\!ybpos \!ybS}

\def\!vsetshadelocation{%         ** vmode: xloc=xpos, yloc=ypos
  \!xloc=\!xpos
  \!yloc=\!ypos}
\def\!hsetshadelocation{%         ** hmode: xloc=ypos, yloc=xpos
  \!xloc=\!ypos
  \!yloc=\!xpos}


% **************************************
% *** TICKS  (Draws ticks on graphs) ***
% **************************************

% ** User commands
% **   \ticksout
% **   \ticksin
% **   \gridlines
% **   \nogridlines
% **   \loggedticks
% **   \unloggesticks
% ** See Subsection 3.4 of the manual

% ** The following is an option of the \axis command
% **   ticks
% **     [in] [out]
% **     [long] [short] [length <LENGTH>]
% **     [width <WIDTH>]
% **     [andacross] [butnotacross]
% **     [logged] [unlogged]
% **     [unlabeled] [numbered] [withvalues VALUE1 VALUE2 ... VALUEk / ]
% **     [quantity Q] [at LOC1 LOC2 ... LOCk / ] [from LOC1 to LOC2 by
% **       LOC_INCREMENT]
% ** See Subsection 3.2 of the manual for the rules.

% ** The various options of the  tick  field are processed by the
% ** \!nextkeyword  command defined below.
% ** For example, `\!nextkeyword short '  expands to  `\!ticksshort',
% ** while `\!nextkeyword withvalues' expands to `\!tickswithvalues'.

\def\!axisticks {%
  \def\!nextkeyword##1 {%
    \expandafter\ifx\csname !ticks##1\endcsname \relax
      \def\!next{\!fixkeyword{##1}}%
    \else
      \def\!next{\csname !ticks##1\endcsname}%
    \fi
    \!next}%
  \!axissetup
    \def\!axissetup{\relax}%
  \edef\!ticksinoutsign{\!ticksinoutSign}%
  \!ticklength=\longticklength
  \!tickwidth=\linethickness
  \!gridlinestatus
  \!setticktransform
  \!maketick
  \!tickcase=0
  \def\!LTlist{}%
  \!nextkeyword}

\def\ticksout{%
  \def\!ticksinoutSign{+}}
\def\ticksin{%
  \def\!ticksinoutSign{-}}
\ticksout

\def\gridlines{%
  \def\!gridlinestatus{\!gridlinestootrue}}
\def\nogridlines{%
  \def\!gridlinestatus{\!gridlinestoofalse}}
\nogridlines

\def\loggedticks{%
  \def\!setticktransform{\let\!ticktransform=\!logten}}
\def\unloggedticks{%
  \def\!setticktransform{\let\!ticktransform=\!donothing}}
\def\!donothing#1#2{\def#2{#1}}
\unloggedticks

% ** \!ticks/ : terminates read of tick options
\expandafter\def\csname !ticks/\endcsname{%
  \!not {\ifx \!LTlist\empty}
    \!placetickvalues
  \fi
  \def\!tickvalueslist{}%
  \def\!LTlist{}%
  \expandafter\csname !axis/\endcsname}

\def\!maketick{%
  \setbox\!boxA=\hbox{%
    \beginpicture
      \!setdimenmode
      \setcoordinatesystem point at {\!zpt} {\!zpt}
      \linethickness=\!tickwidth
      \ifdim\!ticklength>\!zpt
        \putrule from {\!zpt} {\!zpt} to
          {\!ticksinoutsign\!tickxsign\!ticklength}
          {\!ticksinoutsign\!tickysign\!ticklength}
      \fi
      \if!gridlinestoo
        \putrule from {\!zpt} {\!zpt} to
          {-\!tickxsign\!xaxislength} {-\!tickysign\!yaxislength}
      \fi
    \endpicturesave <\!Xsave,\!Ysave>}%
    \wd\!boxA=\!zpt}

\def\!ticksin{%
  \def\!ticksinoutsign{-}%
  \!maketick
  \!nextkeyword}

\def\!ticksout{%
  \def\!ticksinoutsign{+}%
  \!maketick
  \!nextkeyword}

\def\!tickslength<#1> {%
  \!ticklength=#1\relax
  \!maketick
  \!nextkeyword}

\def\!tickslong{%
  \!tickslength<\longticklength> }

\def\!ticksshort{%
  \!tickslength<\shortticklength> }

\def\!tickswidth<#1> {%
  \!tickwidth=#1\relax
  \!maketick
  \!nextkeyword}

\def\!ticksandacross{%
  \!gridlinestootrue
  \!maketick
  \!nextkeyword}

\def\!ticksbutnotacross{%
  \!gridlinestoofalse
  \!maketick
  \!nextkeyword}

\def\!tickslogged{%
  \let\!ticktransform=\!logten
  \!nextkeyword}

\def\!ticksunlogged{%
  \let\!ticktransform=\!donothing
  \!nextkeyword}

\def\!ticksunlabeled{%
  \!tickcase=0
  \!nextkeyword}

\def\!ticksnumbered{%
  \!tickcase=1
  \!nextkeyword}

\def\!tickswithvalues#1/ {%
  \edef\!tickvalueslist{#1! /}%
  \!tickcase=2
  \!nextkeyword}

\def\!ticksquantity#1 {%
  \ifnum #1>1
    \!updatetickoffset
    \!countA=#1\relax
    \advance \!countA -1
    \!ticklocationincr=\!axisLength
      \divide \!ticklocationincr \!countA
    \!ticklocation=\!axisstart
    \loop \!not{\ifdim \!ticklocation>\!axisend}
      \!placetick\!ticklocation
      \ifcase\!tickcase
          \relax %  Case 0: no labels
        \or
          \relax %  Case 1: numbered -- not available here
        \or
          \expandafter\!gettickvaluefrom\!tickvalueslist
          \edef\!tickfield{{\the\!ticklocation}{\!value}}%
          \expandafter\!listaddon\expandafter{\!tickfield}\!LTlist%
      \fi
      \advance \!ticklocation \!ticklocationincr
    \repeat
  \fi
  \!nextkeyword}

\def\!ticksat#1 {%
  \!updatetickoffset
  \edef\!Loc{#1}%
  \if /\!Loc
    \def\next{\!nextkeyword}%
  \else
    \!ticksincommon
    \def\next{\!ticksat}%
  \fi
  \next}

\def\!ticksfrom#1 to #2 by #3 {%
  \!updatetickoffset
  \edef\!arg{#3}%
  \expandafter\!separate\!arg\!nil
  \!scalefactor=1
  \expandafter\!countfigures\!arg/
  \edef\!arg{#1}%
  \!scaleup\!arg by\!scalefactor to\!countE
  \edef\!arg{#2}%
  \!scaleup\!arg by\!scalefactor to\!countF
  \edef\!arg{#3}%
  \!scaleup\!arg by\!scalefactor to\!countG
  \loop \!not{\ifnum\!countE>\!countF}
    \ifnum\!scalefactor=1
      \edef\!Loc{\the\!countE}%
    \else
      \!scaledown\!countE by\!scalefactor to\!Loc
    \fi
    \!ticksincommon
    \advance \!countE \!countG
  \repeat
  \!nextkeyword}

\def\!updatetickoffset{%
  \!dimenA=\!ticksinoutsign\!ticklength
  \ifdim \!dimenA>\!offset
    \!offset=\!dimenA
  \fi}

\def\!placetick#1{%
  \if!xswitch
    \!xpos=#1\relax
    \!ypos=\!axisylevel
  \else
    \!xpos=\!axisxlevel
    \!ypos=#1\relax
  \fi
  \advance\!xpos \!Xsave
  \advance\!ypos \!Ysave
  \kern\!xpos\raise\!ypos\copy\!boxA\kern-\!xpos
  \ignorespaces}

\def\!gettickvaluefrom#1 #2 /{%
  \edef\!value{#1}%
  \edef\!tickvalueslist{#2 /}%
  \ifx \!tickvalueslist\!endtickvaluelist
    \!tickcase=0
  \fi}
\def\!endtickvaluelist{! /}

\def\!ticksincommon{%
  \!ticktransform\!Loc\!t
  \!ticklocation=\!t\!!unit
  \advance\!ticklocation -\!!origin
  \!placetick\!ticklocation
  \ifcase\!tickcase
    \relax % Case 0: no labels
  \or %      Case 1: numbered
    \ifdim\!ticklocation<-\!!origin
      \edef\!Loc{$\!Loc$}%
    \fi
    \edef\!tickfield{{\the\!ticklocation}{\!Loc}}%
    \expandafter\!listaddon\expandafter{\!tickfield}\!LTlist%
  \or %      Case 2: labeled
    \expandafter\!gettickvaluefrom\!tickvalueslist
    \edef\!tickfield{{\the\!ticklocation}{\!value}}%
    \expandafter\!listaddon\expandafter{\!tickfield}\!LTlist%
  \fi}

\def\!separate#1\!nil{%
  \!ifnextchar{-}{\!!separate}{\!!!separate}#1\!nil}
\def\!!separate-#1\!nil{%
  \def\!sign{-}%
  \!!!!separate#1..\!nil}
\def\!!!separate#1\!nil{%
  \def\!sign{+}%
  \!!!!separate#1..\!nil}
\def\!!!!separate#1.#2.#3\!nil{%
  \def\!arg{#1}%
  \ifx\!arg\!empty
    \!countA=0
  \else
    \!countA=\!arg
  \fi
  \def\!arg{#2}%
  \ifx\!arg\!empty
    \!countB=0
  \else
    \!countB=\!arg
  \fi}

\def\!countfigures#1{%
  \if #1/%
    \def\!next{\ignorespaces}%
  \else
    \multiply\!scalefactor 10
    \def\!next{\!countfigures}%
  \fi
  \!next}

\def\!scaleup#1by#2to#3{%
  \expandafter\!separate#1\!nil
  \multiply\!countA #2\relax
  \advance\!countA \!countB
  \if -\!sign
    \!countA=-\!countA
  \fi
  #3=\!countA
  \ignorespaces}

\def\!scaledown#1by#2to#3{%
  \!countA=#1\relax%                          ** get original #
  \ifnum \!countA<0 %                         ** take abs value,
    \def\!sign{-}%                            **   remember sign
    \!countA=-\!countA
  \else
    \def\!sign{}%
  \fi
  \!countB=\!countA%                          ** copy |#|
  \divide\!countB #2\relax%                   ** integer part (|#|/sf)
  \!countC=\!countB%                          ** get sf * (|#|/sf)
    \multiply\!countC #2\relax
  \advance \!countA -\!countC%                ** ctA is now remainder
  \edef#3{\!sign\the\!countB.}%               ** +- integerpart.
  \!countC=\!countA %                         ** Tack on proper number
  \ifnum\!countC=0 %                          **   of zeros after .
    \!countC=1
  \fi
  \multiply\!countC 10
  \!loop \ifnum #2>\!countC
    \edef#3{#3\!zero}%
    \multiply\!countC 10
  \repeat
  \edef#3{#3\the\!countA}%                    ** Add on rest of remainder
  \ignorespaces}

\def\!placetickvalues{%
  \advance\!offset \tickstovaluesleading
  \if!xswitch
    \setbox\!boxA=\hbox{%
      \def\\##1##2{%
        \!dimenput {##2} [B] (##1,\!axisylevel)}%
      \beginpicture
        \!LTlist
      \endpicturesave <\!Xsave,\!Ysave>}%
    \!dimenA=\!axisylevel
      \advance\!dimenA -\!Ysave
      \advance\!dimenA \!tickysign\!offset
      \if -\!tickysign
        \advance\!dimenA -\ht\!boxA
      \else
        \advance\!dimenA  \dp\!boxA
      \fi
    \advance\!offset \ht\!boxA
      \advance\!offset \dp\!boxA
    \!dimenput {\box\!boxA} [Bl] <\!Xsave,\!Ysave> (\!zpt,\!dimenA)
  \else
    \setbox\!boxA=\hbox{%
      \def\\##1##2{%
        \!dimenput {##2} [r] (\!axisxlevel,##1)}%
      \beginpicture
        \!LTlist
      \endpicturesave <\!Xsave,\!Ysave>}%
    \!dimenA=\!axisxlevel
      \advance\!dimenA -\!Xsave
      \advance\!dimenA \!tickxsign\!offset
      \if -\!tickxsign
        \advance\!dimenA -\wd\!boxA
      \fi
    \advance\!offset \wd\!boxA
    \!dimenput {\box\!boxA} [Bl] <\!Xsave,\!Ysave> (\!dimenA,\!zpt)
  \fi}


\normalgraphs
\catcode`!=12 %  *****  THIS MUST NEVER BE OMITTED

% This is postpictex.tex  Version 1.1  9/10/87

% To use the PiCTeX macros under LaTeX, you first need to \input the
% file prepictex.tex, then the main corpus of PiCTeX macros (pictex.tex),
% and finally this file.  Do not \input the file latexpicobjs.tex.

\catcode`@=11 \catcode`!=11

% Save meanings of PiCTeX keywords that duplicate LaTeX keywords
\let\!pictexendpicture=\endpicture
\let\!pictexframe=\frame
\let\!pictexlinethickness=\linethickness
\let\!pictexmultiput=\multiput
\let\!pictexput=\put

% Redefine the PiCTeX \beginpicture macro
\def\beginpicture{%
  \setbox\!picbox=\hbox\bgroup%
  \let\endpicture=\!pictexendpicture
  \let\frame=\!pictexframe
  \let\linethickness=\!pictexlinethickness
  \let\multiput=\!pictexmultiput
  \let\put=\!pictexput
  \let\input=\@@input   % \@@input is LaTeX's saved version of TeX's primitive
  \!xleft=\maxdimen
  \!xright=-\maxdimen
  \!ybot=\maxdimen
  \!ytop=-\maxdimen}

% Reestablish LaTeX's meaning of \frame. This makes
% PiCTeX's meaning of \frame available only inside a PiCture.
\let\frame=\!latexframe

% Make PiCTeX's meaning of \frame available everywhere in the
% guise of \pictexframe
\let\pictexframe=\!pictexframe

% Now do the same for \linethickness
\let\linethickness=\!latexlinethickness
\let\pictexlinethickness=\!pictexlinethickness

% Reset LaTeX's default meaning of \\
\let\\=\@normalcr
\catcode`@=12 \catcode`!=12


\newtheorem{prop}{Proposition}[section]
\newtheorem{lem}[prop]{Lemma}
\newtheorem{sublem}[prop]{Sublemma}
\newtheorem{cor}[prop]{Corollary}
\newtheorem{them}[prop]{Theorem}
\newtheorem{hypo}[prop]{Hypothesis}
\newtheorem{question}[prop]{Question}
\newtheorem{conjecture}[prop]{Conjecture}
\newtheorem{scholum}[prop]{Scholum}

\theoremstyle{definition}

\newtheorem{defn}[prop]{Definition}
\newtheorem{numcon}[prop]{Construction}
\newtheorem{numrmk}[prop]{Remark}
\newtheorem{numrmks}[prop]{Remarks}
\newtheorem{con}[prop]{}

\theoremstyle{remark}

\newtheorem{warning}{Warning}\renewcommand{\thewarning}{}
\newtheorem{note}{Note}\renewcommand{\thenote}{}
\newtheorem{claim}{Claim}\renewcommand{\theclaim}{}
\newtheorem{example}{Example}\renewcommand{\theexample}{}
\newtheorem{examples}{Examples}\renewcommand{\theexamples}{}
\newtheorem{rmk}{Remark}\renewcommand{\thermk}{}
\newtheorem{rmks}{Remarks}\renewcommand{\thermks}{}

\newenvironment{proof}{\begin{trivlist}\item[]{\sc Proof.}}%
            {\nolinebreak $\Box$ \end{trivlist}}
\newenvironment{unsure}{\footnotesize}{}

\newcommand{\noprint}[1]{}

\newcommand{\rtext}[1]{\text{\normalshape #1}}

%\newcommand{\smalltilde}{\tilde}
\renewcommand{\tilde}{\widetilde}
\newcommand{\dual}{\widehat}

\newcommand{\smooth}{\mbox{\tiny sm}}
\newcommand{\parf}{{\mbox{\tiny parf}}}
\newcommand{\coh}{{\mbox{\tiny coh}}}
\newcommand{\cart}{{\mbox{\tiny cart}}}
\newcommand{\sm}{\mbox{\tiny sm}}
\newcommand{\an}{\mbox{\tiny an}}
\newcommand{\etale}{\mbox{\tiny \'{e}t}}
\newcommand{\et}{\mbox{\tiny \'{e}t}}
\newcommand{\Zariski}{\mbox{\tiny Zar}}
\renewcommand{\flat}{\mbox{\tiny fl}}
\newcommand{\fl}{\mbox{\tiny fl}}
\newcommand{\topo}{\mbox{\tiny top}}
\newcommand{\alg}{\mbox{\tiny alg}}
\newcommand{\tors}{\mbox{\tiny tors}}
\newcommand{\semstab}{\mbox{\tiny ss}}
\newcommand{\stab}{\mbox{\tiny stab}}
\newcommand{\virt}{\mbox{\tiny virt}}
\newcommand{\th}{{\operatorname{\rm th}}}
\newcommand{\exptilde}{\operatorname{\widetilde{\phantom{N}}}}
\newcommand{\upst}{^{\ast}}
\newcommand{\op}{{\text{\normalshape op}}}
\newcommand{\upsh}{^{!}}
\newcommand{\lst}{_{\ast}}
\newcommand{\lsh}{_{!}}
\newcommand{\com}{^{\scriptscriptstyle\bullet}}
\newcommand{\lcom}{_{\scriptscriptstyle\bullet}}
\newcommand{\upcom}{^{\scriptscriptstyle\bullet}}
\newcommand{\argument}{{{\,\cdot\,}}}

\newcommand{\XX}{{\frak X}}
\renewcommand{\AA}{{\frak A}}
\newcommand{\BB}{{\frak B}}
\newcommand{\SS}{{\frak S}}
\newcommand{\TT}{{\frak T}}
\newcommand{\YY}{{\frak Y}}
\newcommand{\UU}{{\frak U}}
\newcommand{\PP}{{\frak P}}
\newcommand{\ZZ}{{\frak Z}}
\newcommand{\JJ}{{\frak J}}
\newcommand{\EE}{{\frak E}}
\newcommand{\GG}{{\frak G}}
\newcommand{\NN}{{\frak N}}
\newcommand{\MM}{{\frak M}}
\newcommand{\HH}{{\frak H}}
\newcommand{\VV}{{\frak V}}
\newcommand{\WW}{{\frak W}}
\newcommand{\Gg}{{\frak g}}
\newcommand{\Aa}{{\frak a}}
\newcommand{\Ee}{{\frak e}}
\newcommand{\Pp}{{\frak p}}
\newcommand{\Mm}{{\frak m}}
\newcommand{\Nn}{{\frak n}}
\newcommand{\Qq}{{\frak q}}
\newcommand{\Tt}{{\frak t}}
\newcommand{\Vv}{{\frak v}}
\newcommand{\Ww}{{\frak w}}
\newcommand{\Uu}{{\frak u}}
\newcommand{\Ll}{{\frak l}}
\newcommand{\Cc}{{\frak c}}
\newcommand{\Dd}{{\frak d}}
\newcommand{\CC}{{\frak C}}
\newcommand{\zz}{{\Bbb Z}}
\newcommand{\hh}{{\Bbb H}}
\newcommand{\vv}{{\Bbb V}}
\newcommand{\nz}{{{\Bbb N}_0}}
\newcommand{\aaa}{{\Bbb A}}
\newcommand{\ttt}{{\Bbb T}}
\newcommand{\nn}{{\Bbb N}}
\newcommand{\dd}{{\Bbb D}}
\newcommand{\mm}{{\Bbb M}}
\renewcommand{\ll}{{\Bbb L}}
\newcommand{\qq}{{\Bbb Q}}
\newcommand{\pp}{{\Bbb P}}
\newcommand{\cc}{{\Bbb C}}
\newcommand{\rr}{{\Bbb R}}
\newcommand{\ffq}{{\Bbb F}_q}
\newcommand{\ffp}{{\Bbb F}_p}
\newcommand{\ffqn}{{\Bbb F}_{q^n}}
\newcommand{\ffqb}{{\Bbb {\ol{F}}}_q}
\newcommand{\Gm}{{{\Bbb G}_{\mbox{\tiny\rm m}}}}
\newcommand{\Ga}{{{\Bbb G}_{\mbox{\tiny\rm a}}}}
\newcommand{\zzl}{{{\Bbb Z}_\ell}}
\newcommand{\qql}{{{\Bbb Q}_\ell}}
\newcommand{\qqp}{{{\Bbb Q}_p}}
\newcommand{\zzp}{{{\Bbb Z}_p}}
\newcommand{\qqlb}{{{\Bbb \ol{Q}}_\ell}}
\newcommand{\tT}{{\cal T}}
\newcommand{\ox}{{\cal O}_X}
\renewcommand{\O}{{\cal O}}
\newcommand{\dD}{{\cal D}}
\newcommand{\cC}{{\cal C}}
\newcommand{\eE}{{\cal E}}
\newcommand{\iI}{{\cal I}}
\newcommand{\gG}{{\cal G}}
\newcommand{\fF}{{\cal F}}
\newcommand{\lL}{{\cal L}}
\newcommand{\mM}{{\cal M}}
\newcommand{\hH}{{\cal H}}
\newcommand{\sS}{{\cal S}}
\newcommand{\uU}{{\cal U}}
\newcommand{\oO}{{\cal O}}
\renewcommand{\pt}{{\rm pt}}

%\renewcommand{\theenumi}{\roman{enumi}}

\newcommand{\del}{\partial}
\newcommand{\resto}{{ \mid }}
\newcommand{\st}{\mathrel{\mid}}
\newcommand{\on}{\mathbin{\mid}}
\newcommand{\rk}{\operatorname{\rm rk}}
\newcommand{\ev}{\operatorname{\rm ev}}
\newcommand{\td}{\operatorname{\rm td}}
\newcommand{\wt}{\operatorname{\rm wt}}
\newcommand{\pr}{\operatorname{\rm pr}\nolimits}
\newcommand{\Par}{\operatorname{\rm Par}\nolimits}
\newcommand{\Gr}{\operatorname{\rm Gr}}
\newcommand{\gr}{\operatorname{\rm gr}\nolimits}
\newcommand{\cl}{\operatorname{\rm cl}\nolimits}
\newcommand{\grs}{\operatorname{\rm grs}\nolimits}
\newcommand{\Points}{\operatorname{\rm Points}\nolimits}
\newcommand{\sets}{\operatorname{\rm sets}\nolimits}
\newcommand{\Mod}{\operatorname{\rm Mod}\nolimits}
\newcommand{\Mor}{\operatorname{\rm Mor}\nolimits}
\newcommand{\fg}{\mbox{\scriptsize fg}}
\newcommand{\ch}{\operatorname{\sf X}}
\newcommand{\Quot}{\operatorname{\rm Quot}}
\newcommand{\tr}{\operatorname{\rm tr}\nolimits}
\newcommand{\lc}{\operatorname{\rm lc}}
\newcommand{\Cl}{\operatorname{\rm Cl}}
\renewcommand{\Re}{\operatorname{\rm Re}}
\renewcommand{\Im}{\operatorname{\rm Im}}
\renewcommand{\top}{\operatorname{\rm top}}
\newcommand{\sign}{\operatorname{\rm sign}}
\newcommand{\curl}{\operatorname{\rm curl}}
\newcommand{\grad}{\operatorname{\rm grad}}
\newcommand{\ord}{\operatorname{\rm ord}\nolimits}
\newcommand{\res}{\operatorname{\rm res}\nolimits}
\renewcommand{\div}{\operatorname{\rm div}}
\newcommand{\const}{\operatorname{\rm const}}
\newcommand{\Trans}{\operatorname{\rm Trans}\nolimits}
\newcommand{\Ad}{\operatorname{\rm Ad}\nolimits}
\newcommand{\rank}{\operatorname{\rm rank}\nolimits}
\newcommand{\num}{\operatorname{\rm num}\nolimits}
\newcommand{\red}{\operatorname{\rm red}\nolimits}
\newcommand{\supp}{\operatorname{\rm supp}}
\newcommand{\spann}{\operatorname{\rm span}}
\newcommand{\conv}{\operatorname{\rm conv}}
\newcommand{\Dyn}{\operatorname{\rm Dyn}\nolimits}
\newcommand{\Picc}{\operatorname{\rm Pic}\nolimits}
\newcommand{\Jac}{\operatorname{\rm Jac}\nolimits}
\newcommand{\multdeg}{\operatorname{\rm mult-deg}\nolimits}
\newcommand{\im}{\operatorname{\rm im}}
\newcommand{\ob}{\operatorname{ob}}
\newcommand{\fle}{\operatorname{\rm fl}}
\newcommand{\ab}{\operatorname{\rm ab}}
\newcommand{\source}{\operatorname{\rm source}}
\newcommand{\target}{\operatorname{\rm target}}
\newcommand{\cok}{\operatorname{\rm cok}}
\newcommand{\chr}{\operatorname{\rm char}\nolimits}
\newcommand{\codim}{\operatorname{\rm codim}}
\newcommand{\spec}{\operatorname{\rm Spec}\nolimits}
\newcommand{\Sym}{\operatorname{\rm Sym}\nolimits}
\newcommand{\proj}{\operatorname{\rm Proj}\nolimits}
\newcommand{\id}{\operatorname{\rm id}}
\newcommand{\Hom}{\operatorname{\rm Hom}\nolimits}
\newcommand{\Ext}{\operatorname{\rm Ext}\nolimits}
\newcommand{\Tor}{\operatorname{\rm Tor}\nolimits}
\newcommand{\wierdtensor}{\operatorname{\otimes}\limits}
\newcommand{\ltensor}{\operatorname{\displaystyle\stackrel{L}{\otimes}}}
\newcommand{\sheafhom}{\operatorname{\rm {\mit{ \hH\! om}}}\nolimits}
\newcommand{\sheafepi}{\operatorname{\rm {\mit{ \eE\! pi}}}\nolimits}
\newcommand{\sheafext}{\operatorname{\rm {\mit{ \eE\! xt}}}\nolimits}
\newcommand{\sheaftor}{\operatorname{\rm {\mit{ \tT\! or}}}\nolimits}
\newcommand{\Homu}{\operatorname{\underline{\rm Hom}}\nolimits}
\newcommand{\Isom}{\operatorname{\rm Isom}\nolimits}
\newcommand{\Isomu}{\operatorname{\underline{\rm Isom}}\nolimits}
\newcommand{\Isomext}{\operatorname{\rm Isomext}\nolimits}
\newcommand{\Isomextu}{\operatorname{\underline{\rm Isomext}}\nolimits}
\newcommand{\End}{\operatorname{\rm End}\nolimits}
\newcommand{\Endu}{\operatorname{\underline{\rm End}}\nolimits}
\newcommand{\Aut}{\operatorname{\rm Aut}\nolimits}
\newcommand{\Autu}{\operatorname{\underline{\rm Aut}}\nolimits}
\newcommand{\AAut}{\operatorname{\frak Aut}\nolimits}
\newcommand{\Gal}{\operatorname{\rm Gal}\nolimits}
\newcommand{\Pic}{\operatorname{\rm Pic}\nolimits}
\newcommand{\Stab}{\operatorname{\rm Stab}\nolimits}
\newcommand{\Grass}{\operatorname{\rm Grass}\nolimits}
%\newcommand{\injlim}{\operatorname{\lim\limits_{\longrightarrow}}\limits}
%% FOLLOWING LINE CANNOT BE BROKEN BEFORE 80 CHAR
\renewcommand{\projlim}{\operatorname{{\lim\limits_{
\textstyle\longleftarrow}}}\limits}
\newcommand{\comp}{\mathbin{{\scriptstyle\circ}}}
\newcommand{\ol}{\overline}
\newcommand{\ul}{\underline}
\newcommand{\closure}{^{\textstyle -}}

\newcommand{\ldiag}[1]%
       {\makebox[0cm]{${\scriptstyle#1}\downarrow\phantom{\scriptstyle#1}$}}
\newcommand{\ldiagup}[1]%
       {\makebox[0cm]{${\scriptstyle#1}\uparrow\phantom{\scriptstyle#1}$}}
\newcommand{\rdiag}[1]%
       {\makebox[0cm]{$\phantom{\scriptstyle#1}\downarrow{\scriptstyle#1}$}}
\newcommand{\sediagr}[1]%
       {\makebox[0cm]{$\phantom{\scriptstyle#1}\searrow{\scriptstyle#1}$}}
\newcommand{\nediagr}[1]%
       {\makebox[0cm]{$\phantom{\scriptstyle#1}\nearrow{\scriptstyle#1}$}}
\newcommand{\rdiagup}[1]%
       {\makebox[0cm]{$\phantom{\scriptstyle#1}\uparrow{\scriptstyle#1}$}}
\newcommand{\swdiag}[1]%
       {\makebox[0cm]{$\phantom{\scriptstyle#1}\swarrow{\scriptstyle#1}$}}
\newcommand{\sediag}[1]%
       {\makebox[0cm]{${\scriptstyle#1}\searrow\phantom{\scriptstyle#1}$}}
\newcommand{\nediag}[1]%
       {\makebox[0cm]{${\scriptstyle#1}\nearrow\phantom{\scriptstyle#1}$}}

\newcommand{\longiso}{\stackrel{\textstyle\sim}{\longrightarrow}}
\newcommand{\iso}{\stackrel{\sim}{\rightarrow}}
\newcommand{\longisoback}{\stackrel{\textstyle\sim}{\longleftarrow}}

\newcommand{\scalar}[2]{\langle #1,#2\rangle}

\newcommand{\doublearrowstack}[2]%
{{{{\scriptstyle#1}\atop{\textstyle \longrightarrow}}\atop{{\textstyle
\longrightarrow}\atop{ \scriptstyle#2}}}}
\newcommand{\rightleftarrowstack}[2]%
{{{{\scriptstyle#1}\atop{\textstyle \longrightarrow}}\atop{{\textstyle
\longleftarrow}\atop{ \scriptstyle#2}}}}
\newcommand{\leftrightarrowstack}[2]%
{{{{\scriptstyle#1}\atop{\textstyle \longleftarrow}}\atop{{\textstyle
\longrightarrow}\atop{ \scriptstyle#2}}}}

\newcommand{\distri}[3]{%
        \begin{array}{ccccc}
           &          & #3              &          &    \\
           & \swarrow &                 & \nwarrow &    \\
        #1 &          & \longrightarrow &          & #2
        \end{array}}

\newcommand{\comment}[1]{{\marginpar{\footnotesize #1}}}

\newcommand{\comdia}[9]{%
\begin{array}{ccc}
#1 & \stackrel{#2}{\longrightarrow} & #3 \\
\ldiag{#4} & #5 & \rdiag{#6} \\
#7 & \stackrel{#8}{\longrightarrow} & #9
\end{array}}
\newcommand{\comdiaback}[9]{%
\begin{array}{ccc}
#1 & \stackrel{#2}{\longleftarrow} & #3 \\
\ldiag{#4} & #5 & \rdiag{#6} \\
#7 & \stackrel{#8}{\longleftarrow} & #9
\end{array}}

\newcommand{\comdiaup}[9]{%
\begin{array}{ccc}
#1 & \stackrel{#2}{\longrightarrow} & #3 \\
\ldiagup{#4} & #5 & \rdiagup{#6} \\
#7 & \stackrel{#8}{\longrightarrow} & #9
\end{array}}

\newcommand{\comdiaupback}[9]{%
\begin{array}{ccc}
#1 & \stackrel{#2}{\longleftarrow} & #3 \\
\ldiagup{#4} & #5 & \rdiagup{#6} \\
#7 & \stackrel{#8}{\longleftarrow} & #9
\end{array}}

\newcommand{\comtri}[6]{%
\begin{array}{ccc}
#1 & \stackrel{#2}{\longrightarrow} & #3         \\
   & \sediag{#4}                    & \rdiag{#5} \\
   &                                & #6
\end{array}}

\newcommand{\inversecomtri}[6]{%
\begin{array}{ccc}
#1           &                                &            \\
\ldiag{#2}   & \sediagr{#3}                   &            \\
#4           & \stackrel{#5}{\longrightarrow} & #6
\end{array}}

\newcommand{\comtriup}[6]{%
\begin{array}{ccc}
#1 & \stackrel{#2}{\longrightarrow} & #3         \\
   & \sediag{#4}                    & \rdiagup{#5} \\
   &                                & #6
\end{array}}

\newcommand{\overtoparrow}%
{\makebox[0cm]{\beginpicture
\setcoordinatesystem units <.8cm,.4cm> point at 0 0
\setplotarea x from -3 to 3, y from 0 to 1
\setquadratic
\plot -3 0 0 1 3 0 /
\put{\vector(3,-1){0}}[Bl] at 3 0
\endpicture}}

\newcommand{\underbottomarrow}%
{\makebox[0cm]{\beginpicture
\setcoordinatesystem units <.8cm,.4cm> point at 0 0
\setplotarea x from -3 to 3, y from 0 to 1
\setquadratic
\plot -3 1 0 0 3 1 /
\put{\vector(3,1){0}}[Bl] at 3 1
\endpicture}}

\renewcommand{\t}{_{\tau}}
\newcommand{\s}{_{\sigma}}

\setcounter{secnumdepth}{1}
\setcounter{tocdepth}{2}

\begin{document}
\maketitle
\begin{abstract}
We construct the motivic tree-level system of Gromov-Witten invariants
for convex varieties.
\end{abstract}
%\tableofcontents

\setcounter{section}{-1}
\section{Introduction}

Let $V$ be a projective algebraic manifold. In \cite{KM}, Sec.\ 2,
Gromov-Witten invariants of $V$ were described
axiomatically as a collection of linear maps
\[I^{V}_{g,n,\beta}:\ H\upst(V)^{\otimes n}\longrightarrow
H\upst(\ol{M}_{g,n},\qq),\quad \beta\in H_2(V,\zz)\]
satisfying certain axioms, and a program to construct them by algebro-geometric
(as opposed to symplectic) techniques
was suggested. The program is based upon Kontsevich's notion of a stable map
$(C,x_1,\dots ,x_n,f)$, $f:C\rightarrow V$.
This data consists of an algebraic curve $C$ with $n$ labeled points on it and
a map $f$ such that if an irreducible
component of $C$ is contracted by $f$ to a point, then this component together
with its special points is
Deligne-Mumford stable.  For more details, see \cite{K} and below.

\smallskip

The construction consists of three major steps.

\medskip

A. Construct an orbispace (or rather a stack) of stable maps
$\ol{M}_{g,n}(V,\beta)$ such that $g=\text{genus of }C$,
$f\lst{([C])}=\beta$, and its two morphisms to $V^n$ and $\ol{M}_{g,n}$. On the
level of points, these morphisms are
given respectively by
\begin{eqnarray*}
p:(C,x_1,\dots ,x_n,f) & \longmapsto & (f(x_1),\dots ,f(x_n)), \\
q:(C,x_1,\dots ,x_n,f) & \longmapsto & [(C,x_1,\dots ,x_n)]^{\stab},
\end{eqnarray*}
where the last expression means the  stabilization of $(C,x_1,\dots ,x_n)$.

\medskip

B. Construct a ``virtual fundamental class'' $[\ol{M}_{g,n}(V,\beta)]_{\virt}$,
or ``orientation'' (see
Definition~\ref{domb} below) and use it to define a correspondence in the Chow
ring $C^V_{g,n,\beta}\in A(V^n\times
\ol{M}_{g,n})$.

\smallskip

This step suggested in \cite{K} is quite subtle and has not been spelled out in
full detail. It can be bypassed for
$g=0$ and $V=G/P$ (generalized flag spaces) where the virtual class coincides
with the usual one (see \cite{KM}).

\smallskip

In general, it involves a definition of a new $\zz$-graded supercommutative
structure sheaf on $\ol{M}_{g,n}(V,\beta)$.
The virtual class is obtained as a product of the class of this sheaf and the
inverse Todd class of an appropriate
tangent complex. Geometrically, it serves as a general position argument
furnishing the Dimension Axiom of \cite{KM} and
replacing the deformation of the complex structure used in the symplectic
context.

\medskip

C. Use $C^V_{g,n,\beta}$ in order to construct the induced maps
$I^V_{g,n,\beta}$ on any cohomology satisfying some
version of the standard properties making it functorial on the category of
correspondences.

\medskip

In this approach, the main features of $I^V_{g,n,\beta}$ axiomatized in
\cite{KM} reflect functorial properties of
$\ol{M}_{g,n}(V,\beta)$ and the cotangent complex with respect to degenerations
of stable maps. In particular, the key
``Splitting Axiom'' (or Associativity Equations for $g=0$) expresses the
compatibility between the divisors at infinity
of $\ol{M}_{g,n}(V,\beta)$ and $\ol{M}_{g,n}.$

\smallskip

A neat way to organize this information is to introduce the category of marked
stable modular graphs indexing
degeneration types of stable maps and to treat various modular stacks
$\ol{M}_{g,n}(V,\beta)$ as values of this modular
functor on the simplest one-vertex graphs. Then the check of the axioms in
\cite{KM} essentially boils down to a
calculation of this functor on a family of generating morphisms and objects in
the graph category.

\smallskip

The degeneration type of $(C,x_1,\dots ,x_n, f)$ is described by the graph
whose vertices are the irreducible components
of $C$, edges are singular points of $C$, and tails (``one-vertex edges'') are
$x_1,\dots ,x_n.$ In addition, each
vertex is marked by the homology class in $V$ which is the $f$-image of the
fundamental class of the respective
component of $C$ and by the genus of the normalization of this component. The
description of morphisms is somewhat more
delicate, cf.\ Sec.\ 1 below.

\smallskip

This philosophy is an extension of the operadic picture which already gained
considerable importance from various
viewpoints.  In turn, it leads to a new notion of a $\Gamma$-operad as a
monoidal functor on an appropriate category
$\Gamma$ of graphs, and an algebra over an operad as a morphism of such
functors.  This approach will be developed
elsewhere (see \cite{KapMan}). It clarifies the origin of the proliferation of
the types of operads considered recently
(May's, Markl's, modular, cyclic, ...)

\smallskip

In Part~I of the present paper we treat in this way Step~A, stressing the
functoriality not only with respect to the
degeneration types with fixed $V$ but also with respect to $V$, expressed by
the change of the marking semigroup of
abstract non-negative homology classes. We hope also that our approach will
help to introduce quantum cohomology with
coefficients and to understand better the K\"unneth formula for quantum
cohomology from \cite{KM2}.

\smallskip

Part~II is devoted to Steps~B and~C for $g=0$ and convex manifolds $V$. The
formalism of orientation classes is
introduced axiomatically, but we did not attempt to justify the relevant claims
of \cite{K} in general.

\smallskip

A word of warning and apology is due. The reader will meet several different
categories of marked graphs in this paper
of which the most important are $\GG_s$ (cf.\ Definition~\ref{dmmsg}),
$\tilde{\GG}_s(A)$ (cf.\ Definition~\ref{ecisg}
and the preceding discussion) and $\tilde{\GG}_s(V)_{\cart}$ (cf.\
Definition~\ref{ceicv}). They differ mainly by their
classes of morphisms. Specifically, certain elementary arrows which are
combinatorially ``the same'', run in opposite
directions in different categories, which affects the whole structure of the
morphism semigroups. The reason is that
functorial properties of moduli stacks of maps considered {\em by themselves }
are different form the functorial
properties of their virtual fundamental classes treated {\em as
correspondences}. Since graphs are used mainly as a
bookkeeping device, their categorical properties must reflect this distinction.

\medskip

\subsection{Acknowledgements}

The authors are grateful to the Max-Planck-Institut f\"ur Mathematik in Bonn
where this work was done for support and
stimulating atmosphere.

\vfill\eject
\part{Stacks of Stable Maps}

\section{Graphs} \label{graphs}

\begin{defn} \label{dog}
A {\em graph }$\tau$ is a quadruple $(F\t,V\t,j\t,\del\t)$, where $F\t$ and
$V\t$ are finite sets,
$\del\t:F\t\rightarrow V\t$ is a map and $j\t:F\t\rightarrow F\t$ an
involution. We call $F\t$ the set of {\em flags},
$V\t$ the set of {\em vertices}, $S\t=\{f\in F\t\st j\t f=f\}$ the set of {\em
tails } and $E\t=\{\{f_1,f_2\}\subset
F\t\st\text{$f_2=j\t f_1$}\}$ the set of edges of $\tau$. For $v\in V\t$ let
$F\t(v)=\del^{-1}\t(v)$ and $|v|=\#
F\t(v)$, the {\em valence }of $v$.
\end{defn}

\begin{defn} \label{geomr}
Let $\tau$ be a graph. We define the {\em geometric realization $|\tau|$ of
$\tau$} as follows. We start with the
disjoint union of closed intervals and singletons
\[\coprod_{f\in F\t}[0,{\textstyle{1\over2}}]\amalg\coprod_{v\in V\t}\{|v|\}.\]
We denote the real number $x\in [0,{1\over2}]$ in the component indexed by
$f\in F\t$ by $x_f$.  Then for every $v\in
V\t$ we identify all elements of $\{0_f\st f\in F\t(v)\}$ with $|v|$ and for
every edge $\{f_1,f_2\}$ of $\tau$, we
identify ${1\over2}_{f_1}$ and ${{1\over2}}_{f_2}$. Finally, we remove for
every tail $f\in S\t$ the point
${1\over2}_f$. We consider $|\tau|$ as a topological space with base points
given by $\{|v|\st v\in V\t\}$, the {\em
vertices } of $|\tau|$. It should always be clear from the context whether
$|v|$ denotes the geometric realization of a
vertex or it's valence.
\end{defn}

\begin{defn} \label{doc}
Let $\tau$ and $\sigma$ be graphs. A {\em contraction
}$\phi:\tau\rightarrow\sigma$ is a pair of maps
$\phi^F:F\s\rightarrow F\t$ and $\phi_V:V\t\rightarrow V\s$ such that the
following conditions are satisfied.
\begin{enumerate}
\item $\phi^F$ is injective and $\phi_V$ is surjective.
\item The diagram
\[\begin{array}{ccc}
F\t & \stackrel{\del\t}{\longrightarrow} & V\t \\
\ldiagup{\phi^F} & & \rdiag{\phi_V} \\
F\s & \stackrel{\del\s}{\longrightarrow} & V\s
\end{array}\] commutes.
\item $\phi^F\comp j\s =j\t\comp \phi^F$, so that $\phi$ induces injections
$\phi^S:S\s\rightarrow S\t$ and
$\phi^E:E\s\rightarrow E\t$ on tails and edges.
\item $\phi^S$ is a bijection, so $F\t-\phi^F(F\s)$ consists entirely of edges,
the edges being contracted.
\item Call two vertices $v,w\in V\t$ {\em equivalent}, if there exists an $f\in
F\t-\phi^F(F\s)$ such that $f\in F\t(v)$
and $j\t f\in F\t(w)$. Then pass to the associated equivalence relation on
$V\t$. The map $\phi_V:V\t\rightarrow V\s$
induces a bijection $V\t/\mathop{\sim}\rightarrow V\s$.
\end{enumerate}
For a vertex $v\in V\s$ the graph whose set of flags is
\[\{f\in F\t\st \rtext{$f\not\in\phi^F(F\s)$ and $\phi_V(\del\t f)=v$}\},\]
whose set of vertices is $\phi_V^{-1}(v)$ and whose $j$ and $\del$ are obtained
from $j\t$ and $\del\t$ by restriction,
is called the {\em graph being contracted onto $v$}. If the graphs being
contracted have together exactly one edge, we
call $\phi$ an {\em elementary contraction}.
\end{defn}

\begin{numrmks} \label{rmc}
\begin{enumerate}
\item It is clear how to compose contractions, and that the composition of
contractions is a contraction.
\item If $\phi:\tau\rightarrow\sigma$ and $\phi':\tau\rightarrow\sigma'$ are
contractions with the same set of edges
being contracted, then there exists a unique isomorphism
$\psi:\sigma\rightarrow\sigma'$ such that
$\psi\comp\phi=\phi'$.
\item Every contraction is a composition of elementary contractions.
\item \label{rmcone} To carry out a construction for contractions of graphs,
which is compatible with composition of
contractions, it suffices to perform this construction for elementary
contractions and check that the construction is
independent of the order in which it is realized for two elementary
contractions.
\end{enumerate}
\end{numrmks}

\begin{defn} \label{domg}
A {\em modular graph } is a graph $\tau$ endowed with a map
$g\t:V\t\rightarrow\zz_{\geq0};v\mapsto g(v)$. The number
$g(v)$ is called the {\em genus }of the vertex $v$.
\end{defn}

We say that a semigroup $A$ has {\em indecomposable zero}, if $a+b=0$ implies
$a=0$ and $b=0$, for any two elements
$a,b\in A$.

\begin{defn} \label{doasmg}
Let $\tau$ be a modular graph and $A$ a semigroup with indecomposable zero. An
{\em $A$-structure } on $\tau$ is a map
$\alpha:V_{\tau}\rightarrow A$. The element $\alpha(v)$ is called the {\em
class } of the vertex $v$. The pair
$(\tau,\alpha)$ is called a {\em modular graph with $A$-structure }(or {\em
$A$-graph}, by abuse of language).

A {\em marked graph } is a pair $(A,\tau)$, where $A$ is a semigroup with
indecomposable zero and $\tau$ an $A$-graph.
\end{defn}

\begin{defn} \label{commor}
Let $(\sigma,\alpha)$ and $(\tau,\beta)$ be $A$-graphs. A {\em combinatorial
morphism
$a:(\sigma,\alpha)\rightarrow(\tau,\beta)$} is a pair of maps
$a_F:F\s\rightarrow F\t$ and $a_V:V\s\rightarrow V\t$,
satisfying the following conditions.
\begin{enumerate}
\item \label{commor1} The diagram
\[\comdia{F\s}{\del\s}{V\s}{a_F}{}{a_V}{F\t}{\del\t}{V\t}\]
commutes. In particular, for every $v\in V\s$, letting $w=a_V(v)$, we get an
induced map $a_{V,v}:F\s(v)\rightarrow
F\t(w)$.
\item With the notation of (\ref{commor1}), for every $v\in V\s$ the map
$a_{V,v}:F\s(v)\rightarrow F\t(w)$ is
injective.
\item \label{commor3} Let $f\in F\s$ and $\ol{f}=j\s(f)$. If $f\not=\ol{f}$,
there exists an $n\geq1$ and $2n$ (not
necessarily distinct) flags $f_1,\ldots,f_n,\ol{f}_1,\ldots,\ol{f}_n\in F\t$
such that
\begin{enumerate}
\item $f_1=a_F(f)$ and $\ol{f}_n=a_F(\ol{f})$,
\item $j\t(f_i)=\ol{f}_i$, for all $i=1,\ldots,n$,
\item $\del\t(\ol{f}_i)=\del\t(f_{i+1})$ for all $i=1,\ldots,n-1$,
\item for all $i=1,\ldots,n-1$ we have
\[\ol{f}_i\not= f_{i+1}\Longrightarrow\text{$g({v_i})=0$ and $\beta(v_i)=0$},\]
where $v_i=\del(\ol{f}_i)=\del(f_{i+1})$,
\end{enumerate}
\item for every $v\in V\s$ we have $\alpha(v)=\beta(a_V(v))$,
\item for every $v\in V\s$ we have $g({v})=g({a_V(v)})$.
\end{enumerate}

A {\em combinatorial morphism of marked graphs
}$(B,\sigma,\beta)\rightarrow(A,\tau,\alpha)$ is a pair $(\xi,a)$, where
$\xi:A\rightarrow B$ is a homomorphism of semigroups and
$a:(\sigma,\beta)\rightarrow(\tau,\xi\comp\alpha)$ is a
combinatorial morphism of $B$-graphs.

Usually, we will suppress the subscripts of $a$.
\end{defn}

\begin{rmks}
\begin{enumerate}
\item The composition of two combinatorial morphisms is again a combinatorial
morphism.
\item We say that a combinatorial morphism $a:\sigma\rightarrow\tau$ is {\em
complete}, if for every $v\in V\s$ the map
$a_{V,v}:F\s(v)\rightarrow F\t(a(v))$ is bijective. Examples of complete
combinatorial morphism are
\begin{enumerate}
\item the inclusion of a connected component,
\item the morphism $\sigma\rightarrow\tau$, where $\sigma$ is obtained from
$\tau$ by {\em cutting an edge}, i.e.\
changing $j$ in such a way as to turn a two element orbit into two one element
orbits.
\end{enumerate}
\item Let $\tau$ be an $A$-graph and $f\in S\t$ a tail of $\tau$. Let
$F\s=F\t-\{f\}$, $V\s=V\t$ and define $\del\s$ and
$j\s$ by restricting $\del\t$ and $j\t$. Then $\sigma$ is naturally an
$A$-graph called {\em obtained from $\tau$ by
forgetting the tail $f$}. There is a canonical combinatorial morphism
$\sigma\rightarrow\tau$.
\item Every combinatorial morphism $a:\sigma\rightarrow\tau$ is a composition
$a=b\comp c$, where $b$ is complete and
$c$ is a finite composition of morphisms forgetting tails. If $\sigma$ and
$\tau$ are stable (Definition~\ref{dosag}),
all intermediate graphs in such a factorization are stable.
\item Condition~(\ref{commor3}) of Definition~\ref{commor} can be rephrased in
a more geometric way---see the remark
after Proposition~\ref{cmmgpe}.
\end{enumerate}
\end{rmks}

\begin{defn}
A {\em contraction }$\phi:(\tau,\alpha)\rightarrow(\sigma,\beta)$ of $A$-graphs
is a contraction of graphs
$\phi:\tau\rightarrow\sigma$ such that for every $v\in V\s$ we have
\begin{enumerate}
\item \[g(v)=\sum_{w\in\phi_V^{-1}(v)}g(w)+\dim H^1(|\tau_v|),\]
where $\tau_v$ is the graph being contracted onto $v$,
\item \[\beta(v)=\sum_{w\in\phi_V^{-1}(v)}\alpha(w).\]
\end{enumerate}
\end{defn}

\begin{defn} \label{dosag}
A vertex $v$ of a modular graph with $A$-structure $(\tau,\alpha)$ is called
{\em stable}, if $\alpha(v)=0$ implies
$2g(v)+|v|\geq3$. Otherwise, $v$ is called {\em unstable}. The $A$-graph $\tau$
is called {\em stable}, if all its
vertices are stable.
\end{defn}

We now come to an important construction which we shall call {\em stable
pullback}. Consider the following setup.  We
suppose given a homomorphism of semigroups $\xi:A\rightarrow B$, a contraction
of $A$-graphs
$\phi:\sigma\rightarrow\tau$ and a combinatorial morphism
$a:(B,\rho)\rightarrow(A,\tau)$ of marked graphs.  Moreover,
we assume that $\rho$ is a {\em stable }$B$-graph. We shall construct a stable
$B$-graph $\pi$, together with a
contraction of $B$-graphs $\psi:\pi\rightarrow\rho$ and a combinatorial
morphism of marked graphs
$b:\pi\rightarrow\sigma$. This $B$-graph $\pi$ will be called the {\em stable
pullback }of $\rho$ under $\phi$.
\[\begin{array}{ccccc}
B & & \pi & \stackrel{\psi}{\longrightarrow} & \rho \\
\ldiagup{\xi} & \phantom{\longrightarrow} & \ldiag{b} &  & \rdiag{a} \\
A & & \sigma & \stackrel{\phi}{\longrightarrow} & \tau
\end{array}\]

According with Remark~\ref{rmc}(\ref{rmcone}), we shall assume that $\phi$ is
elementary and contracts the edge
$\{f,\ol{f}\}$ of $\sigma$. Let $v_1=\del\s(f)$, $v_2=\del\s(\ol{f})$ and
$v_0=\phi(v_1)=\phi(v_2)$.

{\em Case I (Contracting a loop). } In this case $v_1=v_2$. Let
$w_1,\ldots,w_n$ be the vertices of $\rho$ that map to
$v_0$ under $a$. Note that $g(w_i)\geq1$, since $g(v_0)\geq1$. Let $\pi$ be
equal to $\rho$ with a loop
$\{f_i,\ol{f}_i\}$ attached to $w_i$, for each $i=1,\ldots,n$ and
$g_\pi(w_i)=g_\rho(w_i)-1$. Clearly, $\pi$ is stable,
the drop in certain genera is made up for by the addition of flags. The
morphism $b:\pi\rightarrow\sigma$ is the obvious
combinatorial morphism mapping every one of the loops $\{f_i,\ol{f}_i\}$ to
$\{f,\ol{f}\}$. The contraction
$\psi:\pi\rightarrow\rho$ simply contracts all the added loops.

\begin{equation*}
\begin{array}{ccc}
\beginpicture%---pi
\setcoordinatesystem units <.3cm,.3cm> point at 3 2.5
\setplotarea x from 0 to 6, y from 0 to 5
\plot 3 4 4 4 /
\plot 3 4 4 3 /
\plot 3 1 4 1 /
\setquadratic
\plot 3 4 2 4.75 1 5 /
\plot 3 4 2 3.25 1 3 /
\plot 3 1 2 1.75 1 2 /
\plot 3 1 2 .25 1 0 /
\circulararc 180 degrees from 1 2 center at 1 1
\circulararc 180 degrees from 1 5 center at 1 4
\shaderectangleson
\setshadegrid span <1mm>
\putrectangle corners at 4 0 and 6 5
\put {\circle*{4}} [Bl] at 3 1
\put {\circle*{4}} [Bl] at 3 4
\axis left invisible label {$\scriptstyle\pi$} /
\axis right invisible label {\phantom{$\scriptstyle\pi$}} /
\axis bottom invisible label {\phantom{.}} /
\endpicture & \stackrel{\psi}{\longrightarrow} &
\beginpicture%---rho
\setcoordinatesystem units <.3cm,.3cm> point at 1.5 2.5
\setplotarea x from 0 to 3, y from 0 to 5
\plot 0 1 1 1 /
\plot 0 4 1 4 /
\plot 0 4 1 3 /
\shaderectangleson
\setshadegrid span <1mm>
\putrectangle corners at 1 0 and 3 5
\put {\circle*{4}} [Bl] at 0 1
\put {\circle*{4}} [Bl] at 0 4
\axis left invisible label {\phantom{$\scriptstyle\rho$}} /
\axis right invisible label {$\scriptstyle\rho$} /
\axis bottom invisible label {\phantom{.}} /
\endpicture \\
\ldiag{b} & & \rdiag{a} \\
\beginpicture%---sigma
\setcoordinatesystem units <.3cm,.3cm> point at 3 2
\setplotarea x from 0 to 6, y from 0 to 4
\plot 3 2 4 3 /
\plot 3 2 4 2 /
\plot 3 2 4 1 /
\setquadratic
\plot 3 2 2 2.75 1 3 /
\plot 3 2 2 1.25 1 1 /
\circulararc 180 degrees from 1 3 center at 1 2
\shaderectangleson
\setshadegrid span <1mm>
\putrectangle corners at 4 0 and 6 4
\put {\circle*{4}} [Bl] at 3 2
\axis left invisible label {$\scriptstyle\sigma$} /
\axis right invisible label {\phantom{$\scriptstyle\sigma$}} /
\axis top invisible label {\phantom{.}} /
\endpicture & \stackrel{\phi}{\longrightarrow} &
\beginpicture%---tau
\setcoordinatesystem units <.3cm,.3cm> point at 1.5 2
\setplotarea x from 0 to 3, y from 0 to 4
\plot 0 2 1 3 /
\plot 0 2 1 2 /
\plot 0 2 1 1 /
\shaderectangleson
\setshadegrid span <1mm>
\putrectangle corners at 1 0 and 3 4
\put {\circle*{4}} [Bl] at 0 2
\axis left invisible label {\phantom{$\scriptstyle\tau$}} /
\axis right invisible label {$\scriptstyle\tau$} /
\axis top invisible label {\phantom{.}} /
\endpicture
\end{array}
\end{equation*}

{\em Case II (Contracting a non-looping edge). } In this case $v_1\not=v_2$.
Again, let $w_1,\ldots, w_n$ be the
vertices of $\rho$ that map to $v_0$ under $a$. First we shall construct an
intermediate graph $\pi'$. Let us fix an
$i=1,\ldots,n$. Construct $\pi'$ from $\rho$ by replacing $w_i$ with two
vertices $w_i'$ and $w_i''$, connected by an
edge $\{f_i,\ol{f}_i\}$, such that $\del(f_i)=w_i'$ and $\del(\ol{f}_i)=w_i''$.
Let $f$ be a flag of $\rho$ such that
$\del_{\rho}(f)=w_i$. If $\del\s\phi^F(a(f))=v_1$ we attach $f$ to $w_i'$ and
if $\del\s\phi^F(a(f))=v_2$ we attach $f$
to $w_i''$. Set $g(w_i')=g(v_1)$, $g(w_i'')=g(v_2)$,
$\beta(w_i')=\xi(\alpha(v_1))$ and
$\beta(w_i'')=\xi(\alpha(v_2))$. This defines the $B$-graph $\pi'$. The problem
with $\pi'$ is that it might not be
stable. So to construct $\pi$ we proceed as follows. Fix an $i=1,\ldots,n$. If
$w_i'$ and $w_i''$ are stable vertices of
$\pi'$ we do not change anything. If either of $w_i'$ or $w_i''$ is unstable,
we go back to where we started, by
contracting $\{f_i,\ol{f}_i\}$ again, obtaining the stable vertex $w_i$. This
finally finishes the construction of
$\pi$. The contraction $\psi:\pi\rightarrow\rho$ is defined by contracting all
the edges that were just inserted into
$\rho$ to construct $\pi$. There is an obvious combinatorial morphism
$b':\pi'\rightarrow\sigma$ mapping the edge
$\{f_i,\ol{f}_i\}$ to $\{f,\ol{f}\}$. Moreover, we define a combinatorial
morphism $c:\pi\rightarrow\pi'$ as follows. If
$i=1,\ldots,n$ is an index such that either of $w_i'$ or $w_i''$ is an unstable
vertex of $\pi'$, we map the vertex
$w_i$ of $\pi$ to the stable one of the two, say $w_i'$, to fix notation. If
$f$ is a flag of $\rho$ such that
$\del{f}=w_i$, then $f$ is also considered as a flag of $\pi$ and $\pi'$, and
under $c$ we map $f$ to itself, if
$\del_{\pi'}(f)=w_i'$, and to $f_i$, otherwise. Finally,
$b:\pi\rightarrow\sigma$ is defined as the composition
$b=b'\comp c$.

\begin{equation} \label{tadsp}
\begin{array}{ccc}
\beginpicture%---pi
\setcoordinatesystem units <.3cm,.3cm> point at 4.5 3
\setplotarea x from 0 to 9, y from 0 to 6
\plot 2 5 7 5 /
\plot 2 4 3 5 /
\plot 6 5 7 4 /
\plot 2 3 3 3 /
\plot 2 1 7 1 /
\plot 6 1 7 2 /
\shaderectangleson
\setshadegrid span <1mm>
\putrectangle corners at 0 0 and 2 6
\putrectangle corners at 7 0 and 9 6
\put {\circle*{4}} [Bl] at 3 5
\put {\circle*{4}} [Bl] at 6 5
\put {\circle*{4}} [Bl] at 3 3
\put {\circle*{4}} [Bl] at 6 1
\axis left invisible label {$\scriptstyle\pi$} /
\axis right invisible label {\phantom{$\scriptstyle\pi$}} /
\axis bottom invisible label {\phantom{.}} /
\endpicture & \stackrel{\psi}{\longrightarrow} &
\beginpicture%---rho
\setcoordinatesystem units <.3cm,.3cm> point at 3 3
\setplotarea x from 0 to 6, y from 0 to 6
\plot 2 5 4 5 /
\plot 2 4 3 5 /
\plot 3 5 4 4 /
\plot 3 3 2 3 /
\plot 2 1 4 1 /
\plot 3 1 4 2 /
\shaderectangleson
\setshadegrid span <1mm>
\putrectangle corners at 0 0 and 2 6
\putrectangle corners at 4 0 and 6 6
\put {\circle*{4}} [Bl] at 3 5
\put {\circle*{4}} [Bl] at 3 3
\put {\circle*{4}} [Bl] at 3 1
\axis left invisible label {\phantom{$\scriptstyle\rho$}} /
\axis right invisible label {$\scriptstyle\rho$} /
\axis bottom invisible label {\phantom{.}} /
\endpicture \\
\ldiag{c} & & \\
\beginpicture%---pi'
\setcoordinatesystem units <.3cm,.3cm> point at 4.5 3
\setplotarea x from 0 to 9, y from 0 to 6
\plot 2 5 7 5 /
\plot 2 4 3 5 /
\plot 6 5 7 4 /
\plot 2 3 6 3 /
\plot 2 1 7 1 /
\plot 6 1 7 2 /
\shaderectangleson
\setshadegrid span <1mm>
\putrectangle corners at 0 0 and 2 6
\putrectangle corners at 7 0 and 9 6
\put {\circle*{4}} [Bl] at 3 5
\put {\circle*{4}} [Bl] at 6 5
\put {\circle*{4}} [Bl] at 3 3
\put {\circle*{4}} [Bl] at 6 3
\put {\circle*{4}} [Bl] at 3 1
\put {\circle*{4}} [Bl] at 6 1
\axis left invisible label {$\scriptstyle\pi'$} /
\axis right invisible label {\phantom{$\scriptstyle\pi'$}} /
\axis bottom invisible label {\phantom{.}} /
\axis top invisible label {\phantom{.}} /
\endpicture & & \rdiag{a} \\
\ldiag{b'} & & \\
\beginpicture%---sigma
\setcoordinatesystem units <.3cm,.3cm> point at 4.5 2
\setplotarea x from 0 to 9, y from 0 to 4
\plot 2 2 7 2 /
\plot 2 3 3 2 2 1 /
\plot 7 3 6 2 7 1 /
\shaderectangleson
\setshadegrid span <1mm>
\putrectangle corners at 0 0 and 2 4
\putrectangle corners at 7 0 and 9 4
\put {\circle*{4}} [Bl] at 3 2
\put {\circle*{4}} [Bl] at 6 2
\axis left invisible label {$\scriptstyle\sigma$} /
\axis right invisible label {\phantom{$\scriptstyle\sigma$}} /
\axis top invisible label {\phantom{.}} /
\endpicture & \stackrel{\phi}{\longrightarrow} &
\beginpicture%---tau
\setcoordinatesystem units <.3cm,.3cm> point at 3 2
\setplotarea x from 0 to 6, y from 0 to 4
\plot 2 2 4 2 /
\plot 2 3 3 2 2 1 /
\plot 4 3 3 2 4 1 /
\shaderectangleson
\setshadegrid span <1mm>
\putrectangle corners at 0 0 and 2 4
\putrectangle corners at 4 0 and 6 4
\put {\circle*{4}} [Bl] at 3 2
\axis left invisible label {\phantom{$\scriptstyle\tau$}} /
\axis right invisible label {$\scriptstyle\tau$} /
\axis top invisible label {\phantom{.}} /
\endpicture
\end{array}
\end{equation}

Iterating this construction leads to the construction of a stable pullback for
arbitrary contractions of $A$-graphs.

\begin{rmk}
Let
\[\begin{array}{ccccc}
B & & \pi & \stackrel{\psi}{\longrightarrow} & \rho \\
\ldiagup{\xi} & \phantom{\longrightarrow} & \ldiag{b} &  & \rdiag{a} \\
A & & \sigma & \stackrel{\phi}{\longrightarrow} & \tau
\end{array}\]
be a stable pullback.
\begin{enumerate}
\item The diagram
\[\comdia{V_\pi}{\psi_V}{V_\rho}{b}{}{a}{V\s}{\phi_V}{V\t}\]
commutes.
\item The diagram
\[\comdiaback{F_\pi}{\psi^F}{F_\rho}{b}{}{a}{F\s}{\phi^F}{F\t}\]
does {\em not }commute (except for special cases, e.g.\ if the $B$-graph $\pi'$
constructed above is stable).
\end{enumerate}
\end{rmk}

\begin{prop} \label{cspcn}
Stable pullback is independent of the order in which $\phi$ is decomposed into
elementary contractions.
Moreover, if
\[\begin{array}{ccccc}
B & & \pi & \stackrel{\psi}{\longrightarrow} & \rho \\
\ldiagup{\xi} & \phantom{\longrightarrow} & \ldiag{b} &  & \rdiag{a} \\
A & & \sigma & \stackrel{\phi}{\longrightarrow} & \tau
\end{array}\]
and
\[\begin{array}{ccccc}
B & & \pi' & \stackrel{\psi'}{\longrightarrow} & \pi \\
\ldiagup{\xi} & \phantom{\longrightarrow} & \ldiag{b'} &  & \rdiag{b} \\
A & & \sigma' & \stackrel{\phi'}{\longrightarrow} & \sigma
\end{array}\]
are stable pullbacks, then
\[\begin{array}{ccccc}
B & & \pi' & \stackrel{\psi\comp\psi'}{\longrightarrow} & \rho \\
\ldiagup{\xi} & \phantom{\longrightarrow} & \ldiag{b'} &  & \rdiag{a} \\
A & & \sigma' & \stackrel{\phi\comp\phi'}{\longrightarrow} & \tau
\end{array}\]
is a stable pullback, too.
\end{prop}
\begin{pf}
To check that stable pullback is well-defined, it suffices by
Remark~\ref{rmc}(\ref{rmcone}) to check that the above
construction yields the same result for both orders in which two elementary
contractions can be composed. This is a
straightforward, though maybe slightly tedious calculation.  The compatibility
of stable pullback with compositions of
contractions follows trivially from the definition.
\end{pf}

\begin{prop} \label{cspcb}
If
\[\begin{array}{ccccc}
B & & \pi & \stackrel{\psi}{\longrightarrow} & \rho \\
\ldiagup{\xi} & \phantom{\longrightarrow} & \ldiag{b} &  & \rdiag{a} \\
A & & \sigma & \stackrel{\phi}{\longrightarrow} & \tau
\end{array}\]
and
\[\begin{array}{ccccc}
C & & \pi' & \stackrel{\chi}{\longrightarrow} & \rho' \\
\ldiagup{\eta} & \phantom{\longrightarrow} & \ldiag{b'} &  & \rdiag{a'} \\
B & & \pi & \stackrel{\psi}{\longrightarrow} & \rho
\end{array}\]
are stable pullbacks, then
\[\begin{array}{ccccc}
C & & \pi' & \stackrel{\chi}{\longrightarrow} & \rho' \\
\ldiagup{\eta\comp\xi} & \phantom{\longrightarrow} & \ldiag{b\comp b'} &  &
\rdiag{a\comp a'} \\
A & & \sigma & \stackrel{\phi}{\longrightarrow} & \tau
\end{array}\]
is a stable pullback, too.
\end{prop}
\begin{pf}
Of course, it suffices to consider the case that $\phi$ is an elementary
contraction. Then the claim follows immediately
from the construction.
\end{pf}


We are now ready to define the notion of {\em morphism } of marked stable
graphs.

\begin{defn} \label{dmmsg}
Let $(A,\tau)$ and $(B,\sigma)$ be marked stable graphs. A {\em morphism } from
$(A,\tau)$ to $(B,\sigma)$ is a
quadruple $(\xi,a,\tau',\phi)$, where $\xi:A\rightarrow B$ is a homomorphism of
semigroups, $\tau'$ is a stable
$B$-graph, $a:\tau'\rightarrow\tau$ makes $(\xi,a)$ a combinatorial morphism of
marked graphs, and
$\phi:\tau'\rightarrow\sigma$ is a contraction of $B$-graphs. We also say that
$(a,\tau',\phi)$ is a morphism of marked
stable graphs {\em covering }$\xi$.

Let $(\xi,a,\tau',\phi):(A,\tau)\rightarrow(B,\sigma)$ and
$(\eta,b,\sigma',\psi):(B,\sigma)\rightarrow(C,\rho)$ be
morphisms of stable marked graphs. Then we define the {\em composition }
$(\eta,b,\sigma',\psi)\comp(\xi,a,\tau',\phi):(A,\tau)\rightarrow(C,\rho)$ to
be $(\eta\xi,a c,\tau'',\psi\chi)$, where
$(c,\tau'',\chi)$ is the stable pullback of $\sigma'$ under $\phi$.
\[\begin{array}{ccccccc}
C & & \tau'' & \stackrel{\chi}{\longrightarrow} & \sigma' &
\stackrel{\psi}{\longrightarrow} & \rho \\
\ldiagup{\eta} & \phantom{\longrightarrow} & \ldiag{c} &  & \rdiag{b} & & \\
B & & \tau' & \stackrel{\phi}{\longrightarrow} & \sigma & & \\
\ldiagup{\xi} & & \ldiag{a} & & & & \\
A & & \tau & & & &
\end{array}\]
\end{defn}

\begin{rmks}
\begin{enumerate}
\item In reality a morphism is an isomorphy class of quadruples as in this
definition. But we shall always stick to the
abuse of language begun here.
\item The composition of morphisms is associative by Propositions~\ref{cspcn}
and~\ref{cspcb}.
\item Every combinatorial morphism of marked graphs whose source and target are
stable defines a morphism of marked
stable graphs, but in the {\em opposite } direction.
\item Every contraction of $A$-graphs whose source (and hence target) is stable
defines a morphism of marked stable
graphs (in the same direction).
\end{enumerate}
\end{rmks}

The category of stable marked graphs shall be denoted by $\GG_s$. Let $\AA$ be
the category of (additive) semigroups with
indecomposable zero element. By projecting onto the first component, we get a
functor $\Aa:\GG_s\rightarrow\AA$. For
$A\in\ob\AA$ let $\GG_s(A)$ be the fiber of $\Aa$ over $A$, i.e.\ the category
of stable $A$-graphs.

\begin{prop} \label{ppes}
Let $\tau$ be an $A$-graph. Then there exists a stable $A$-graph $\tau^s$,
together with a combinatorial morphism
$\tau^s\rightarrow\tau$, such that every combinatorial morphism
$\sigma\rightarrow\tau$, where $\sigma$ is a stable
$A$-graph, factors uniquely through $\tau^s$. We call $\tau^s$ the {\em
stabilization  }of $\tau$.
\end{prop}
\begin{pf}
let $\alpha$ denote the $A$-structure of $\tau$.

{\em Case I\@. }Assume that $\tau$ has a vertex $v_0$ such that
 $g(v_0)=0$,
 $\alpha(v_0)=0$,
 $v_0$ has a unique flag $f_1$, and
 $f_2:=j(f_1)\not=f_1$.
Let $\tau'\rightarrow\tau$ be the `subgraph' defined by
 $F_{\tau'}=F\t-\{f_1\}$,
 $V_{\tau'}=V\t-\{v_0\}$,
 $\del_{\tau'}=\del\t\resto F_{\tau'}$,
 $j_{\tau'}\resto F_{\tau'}-\{f_2\}=j\t\resto F_{\tau'}-\{f_2\}$ and
$j_{\tau'}(f_2)=f_2$.

{\em Case II\@. }Assume that $\tau$ has a vertex $v_0$ such that
 $g(v_0)=0$,
 $\alpha(v_0)=0$,
 $v_0$ has exactly two flags, $f_1$ and $f_2$,
 $f_1$ is a tail of $\tau$ and
 $f_3:=j(f_2)\not=f_2$.
Let $\tau'\rightarrow\tau$ be the `subgraph' defined by
 $F_{\tau'}=F\t-\{f_1,f_2\}$,
 $V_{\tau'}=V\t-\{v_0\}$,
 $\del_{\tau'}=\del\t\resto F_{\tau'}$,
 $j_{\tau'}\resto F_{\tau'}-\{f_3\}=j\t\resto F_{\tau'}-\{f_3\}$ and
$j_{\tau'}(f_3)=f_3$.

{\em Case III\@. }Assume that $\tau$ has a vertex $v_0$ such that
 $g(v_0)=0$,
 $\alpha(v_0)=0$,
 $v_0$ has exactly two flags, $f_1$ and $f_2$,
 $\ol{f}_1:=j(f_1)\not=f_1$ and $\ol{f}_2:=j(f_2)\not=f_2$.
Let $\tau'\rightarrow\tau$ be the `subgraph' defined by
 $F_{\tau'}=F\t-\{f_1,f_2\}$,
 $V_{\tau'}=V\t-\{v_0\}$,
 $\del_{\tau'}=\del\t\resto F_{\tau'}$,
 $j_{\tau'}\resto F_{\tau'}-\{\ol{f}_1,\ol{f}_2\}=j\t\resto
F_{\tau'}-\{\ol{f}_1,\ol{f}_2\}$ and $j_{\tau'}(\ol{f}_1)=\ol{f}_2$.

{\em Case IV\@. }Assume that $\tau$ has a vertex $v_0$ such that
 $2g(v_0)+|v_0|<3$,
 $\alpha(v_0)=0$ and
 $F\t(v_0)$ is a union of orbits of $j\t$.
Let $\tau'\rightarrow\tau$ be the `subgraph' defined by
 $F_{\tau'}=F\t-F\t(v_0)$,
 $V_{\tau'}=V\t-\{v_0\}$,
 $\del_{\tau'}=\del\t\resto F_{\tau'}$ and
 $j_{\tau'}=j\t\resto F_{\tau'}$.

In each of these four cases every combinatorial morphism
$\sigma\rightarrow\tau$, with $\sigma$ stable factors uniquely
through $\tau'$. By induction on the number of vertices of $\tau$, the graph
$\tau'$ has a stabilization, which is thus
also a stabilization of $\tau$. If $\tau$ has no vertices $v_0$ of the kind
covered by the above four cases, $\tau$ is stable
and $\tau$ itself may serve as stabilization of $\tau$.
\end{pf}

See Section~10 in \cite[Exp.\ ~VI]{sga1}) for the definition of {\em
cofibration } of categories.

\begin{prop} \label{goscf}
The functor $\Aa:\GG_s\rightarrow\AA$ is a cofibration.
\end{prop}
\begin{pf}
Let $\xi:A\rightarrow B$ be a homomorphism in $\AA$, and $(\tau,\alpha)$ a
stable $A$-graph. We need to construct a
stable $B$-graph $\sigma=\xi\lst\tau$, together with a morphism
$(a,\tau',\phi):(A,\tau)\rightarrow(B,\sigma)$ covering
$\xi$, with the following universal mapping property. Whenever
$\eta:B\rightarrow C$ is another homomorphism in $\AA$,
$\rho$ is a stable $C$-graph and $(b,\tau'',\psi):(A,\tau)\rightarrow(C,\rho)$
is a morphism covering $\eta\comp\xi$,
there exists a unique morphism $(c,\sigma',\chi):(B,\sigma)\rightarrow(C,\rho)$
covering $\eta$, such that
$(c,\sigma',\chi)\comp(a,\tau',\phi)=(b,\tau'',\psi)$, i.e.\ such that $\tau''$
is the stable pullback of $\sigma'$
under $\phi$.

In fact, it is not difficult to see that the stabilization of
$(\tau,\xi\comp\alpha)$ satisfies this universal mapping
property.
\end{pf}

\begin{numrmk} \label{pfabs}
Choosing a {\em clivage normalis\'e }(see Definition~7.1 in \cite[Exp.\
{}~VI]{sga1}) of
$\GG_s$ over $\AA$ amounts to choosing a pushforward functor
$\xi\lst:\GG_s(A)\rightarrow\GG_s(B)$ for any homomorphism
$\xi:A\rightarrow B$ in $\AA$. We may call $\xi\lst$ {\em stabilization } with
respect to $\xi$. If $B=\{0\}$, we speak
of {\em absolute stabilization }(or simply {\em stabilization}, if no confusion
seems likely to arise).
\end{numrmk}


\section{Prestable Curves}

We recall the definition of prestable curves. A morphism of prestable curves is
defined in such a way that it has degree
at most one and contracts at most rational components.

\begin{defn} \label{dopsc}
A {\em prestable curve }over the scheme $T$ is a flat proper morphism
$\pi:C\rightarrow T$ of schemes such that the
geometric fibers of $\pi$ are reduced, {\em connected}, one-dimensional and
have at most ordinary double points (nodes)
as singularities.  The {\em genus } of a prestable curve $C\rightarrow T$ is
the map $t\mapsto\dim H^1(C_t,\O_{C_t})$,
which is a locally constant function $g:T\rightarrow\zz_{\geq0}$. If $L$ is a
line bundle on $C$, then the {\em degree }
of $L$ is the locally constant function $\deg L:T\rightarrow\zz_{\geq0}$ given
by $t\mapsto\chi(L_t)+g-1$.

A {\em morphism }$p:C\rightarrow D$ of prestable curves over $T$ is a
$T$-morphism of schemes, such that for every
geometric point $t$ of $T$ we have
\begin{enumerate}
\item if $\eta$ is the generic point of an irreducible component of $D_t$, then
the fiber of $p_t$ over $\eta$ is a finite
$\eta$-scheme of degree at most one,
\item if $C'$ is the normalization of an irreducible component of $C_t$, then
$p_t(C')$ is a single point only if $C'$
is rational.
\end{enumerate}
\end{defn}

If $V$ is a scheme and $f:C\rightarrow V$ a morphism, then $L\mapsto\deg f\upst
L$ defines a locally constant function
$T\rightarrow\Hom_\zz(\Pic V,\zz)$ which we shall call the {\em homology class
} of $f$, by abuse of language, denoted
$f\lst[C]$.

If $V$ is a scheme admitting an ample invertible sheaf let
\[H_2(V)^+=\{\alpha\in\Hom_\zz(\Pic V,\zz)\st \text{$\alpha(L)\geq0$ whenever
$L$ is ample}\}.\]
Note that $H_2(V)^+$ is a semigroup with indecomposable zero. This is because
if $V$ admits an ample invertible sheaf
then $\Pic V$ is generated by ample invertible sheaves (see Remarque~4.5.9 in
\cite{ega2}). So if $f:C\rightarrow V$ is
a morphism from a prestable curve into $V$, then the homology class is a
locally constant function $T\rightarrow
H_2(V)^+$.

\begin{lem} \label{tloq}
Let $f:X\rightarrow Y$ be a proper surjective morphism of $T$-schemes such that
$f\lst\O_X=\O_Y$. Let $g:X\rightarrow U$
be another morphism of $T$-schemes, such that for every geometric point $t$ of
$T$ the map $g_t:X_t\rightarrow U_t$ is
constant (as a map of underlying Zariski topological spaces) on the fibers of
$f_t:X_t\rightarrow Y_t$. Then $g$ factors
uniquely through $f$.
\end{lem}
\begin{pf}
This follows easily, for example, from Lemma~8.11.1 in \cite{ega2}.
\end{pf}

\begin{cor} \label{zc}
Let $C$ be a prestable curve over $T$ and $f:C\rightarrow V$ a morphism, where
$V$ is a scheme admitting
an ample invertible sheaf. Then $f\lst[C]=0$ if and only if $f$ factors through
$T$.  \qed
\end{cor}

We shall need the following two results about gluing marked prestable curves at
the marks.
%
\begin{prop} \label{glue}
Let $T$ be a scheme and $C_1$, $C_2$ two prestable curves over $T$. Let $x_1\in
C_1(T)$ and $x_2\in C_2(T)$ be sections
such that for every geometric point $t$ of $T$ we have that $x_1(t)$ and
$x_2(t)$ are in the smooth locus of $C_{1,t}$
and $C_{2,t}$, respectively. Then there exists a prestable curve $C$ over $T$,
together with $T$-morphisms
$p_1:C_1\rightarrow C$ and $p_2:C_2\rightarrow C$, such that
\begin{enumerate}
\item $p_1(x_1)=p_2(x_2)$,
\item $C$ is universal among all $T$-schemes with this property.
\end{enumerate}
The curve $C$ is uniquely determined (up to unique isomorphism) and will be
called {\em obtained by gluing $C_1$ and
$C_2$ along the sections $x_1$ and $x_2$}, notation
\[C=C_1\amalg_{x_1,x_2}C_2.\]

If $u:S\rightarrow T$ is a morphism of schemes, then $C_S$ is the curve
obtained by gluing $C_{1,S}$ and $C_{2,S}$ along
$x_{1,S}$ and $x_{2,S}$. If $g_i$ is the genus of $C_i$, for $i=1,2$, then for
the genus $g$ of $C$ we have $g=g_1+g_2$.
If, for $i=1,2$, $f_i:C_i\rightarrow V$ is a morphism into a scheme such that
$f_1(x_1)=f_2(x_2)$, and $f:C\rightarrow
V$ is the induced morphism, we have $f\lst[C]={f_1}\lst[C_1]+{f_2}\lst[C_2]$ in
$\Hom_\zz(\Pic V,\zz)$.  \qed
\end{prop}
%
\begin{prop} \label{glue1}
Let $T$ be a scheme and $C$ a prestable curve over $T$. Let $x_1\in C(T)$ and
$x_2\in C(T)$ be sections
such that for every geometric point $t$ of $T$ we have that $x_1(t)$ and
$x_2(t)$ are in the smooth locus of $C_{t}$ and
$x_1(t)\not= x_2(t)$. Then there exists a prestable curve $\tilde{C}$ over $T$,
together with a $T$-morphism
$p:C\rightarrow \tilde{C}$, such that
\begin{enumerate}
\item $p(x_1)=p(x_2)$,
\item $\tilde{C}$ is universal among { all }$T$-schemes with this property.
\end{enumerate}
The curve $\tilde{C}$ is uniquely determined (up to unique isomorphism) and
will be called {\em obtained by gluing $C$
with itself along the sections $x_1$ and $x_2$}, notation
\[\tilde{C}=C/{x_1\sim x_2}.\]

If $u:S\rightarrow T$ is a morphism of schemes, then $(\tilde{C})_S$ is the
curve obtained by gluing $C_{S}$ with itself along
$x_{1,S}$ and $x_{2,S}$. If $g$ is the genus of $C$, then for the genus
$\tilde{g}$ of $\tilde{C}$ we have $\tilde{g}=g+1$.
If $f:C\rightarrow V$ is a morphism into a scheme such that $f(x_1)=f(x_2)$,
and $\tilde{f}:\tilde{C}\rightarrow
V$ is the induced morphism, we have $\tilde{f}\lst[\tilde{C}]={f}\lst[C]$ in
$\Hom_\zz(\Pic V,\zz)$. \qed
\end{prop}

\begin{defn} \label{mpc}
Let $\tau$ be a modular graph. A {\em $\tau$-marked prestable curve over $T$}
is a pair $(C,x)$, where $C=(C_v)_{v\in
V\t}$ is a family of prestable curves $\pi_v:C_v\rightarrow T$ and
$x=(x_i)_{i\in F\t}$ is a family of sections
$x_i:T\rightarrow C_{\del\t(i)}$, such that for every geometric point $t$ of
$T$ we have
\begin{enumerate}
\item $x_i(t)$ is in the smooth locus of $C_{\del\t(i),t}$, for all $i\in F\t$,
\item $x_i(t)\not=x_j(t)$, if $i\not=j$, for $i,j\in F\t$,
\item $g(C_{v,t})=g(v)$ for all $v\in V\t$.
\end{enumerate}

We define a {\em marked prestable curve over $T$} to be a triple $(\tau,C,x)$,
where $\tau$ is a modular graph and
$(C,x)$ a $\tau$-marked prestable curve over $T$.
\end{defn}

\section{Stable Maps}

We now come to the definition of stable maps, the central concept of this work,
which is due to Kontsevich.

Fix a field $k$ and let $\VV$ be the category of smooth projective (not
necessarily connected) varieties over $k$.
Consider the covariant functor
\begin{eqnarray*}
H_2^+:\VV & \longrightarrow & \AA \\
V & \longmapsto & H_2(V)^+,
\end{eqnarray*}
where $\AA$ is the category of semigroups with indecomposable zero (see
Section~\ref{graphs}).
Define the category $\VV\GG_s$ as the fibered product (see Section~3 in
\cite[Exp.\ ~VI]{sga1})
\[\comdia{\VV\GG_s}{}{\GG_s}{}{\Box}{\Aa}{\VV}{H_2^+}{\AA.}\]
To spell this definition out, we have
\begin{enumerate}
\item objects of $\VV\GG_s$ are pairs $(V,\tau)$, where $V$ is a smooth
projective variety over $k$
and $\tau$ is a stable $H_2(V)^+$-graph,
\item a morphism $(V,\tau)\rightarrow (W,\sigma)$ is a quadruple
$(\xi,a,\tau',\phi)$, where
$\xi:V\rightarrow W$ is a morphism of $k$-varieties and
$(H_2^+(\xi),a,\tau',\phi)$ is a morphism in $\GG_s$ as defined in
Definition~\ref{dmmsg}.
\end{enumerate}

\begin{numrmk} \label{ugscf}
By Corollary~6.9 of \cite[Exp.\ ~VI]{sga1} and Proposition~\ref{goscf} the
category $\VV\GG_s$ is a cofibered category
over $\VV$.
\end{numrmk}

\begin{defn} \label{dsm}
Let $(V,\tau,\alpha)$ be an object of $\VV\GG_s$ and $T$ a $k$-scheme. A {\em
stable $(V,\tau,\alpha)$-map over $T$} is
a triple $(C,x,f)$, where $(C,x)$ is a $\tau$-marked prestable curve over $T$
and $f=(f_v)_{v\in V\t}$ is a family of
$k$-morphisms $f_v:C_v\rightarrow V$, such that the following conditions are
satisfied.
\begin{enumerate}
\item For every $i\in F\t$ we have
$f_{\del\t(i)}(x_i)=f_{\del\t(j\t(i))}(x_{j\t(i)})$ as $k$-morphisms from $T$
to $V$.
\item For all $v\in V\t$ we have that ${f_v}\lst[C_{v}]=\alpha(v)$ in
$H_2(V)^+$.
\item For every geometric point $t$ of $T$ and every $v\in V\t$ the {\em
stability condition }is satisfied. This means that if
$C'$ is the normalization of a component of $C_{v,t}$ that maps to a point
under $f_{v,t}:C_{v,t}\rightarrow V_t$, then
\begin{enumerate}
\item if the genus of $C'$ is zero, then $C'$ has at least three special
points,
\item if the genus of $C'$ is one, then $C'$ has at least one special point.
\end{enumerate}
Here, a point of $C'$ is called {\em special}, if it maps in $C_{v,t}$ to a
marked point or a node.
\end{enumerate}

We define a {\em stable map over $T$} to be a sextuple $(V,\tau,\alpha,C,x,f)$,
where $(V,\tau,\alpha)$ is an object
of $\VV\GG_s$ and $(C,x,f)$ is a stable $(V,\tau,\alpha)$-map over $T$.

A {\em morphism }$(V,\tau,\alpha,C,x,f)\rightarrow (W,\sigma,\beta,D,y,h)$ of
stable maps over $T$ is a quintuple
$(\xi,a,\tau',\phi,p)$, where $(\xi,a,\tau',\phi):(V,\tau,\alpha)\rightarrow
(W,\sigma,\beta)$ is a morphism in
$\VV\GG_s$ and $p=(p_v)_{v\in V_{\tau'}}$ is a family of morphisms of prestable
curves $p_v:C_{a(v)}\rightarrow
D_{\phi_V(v)}$, such that the following are true.
\begin{enumerate}
\item \label{mpc1}For every $i\in F\s$ we have
$p_{\del(\phi^F(i))}(x_{a\phi^F(i)})=y_i$,
\item \label{pre} If $\{i_1,i_2\}$ is an edge of $\tau'$ which is being
contracted by $\phi$, then
$p_{v_1}(x_{a(i_1)})=p_{v_2}(x_{a(i_2)})$, where $v_1=\del i_1$ and $v_2=\del
i_2$. So, in particular, if $v_1\not=v_2$
there exists an induces morphism
\[p_{12}:C_{a(v_1)}\amalg_{x_{a(i_1)},x_{a(i_2)}}C_{a(v_2)}\rightarrow D_w,\]
where $w=\phi(v_1)=\phi(v_2)$.
\item \label{post} With the notation of (\ref{pre}), if $v_1\not=v_2$, the
morphism $p_{12}$ is a morphism of prestable
curves.
\item For every $v\in V_{\tau'}$ the diagram
\[\comdia{C_{a(v)}}{f_{a(v)}}{V}{p_v}{}{\xi}{D_{\phi(v)}}{h_{\phi(v)}}{W}\]
commutes.
\end{enumerate}
In this situation we also say that $p:(C,x,f)\rightarrow(D,y,h)$ is a morphism
of stable maps {\em covering
}the morphism $(\xi,a,\tau',\alpha)$ in $\VV\GG_s$.

To define the {\em composition } of morphisms, let
$(\xi,a,\tau',\phi,p):(V,\tau,\alpha,C,x,f)\rightarrow(W,\sigma,\beta,D,y,h)$
and
%% FOLLOWING LINE CANNOT BE BROKEN BEFORE 80 CHAR
$(\eta,b,\sigma',\psi,q):(W,\sigma,\beta,D,y,h)\rightarrow(U,\rho,
\gamma,E,z,e)$ be morphisms of stable maps over
$T$. We already know how to compose the morphisms $(\xi,a,\tau',\phi)$ and
$(\eta,b,\sigma',\psi)$ in $\VV\GG_s$. Use
notation as in Definition~\ref{dmmsg}. Then this composition is
$(\eta\xi,ac,\tau'',\psi\chi)$. Define the family
$r=(r_u)_{u\in V_{\tau''}}$ of morphisms of prestable curves
$r_u:C_{ac(u)}\rightarrow E_{\psi\chi(u)}$ as
$r_u=q_{\chi(u)}\comp p_{c(u)}$, which is well-defined, since
$\phi_Vc(u)=a\chi_V(u)$. Then we define our composition as
%% FOLLOWING LINE CANNOT BE BROKEN BEFORE 80 CHAR
\[(\eta,b,\sigma',\psi,q)\comp(\xi,a,\tau',\phi,p) =
(\eta\xi,ac,\tau'',\psi\chi,r).\]
\end{defn}

\begin{prop} \label{cmsmm}
The composition of morphisms of stable maps is a morphism of stable maps.
\end{prop}
\begin{pf}
The proof will be given at the same time as the proof of Theorem~\ref{mbfc}
below.
\end{pf}

\begin{defn} \label{sovgs}
Let $V\in\ob\VV$ be a variety, $\beta\in H_2(V)^+$ a homology class and
$g,n\geq0$ integers. Then $(V,g,n,\beta)$ shall
denote the object $(V,\tau,\beta)$ of $\VV\GG_s$ whose modular graph $\tau$ is
given by $F\t=\ul{n}$,
$V\t=\{\varnothing\}$, $\del\t:F\t\rightarrow V\t$ the unique map,
$j\t=\id_{\ul{n}}$ and $g(\varnothing)=g$. The
$H_2(V)^+$-structure on $\tau$ is given by $\beta(\varnothing)=\beta$. A stable
$(V,g,n,\beta)$-map is also called a
{\em stable map from an $n$-pointed curve (of genus $g$) to $V$ (of class
$\beta$)}. Here we use the notation
$\ul{n}=\{1,\ldots,n\}$.
\end{defn}

\begin{lem} \label{lcsn}
Over an algebraically closed field, let $(C,x,f)$ be a stable map from an
$n$-pointed curve of genus $g$ to $V$ of class
$\beta$ and let $(D,y,h)$ be a stable map from an $m$-pointed curve of genus
$g$ to $V$ of class $\beta$, where $m\leq
n$. Let $p:C\rightarrow D$ be a morphism such that $p(x_i)=y_i$ for $i\leq m$
and $hp=f$. If $C'\subset C$ is a
subcurve (a connected union of irreducible components), such that
\begin{enumerate}
\item letting $C''$ be the closure of the complement of $C'$ in $C$, the curves
$C'$ and $C''$ have exactly one node in
common,
\item $g(C')=0$,
\item $f(C')$ is a point,
\item for $i\leq m$ the $x_i$ do not lie on $C'$ except for at most one of
them,
\end{enumerate}
then $p$ maps $C'$ to a point in $D$. \qed
\end{lem}

Let us denote the category of stable maps over $T$ by $\ol{M}(T)$. It comes
together with a functor
\[\ol{M}(T)\longrightarrow\VV\GG_s,\]
defined by projecting onto the first components. For a morphism $u:S\rightarrow
T$, pulling back defines a
$\VV\GG_s$-functor
\[u\upst:\ol{M}(T)\longrightarrow \ol{M}(S).\]

\begin{them} \label{mbfc}
For every $k$-scheme $T$ the functor $\ol{M}(T)\rightarrow\VV\GG_s$ is a
cofibration, whose fibers are groupoids. In
other words, $\ol{M}(T)$ is cofibered in groupoids over $\VV\GG_s$.

For every base change $u:S\rightarrow T$ the $\VV\GG_s$-functor
$u\upst:\ol{M}(T)\rightarrow \ol{M}(S)$ is cocartesian.
\end{them}
\begin{pf}
To prove that $\ol{M}(T)\rightarrow\VV\GG_s$ is a cofibration, we need to prove
the following. Let
$(\xi,a,\tau',\phi):(V,\tau)\rightarrow(W,\sigma)$ be a morphism in $\VV\GG_s$
and $(C,x,f)$ a stable $(V,\tau)$-map
over $T$. Then there exists a {\em pushforward }$(D,y,h)$ of $(C,x,f)$ under
$(\xi,a,\tau',\phi)$. This pushforward
comes with a morphism $p:(C,x,f)\rightarrow(D,y,h)$ of stable maps covering
$(\xi,a,\tau',\phi)$ and is characterized by
the following universal mapping property. Whenever
$(\eta,b,\sigma',\psi):(W,\sigma)\rightarrow(U,\rho)$ is another
morphism in $\VV\GG_s$, $(E,z,e)$ a stable $(U,\rho)$-map over $T$ and
$r:(C,x,f)\rightarrow(E,z,e)$ a morphism of
stable maps covering $(\eta\xi,ac,\tau'',\psi\chi):(V,\tau)\rightarrow(U,\rho)$
(in the notation of
Definition~\ref{dmmsg}), there exists a unique morphism of stable maps
$q:(D,y,h)\rightarrow(E,z,e)$ covering
$(\eta,b,\sigma',\psi):(W,\sigma)\rightarrow(U,\rho)$ such that $r=q\comp p$.
\begin{equation} \label{umprq}
\begin{array}{ccccc}
 & & \stackrel{r}{\overtoparrow} & & \\
(C,x,f) & \stackrel{p}{\longrightarrow} & (D,y,h) &
\stackrel{q}{\longrightarrow} & (E,z,e) \\
\mid & & \mid & & \mid \\
(V,\tau) & \stackrel{(\xi,a,\tau',\phi)}{\longrightarrow} & (W,\sigma) &
\stackrel{(\eta,b,\sigma',\psi)}{\longrightarrow} & (U,\rho) \\
 & & {\underbottomarrow\atop(\eta\xi,ac,\tau'',\psi\chi)} &  &
\end{array}
\end{equation}
To prove that $u\upst:\ol{M}(T)\rightarrow\ol{M}(S)$ is always cocartesian, we
need to prove that this pushforward
commutes with base change.

Recall that we also wish to prove Proposition~\ref{cmsmm}, i.e.\ that if
morphisms of stable maps
$p:(C,x,f)\rightarrow(D,y,h)$ and $q:(D,y,h)\rightarrow(E,z,e)$ are given as in
(\ref{umprq}), then the composition
$r:(C,x,f)\rightarrow(E,z,e)$ is also a morphism of stable maps.

Purely formal considerations tell us that to prove these three facts, we may
decompose the morphism
$(\xi,a,\tau',\phi):(V,\tau)\rightarrow(W,\sigma)$ into a composition of other
morphisms in any way we wish and prove the
three facts for the factors of this decomposition. We shall thus consider the
following five cases.

{\em Case I (Changing $V$). } In this case $\sigma=\xi\lst\tau$. This means
that $\sigma$ is the pushforward of $\tau$
under $\xi:V\rightarrow W$, using the fact that $\VV\GG_s\rightarrow\VV$ is a
cofibration (Remark~\ref{ugscf}). In other
words, $\sigma$ is the stabilization of $\tau$ with respect to the induced
$H_2(W)^+$-structure (Proposition~\ref{ppes}).
Thus $\tau'=\sigma$ and $\phi=\id\s$.

In all other cases $W=V$ and $\xi$ is the identity. In the next two cases
$a=\id\t$ and $\tau'=\tau$.

{\em Case II (Contracting and edge). } The contraction
$\phi:\tau\rightarrow\sigma$ contracts exactly one edge
$\{i_1,i_2\}\subset F\t$ and we have $v_1\not=v_2$, where $v_1=\del(i_1)$ and
$v_2=\del(i_2)$. To fix notation, let
$v_0=\phi(v_1)=\phi(v_2)$.

{\em Case III (Contracting a loop). } This is the same as Case~II, except that
we have $v_1=v_2$.

In the last two cases $\tau'=\sigma$ and $\phi=\id\s$.

{\em Case IV (Complete combinatorial). } The combinatorial morphism
$a:\sigma\rightarrow\tau$ has the property that
$a:F\s(v)\rightarrow F\t(a(v))$ is a bijection, for all $v\in V\s$.

{\em Case V (Removing a tail). } In this case, $V\s=V\t$, we have given a
vertex $v_0\in V\t$ and a {\em tail }$i_0\in
F\t(v_0)$ of $\tau$ and we have
\begin{enumerate}
\item $F\s= F\t-\{i_0\}$,
\item $\del\s=\del\t\resto F\s$,
\item $j\s=j\t\resto F\s$.
\end{enumerate}
\renewcommand{\qed}{}\end{pf}

Note that the proof of Proposition~\ref{cmsmm} is only interesting (if at all)
for Case~II, since only in this case
carrying out the composition of $(\xi,a,\tau',\phi)$ and
$(\eta,b,\sigma',\psi)$ involves the second case of the
construction of stable pullback (Section~\ref{graphs}).

{\em Case I\@. } First we note the following trivial lemma.

\begin{lem} \label{tcvw}
Assume that $\tau$ is stable with respect to the induced $H_2(W)^+$-structure,
so that $\sigma=\tau$ and $a=\id\t$. Then
if $(C,x,\xi\comp f)$ satisfies the stability condition it may serve as
pushforward of $(C,x,f)$ under $\xi$. \qed
\end{lem}

We shall now reduce Case~I to Cases~IV and~V\@. By the claimed compatibility
with base change, we may construct the
pushforward locally, and pass to an \'etale cover of $T$, whenever desirable.
Thus we add tails to $\tau$, obtaining
$\tilde{\tau}$, and corresponding sections of $C$, obtaining $(C,\tilde{x})$
until $\tilde{\tau}$ with the induced
$H_2(W)^+$-structure is stable and $(C,\tilde{x},\xi\comp f)$ satisfies the
stability condition. Then we have the
commutative diagram
\begin{equation} \label{divgs}
\comdia{(V,\tilde{\tau})}{}{(V,\tau)}{}{}{}{(W,\tilde{\tau})}{}{(W,\sigma)}
\end{equation}
in $\VV\GG_s$. The top row of (\ref{divgs}) is covered by
$(C,\tilde{x},f)\rightarrow(C,x,f)$, and clearly $(C,x,f)$ is
the pushforward of $(C,\tilde{x},f)$ under
$(V,\tilde{\tau})\rightarrow(V,\tau)$ (see also Case~V). The first column of
(\ref{divgs}) is covered by $(C,\tilde{x},f)\rightarrow(C,\tilde{x},\xi\comp
f)$, which is a pushforward by
Lemma~\ref{tcvw}. Now the pushforward of $(C,\tilde{x},\xi\comp f)$ under
$(W,\tilde{\tau})\rightarrow(W,\sigma)$ will
also be the sought after pushforward of $(C,x,f)$ under
$(V,\tau)\rightarrow(W,\sigma)$. But
$(W,\tilde{\tau})\rightarrow(W,\sigma)$ is covered by Cases~IV and~V, achieving
the reduction. \qed

{\em Case II\@. } The diagram defining the composition of $\phi$ and
$(b,\sigma',\psi)$ is
\[\begin{array}{ccccc}
\tau' & \stackrel{\chi}{\longrightarrow} & \sigma' &
\stackrel{\psi}{\longrightarrow} & \rho \\
\ldiag{c} & & \rdiag{b} & & \\
\tau & \stackrel{\phi}{\longrightarrow} & \sigma. & &
\end{array}\]
Let us first deal with the proof of Proposition~\ref{cmsmm}.

\begin{lem} \label{lc2}
For every $i\in F_{\sigma'}$ we have
\[q_{\del(i)}p_{\del
c\chi^F(i)}(x_{c\chi^F(i)})=q_{\del(i)}p_{\del\phi^Fb(i)}(x_{\phi^Fb(i)}).\]
\end{lem}
\begin{pf}
Assume that $c\chi^F(i)\not=\phi^Fb(i)$, since otherwise there is nothing to
prove. In this case, necessarily,
$c\chi^F(i)$ is being contracted by $\phi$. Without loss of generality, let
$c\chi^F(i)=i_1$, so the situation is as in
the following diagram (cf.\ ~(\ref{tadsp})).
\begin{equation} \label{fdgc}
\begin{array}{ccc}
\beginpicture%---tau'
\setcoordinatesystem units <.5cm,.3cm> point at 4.5 2
\setplotarea x from 0 to 9, y from 0 to 4
\plot 1.5 2 7.5 2 /
\plot 6 2 7.5 1 /
\shaderectangleson
\setshadegrid span <1mm>
\putrectangle corners at 0 0 and 1.5 4
\putrectangle corners at 7.5 0 and 9 4
\put {\circle*{4}} [Bl] at 6 2
\put {$\scriptstyle\chi^F(i)$} at 5.5 2.6
\axis left invisible label {$\scriptstyle\tau'$} /
\axis right invisible label {\phantom{$\scriptstyle\tau'$}} /
\axis bottom invisible label {\phantom{.}} /
\endpicture & \stackrel{\chi}{\longrightarrow} &
\beginpicture%---sigma'
\setcoordinatesystem units <.5cm,.3cm> point at 3 2
\setplotarea x from 0 to 6, y from 0 to 4
\plot 1.5 2 4.5 2 /
\plot 3 2 4.5 1 /
\shaderectangleson
\setshadegrid span <1mm>
\putrectangle corners at 0 0 and 1.5 4
\putrectangle corners at 4.5 0 and 6 4
\put {\circle*{4}} [Bl] at 3 2
\put {$\scriptstyle i$} at 2.5 2.6
\axis left invisible label {\phantom{$\scriptstyle\sigma'$}} /
\axis right invisible label {$\scriptstyle\sigma'$} /
\axis bottom invisible label {\phantom{.}} /
\endpicture \\
\ldiag{} & & \\
\beginpicture%---tilde{tau}'
\setcoordinatesystem units <.5cm,.3cm> point at 4.5 2
\setplotarea x from 0 to 9, y from 0 to 4
\plot 1.5 2 7.5 2 /
\plot 6 2 7.5 1 /
\shaderectangleson
\setshadegrid span <1mm>
\putrectangle corners at 0 0 and 1.5 4
\putrectangle corners at 7.5 0 and 9 4
\put {\circle*{4}} [Bl] at 3 2
\put {\circle*{4}} [Bl] at 6 2
\axis left invisible label {$\scriptstyle\tilde{\tau}'$} /
\axis right invisible label {\phantom{$\scriptstyle\tilde{\tau}'$}} /
\axis top invisible label {\phantom{.}} /
\axis bottom invisible label {\phantom{.}} /
\endpicture & & \rdiag{b} \\
\ldiag{} & & \\
\beginpicture%---tau
\setcoordinatesystem units <.5cm,.3cm> point at 4.5 2
\setplotarea x from 0 to 9, y from 0 to 4
\plot 1.5 2 7.5 2 /
\plot 3 2 1.5 1 /
\plot 7.5 3 6 2 7.5 1 /
\shaderectangleson
\setshadegrid span <1mm>
\putrectangle corners at 0 0 and 1.5 4
\putrectangle corners at 7.5 0 and 9 4
\put {\circle*{4}} [Bl] at 3 2
\put {\circle*{4}} [Bl] at 6 2
\put {$\scriptstyle\phi^Fb(i)$} at 2.5 2.6
\put {$\scriptstyle c\chi^F(i)$} at 5.5 2.6
\put {$\scriptstyle i_1$} at 5.5 1.4
\put {$\scriptstyle i_2$} at 3.5 1.4
\axis left invisible label {$\scriptstyle\tau$} /
\axis right invisible label {\phantom{$\scriptstyle\tau$}} /
\axis top invisible label {\phantom{.}} /
\endpicture & \stackrel{\phi}{\longrightarrow} &
\beginpicture%---sigma
\setcoordinatesystem units <.5cm,.3cm> point at 3 2
\setplotarea x from 0 to 6, y from 0 to 4
\plot 1.5 2 4.5 2 /
\plot 3 2 1.5 1 /
\plot 4.5 3 3 2 4.5 1 /
\shaderectangleson
\setshadegrid span <1mm>
\putrectangle corners at 0 0 and 1.5 4
\putrectangle corners at 4.5 0 and 6 4
\put {\circle*{4}} [Bl] at 3 2
\put {$\scriptstyle b(i)$} at 2.5 2.6
\axis left invisible label {\phantom{$\scriptstyle\sigma$}} /
\axis right invisible label {$\scriptstyle\sigma$} /
\axis top invisible label {\phantom{.}} /
\endpicture
\end{array}
\end{equation}
Here, ${\tau}'$ is the stable pullback and $\tilde{\tau}'$ the intermediate
graph used in the construction of
$\tau'$. Using the fact that $p$ is a morphism of stable maps we get a morphism
$p_{12}:C_{12}\rightarrow
D_{v_0}$ of prestable curves, where
\[C_{12}=C_{v_1}\amalg_{x_{i_1},x_{i_2}}C_{v_2}.\]
Compose this with $q_{\del(i)}:D_{v_0}\rightarrow E_{\psi\del(i)}$. Let
$f_{12}:C_{12}\rightarrow V$ be the map induced
from $f_{v_1}$ and $f_{v_2}$ and $\tilde{x}=x\resto F\t(v_1)\cup
F\t(v_2)-\{i_1,i_2\}$. Then $(C_{12},\tilde{x},f_{12})$
is a stable map and
\[q_{\del(i)}\comp
p_{12}:(C_{12},\tilde{x},f_{12})\longrightarrow(E_{\psi\del(i)},z\resto
F_\rho(\psi\del(i)),e_{\psi\del(i)})\] is a morphism of stable maps to which
Lemma~\ref{lcsn} applies, with $C'=C_{v_2}$
and $x_{\phi^Fb(i)}\in C'$ being the only marked point coming from
$F_\rho(\psi\del(i))$, if there exists such a point
at all (this is because $\tau'\not=\tilde{\tau}'$). So by Lemma~\ref{lcsn}
$q_{\del(i)}p_{v_2}(C_{v_2})$ is a point in
$E_{\psi\del(i)}$. To be precise, this holds if $T$ is the spectrum of an
algebraically closed field. For the general
case, applying Lemma~\ref{tloq} yields that $q_{\del(i)}\comp p_{v_2}$ factors
thought $T$. In particular,
%% FOLLOWING LINE CANNOT BE BROKEN BEFORE 80 CHAR
\[q_{\del(i)}p_{v_2}(x_{\phi^Fb(i)})= q_{\del(i)}p_{v_2}(x_{i_2})=
q_{\del(i)}p_{v_1}(x_{i_1}),\]
which is what we set out to prove.
\end{pf}

Let us check that $r:(C,x,f)\rightarrow(E,z,e)$ is a morphism of stable maps,
i.e.\ satisfies Properties~(1) through~(4)
from Definition~\ref{dsm}.

{\em Property (1). }  Let $i\in F_\rho$. The we have
\begin{align*}
r_{\del\chi^F\psi^F(i)}(x_{c\chi^F\psi^F(i)}) &=  q_{\del\psi^F(i)}p_{\del
c\chi^F\psi^F(i)}(x_{c\chi^F\psi^F(i)}) \\
\intertext{by Definition~\ref{dsm},}
 &=  q_{\del\psi^F(i)}p_{\del\phi^Fb\psi^F(i)}(x_{\phi^Fb\psi^F(i)}) \\
\intertext{by Lemma~\ref{lc2},}
 &=  q_{\del\psi^F(i)}(y_{b\psi^F(i)}) \\
 &=  z_i,
\end{align*}
since $p$ and $q$ are morphisms of stable maps.

{\em Property (2). }  Let $\{j_1,j_2\}$ be an edge of $\tau'$ which is being
contracted by $\psi\chi$. Let
$u_1=\del j_1$ and $u_2=\del j_2$.

{\em Case 1. } Let $\{j_1,j_2\}$ be contracted by $\chi$. Then
$\{c(j_1),c(j_2)\}$ is being contracted by $\phi$. So
without loss of generality $c(j_1)=i_1$ and $c(j_2)=i_2$. Then
\begin{align*}
r_{u_1}(x_{i_1}) & =  q_{\chi(u_1)}p_{v_1}(x_{i_1}) \\
 & =  q_{\chi(u_2)}p_{v_2}(x_{i_2}) \\
 & =  r_{u_2}(x_{i_2}),
\end{align*}
since $p$ is a morphism of stable maps and $\chi(u_1)=\chi(u_2)$.

{\em Case 2. } If $\{j_1,j_2\}$ is not contracted by $\chi$, then there exists
a unique edge $\{j_1',j_2'\}$ of
$\sigma'$ being contracted by $\psi$, such that $j_1=\chi^F(j_1')$ and
$j_2=\chi^F(j_2')$. Then
\begin{align*}
r_{u_1}(x_{c(j_1)}) & =  q_{\chi(u_1)}p_{c(u_1)}(x_{c(j_1)}) \\
 & =  q_{\chi(u_1)}p_{\del\phi^Fb(j_1')}(x_{\phi^Fb(j_1')}) \\
\intertext{by Lemma~\ref{lc2},}
 & =  q_{\chi(u_1)}(y_{b(j_1')}) \\
 & =  q_{\chi(u_2)}(y_{b(j_2')}) \\
\intertext{since $q$ is a morphism of stable maps,}
 & =  r_{u_2}(x_{c(j_2)}),
\end{align*}
by symmetry.

{\em Property (3). } This follows from the fact that the composition of
morphisms of prestable curves is again
a morphism of prestable curves.

{\em Property (4). } Straightforward.

This finishes the proof of Proposition~\ref{cmsmm} in Case~II\@. Let us now
construct the pushforward $(D,y,h)$ of
$(C,x,f)$ under $\phi$.

Let $w\in V\s$. If $w\not=v_0$, let $v$ be the unique vertex $v\in V\t$ such
that $\phi_V(v)=w$ and set $D_w=C_v$. If
$w=v_0$ set
\[D_{v_0}=C_{v_1}\amalg_{x_{i_1},x_{i_2}}C_{v_2}.\]
This defines a family of prestable curves $D$. For every $v\in V\t$ let
$p_v:C_v\rightarrow D_{\phi(v)}$ be the canonical
map. Define sections $y_i$, for $i\in F\s$, by
\[y_i=p_{\del\phi^F(i)}(x_{\phi^F(i)}).\]
Finally, define for every $w\in V\s$ a map $g_w:D_w\rightarrow V$ from $f$ (by
using Proposition~\ref{glue}, if
$w=v_0$). Essentially by definition, $(D,y,h)$ is a stable $(V,\sigma)$-map and
$p:(C,x,f)\rightarrow(D,y,h)$ is a
morphism of stable maps covering $\phi:(V,\tau)\rightarrow(V,\sigma)$. It
remains to check that $(D,y,h)$ satisfies the
universal mapping property of a pushforward under $\phi$. So let
$r:(C,x,f)\rightarrow(E,z,e)$ as in
Diagram~(\ref{umprq}) be given.

Let $u\in V_{\sigma'}$. We need to define a unique morphism
$q_u:D_{b(u)}\rightarrow E_{\psi(u)}$ such that for every
$u'\in V_{\tau'}$, satisfying $\chi(u')=u$, the diagram
\[\begin{array}{ccc}
C_{c(u')} & & \\
\ldiag{p_{c(u')}} & \sediagr{r_{u'}} & \\
D_{b(u)} & \stackrel{q_u}{\longrightarrow} & E_{\psi(u)}
\end{array}\]
commutes. If $b(u)\not=v_0$, necessarily, $q_u=r_{u'}$. So let $b(u)=v_0$. If
there are two vertices $u_1'$ and $u_2'$
such that $\chi(u_1')=\chi(u_2')=u$, then we have two maps
$r_{u_1'}:C_{v_1}\rightarrow E_{\psi(u)}$ and
$r_{u_2'}:C_{v_2}\rightarrow E_{\psi(u)}$ giving rise to a unique map
$q_u:D_{v_0}\rightarrow E_{\psi(u)}$. If there is
only one vertex $u_1'$ of $\tau'$ such that $\chi(u_1')=u$, then we are in a
situation as in Diagram~(\ref{fdgc}), and
by Lemma~\ref{lcsn} $q_u$ has to map $C_{v_2}\subset D_{v_0}$ to a single point
of $E_{\psi(u)}$ and
$r_{u_1'}:C_{v_1}\rightarrow E_{\psi(u)}$ suffices to determine
$q_u:D_{v_0}\rightarrow E_{\psi(u)}$ uniquely. This
defines $q:D\rightarrow E$ satisfying all properties required of a morphism of
stable maps, as some routine
considerations reveal. This finishes Case~II\@. \qed

{\em Case III\@. } This case is similar to Case~II, but much simpler, because
the construction of the composition of
$\phi$ and $(b,\sigma,\psi)$ is simpler, and thus for every $i\in F_{\sigma'}$
we have $c\chi^F(i)=\phi^Fb(i)$. We use
Proposition~\ref{glue1} instead of Proposition~\ref{glue} to construct the
pushforward of $(C,x,f)$ under $\phi$, gluing
the two sections $x_{i_1}$ and $x_{i_2}$ of $C_{v_1}=C_{v_2}$ to obtain
$D_{v_0}$. \qed

{\em Case IV\@. } To construct the pushforward, set $D_v=C_{a(v)}$,
$p_v:C_{a(v)}\rightarrow D_v$ the identity and
$h_v=f_{a(v)}$, for every $v\in V\s$. Moreover, for $i\in F\s$ set
$y_i=x_{a(i)}$. To check that $(D,y,h)$ is a stable
map and $p:(C,x,f)\rightarrow(D,y,h)$ a morphism of stable maps, the only fact
to check is that for every $i\in F\s$ we
have $h_{\del(i)}(y_i)=h_{\del(j(i))}(y_{j(i)})$, in other words
\[f_{\del a(i)}(x_{a(i)})=f_{\del(j(i))}(x_{a(j(i))}).\]
Here, Condition~(\ref{commor3}) in the definition of combinatorial morphism of
$A$-graphs (Definition~\ref{commor})
enters in. It implies this claim together with Corollary~\ref{zc}. The
universal mapping property of $(D,y,h)$ is easily
verified. \qed

{\em Case V\@. }  Before we can treat this case, we need a few preparations.

\begin{prop} \label{mplsm}
Let $(C,x_1,\ldots,x_n,f)$ be a stable map over a field from a curve of genus
$g$ to $V$ and $M$ an ample invertible
sheaf on $V$. Then
\[L=\omega_C(x_1+\ldots+x_n)\otimes f\upst M^{\otimes3}\]
is ample on $C$. Here $\omega_C$ is the dualizing sheaf of $C$.
\end{prop}
\begin{pf}
Let us first consider the case that $C$ has no nodes, so that $C$ is
irreducible and non-singular. Then is suffices to
prove that $\deg L>0$.

{\em Case 1. }The image $f(C)$ is a point. Then $\deg f\upst M=0$ and we have
\[\deg L=\deg\omega_C+n=2g-2+n\geq1,\]
by the stability condition.

{\em Case 2. }The image $f(C)$ is not a point. Then $\deg f\upst M\geq1$ and so
\[\deg L=2g-2+n+3\deg f\upst M\geq 2g-2+n+3=2g+n+1>0.\]

So suppose now that $C$ has a node $P$. Let $q:C'\rightarrow C$ be the curve
obtained by blowing up $P$ and let
$P_1,P_2\in C'$ be the two points lying over $P$. Let $L'=q\upst L$ and
$f'=f\comp q$.

{\em Case 1. }The curve $C'$ is connected. Then
$(C',x_1,\ldots,x_n,P_1,P_2,f')$ is a stable map and
\[L'=\omega_{C'}(x_1+\ldots+x_n+P_1+P_2)\otimes {f'}\upst M^{\otimes3}.\]

{\em Case 2. }The curve $C'$ is disconnected. Let $C_1'$ and $C_2'$ be the two
components of $C'$ and $L_1',L_2'$ the
restriction of $L'$ to $C_1'$ and $C_2'$, respectively. Let
$f_i':C_i'\rightarrow V$ for $i=1,2$ be the map induced by
$f'$. Without loss of generality assume that $x_1,\ldots,x_r\in  C_1'$ and
$x_{r+1},\ldots,x_n\in C_2'$, for some $0\leq
r\leq n$ and $P_1\in C_1'$, $P_2\in C_2'$. Then
$(C_1',x_1,\ldots,x_r,P_1,f_1')$ and $(C_2',x_{r+1},\ldots,x_n,P_2,f_2')$
are stable maps and
\begin{eqnarray*}
L_1' & = & \omega_{C_1'}(x_1+\ldots +x_r+P_1)\otimes {f_1'}\upst M^{\otimes3}\\
L_2' & = & \omega_{C_2'}(x_{r+1}+\ldots+x_n+P_2)\otimes {f_2'}\upst
M^{\otimes3}.
\end{eqnarray*}

Thus by induction on the the number of nodes we may assume that $L'$ is ample
on $C'$. Let $\fF$ be a coherent sheaf on
$C$ and $\fF'=q\upst\fF$. Then there exists an $n_0$ such that for all $n\geq
n_0$ we have that $\fF'\otimes{L'}^{\otimes
n}(-P_1)$, $\fF'\otimes{L'}^{\otimes n}(-P_2)$ and $\fF'\otimes{L'}^{\otimes
n}(-P_1-P_2)$ are generated by global
sections. This implies that $\fF\otimes L^{\otimes n}$ is generated by global
sections. So $L$ is ample.
\end{pf}

We will now consider the following setup. Let $(C,x_1,\ldots,x_{n+1},f)$ be a
stable map over $T$ from an
$(n+1)$-pointed curve of genus $g$ to $V$ of class $\beta\in H_2(V)^+$, where
$2g+n\geq3$ if $\beta=0$ (otherwise,
$n\geq0$).

If $t$ is a geometric point of $T$ and $C'$ a component of $C_t$, then we say
that $C'$ {\em is to be contracted}, if,
after removing $x_{n+1}$, the normalization of $C'$ violates the stability
condition. Equivalently,
\begin{enumerate}
\item $C'$ is rational without self intersection (so that $C'$ is equal to its
normalization),
\item $x_{n+1}\in C'$,
\item $C'$ has exactly two special points besides $x_{n+1}$, at least one of
which is not a marked point, but a node,
\item $f_t(C')$ is a single point of $V$.
\end{enumerate}
Pictorially, the only two possible components to be contracted look as follows.
\[\beginpicture

\setcoordinatesystem units <1cm,.71cm> point at 4 0
\setplotarea x from -1 to 3, y from -3 to 1
\plot 0 1 0 -3 /
\plot -1 0 3 0 /
\put {\circle*{8}} [Bl] at 0 -3
\put {$\boldsymbol{\times}$} at 1 0
\put {$\boldsymbol{\times}$} at 2 0
\put {$x_{n+1}$} at 2.2 -.4

\setcoordinatesystem units <1cm,.71cm> point at -3 0
\setplotarea x from -2 to 2, y from -3 to 1
\plot -2 0 2 0 /
\setquadratic
\plot 0 -3 -.75 -1 -1 1 /
\plot 0 -3  .75 -1  1 1 /
\put {\circle*{8}} [Bl] at 0 -3
\put {$\boldsymbol{\times}$} at 0 0
\put {$x_{n+1}$} at 0.2 -.4

\endpicture\]
Note that every geometric fiber of $\pi:C\rightarrow T$ has at most one
component to be contracted.

We say a $T$-morphism $q:C\rightarrow U$, for a $T$-scheme $U$, {\em contracts
the components to be contracted}, if for
every geometric point $t$ of $T$ the map (of underlying Zariski topological
spaces) $q_t:C_t\rightarrow U_t$ maps every
component to be contracted to a single point in $U_t$. For example,
$f:C\rightarrow V_T$ contracts the components to be
contracted.

\begin{prop} \label{cae}
There exists a universal morphism $p:C\rightarrow\tilde{C}$ contracting the
components to be contracted. Let
$\tilde{f}:\tilde{C}\rightarrow V$ be the unique map given by the universal
mapping property of $(\tilde{C},p)$. Then
$(\tilde{C},p(x_1),\ldots, p(x_n),\tilde{f})$ is a stable map from an
$n$-pointed curve of genus $g$ to $V$ of class
$\beta$.
\end{prop}
\begin{pf}
This is a variation on Section~1 of \cite{knudsen}.  Let us first prove the
proposition for the case that
$T$ is the spectrum of an algebraically closed field. Let $C'$ be a component
to be contracted.

{\em Case 1. }The component $C'$ has one node. We define $\tilde{C}=C-(C'-C)$
and let $p:C\rightarrow \tilde{C}$ be the
obvious map. Clearly, $\O_{\tilde{C}}=p\lst\O_C$, so $\tilde{C}$ satisfies the
universal mapping property by
Lemma~\ref{tloq}. The rest is trivial.

{\em Case 2. }The component $C'$ has two nodes. We define $\ol{C}=C-(C'-C)$ and
let $P_1$ and $P_2$ be the two points of
$\ol{C}$ that intersect $C'$. Then we set $\tilde{C}=\ol{C}/P_1\sim P_2$, i.e.\
we identify the two points $P_1$ and
$P_2$. We then proceed similarly as in Case~1.

\begin{lem} \label{clocct}
Let $T$ be the spectrum of an algebraically closed field and let $\tilde{C}$ be
the universal curve contracting the
components of $C$ to be contracted.  Choose an ample invertible sheaf $M$ on
$V$. Let
\[L=\omega_C(x_1+\ldots+x_n)\otimes f\upst M^{\otimes3}\]
and
\[\tilde{L}=\omega_{\tilde{C}}(p(x_1)+\ldots+p(x_n))\otimes \tilde{f}\upst
M^{\otimes3}.\]
Then for all $k\geq0$ we have
\begin{enumerate}
\item $\tilde{L}^{\otimes k}=p\lst L^{\otimes k}$,
\item $p\upst \tilde{L}^{\otimes k}={L}^{\otimes k}$,
\item $R^1f\lst  L^{\otimes k}=0$,
\item $H^i(\tilde{C},\tilde{L}^{\otimes k})=H^i(C,L^{\otimes k})$, for $i=0,1$.
\end{enumerate}
\end{lem}
\begin{pf}
This is analogous to Lemma~1.6 of \cite{knudsen}.
\end{pf}

\begin{lem} \label{tfll}
Let $T$ be the spectrum of an algebraically closed field and let $L$ be defined
as in Lemma~\ref{clocct}. Define the
open subset $U$ of $C$ by
\[U=\{x\in C\st \rtext{$x$ is smooth and $x$ is not in any component to be
contracted}\}.\]
Then for $k$ sufficiently large we have
\begin{enumerate}
\item $L^{\otimes k}$ is generated by global sections, \label{genglob}
\item $H^1(C,L^{\otimes k})=0$,
\item $L^{\otimes k}$ is normally generated, \label{normgen}
\item $L^{\otimes k}(-P)$ is generated by global sections for all $P\in U$.
\label{mpglob}
\end{enumerate}
(The sheaf $L$ is normally generated if $\Gamma(C,L)^{\otimes
k}\rightarrow\Gamma(C,L^{\otimes k})$ is surjective, for
all $k\geq1$.)
\end{lem}
\begin{pf}
Let $\tilde{C}$ and $\tilde{L}$ be as in Lemma~\ref{clocct}. Note that one can
apply Proposition~\ref{mplsm} to
$\tilde{C}$ and $\tilde{L}$. Then the results are implied by
Lemma~\ref{clocct}.
\end{pf}

We can now proceed with the proof of Proposition~\ref{cae} for general base
$T$. Choose an ample invertible sheaf $M$ on
$V$ and consider on $C$ the invertible sheaf
\[L=\omega_{C/T}(x_1+\ldots+x_n)\otimes f\upst M^{\otimes3},\]
where $\omega_{C/T}$ is the relative dualizing sheaf of $C$ over $T$.  Then
form
\[\sS=\bigoplus_{k\geq0}\pi\lst L^{\otimes k},\]
where $\pi:C\rightarrow T$ is the structure map, and let
\[\tilde{C}=\proj\sS.\]

{\em Claim 1. }The formation of $\tilde{C}$ commutes with base change.
\begin{pf}
Clearly, the formation of $L^{\otimes k}$ commutes with base change.  That the
formation of $\pi\lst L^{\otimes k}$
commutes with base change for $k$ sufficiently large follows from the fact that
$H^1(C_t,L^{\otimes k})=0$, for all
$t\in T$, by Lemma~\ref{tfll}. For $k=0$ this is trivially true. Thus the
formation of
\[\sS^{(d)}=\bigoplus_{d\mid k}\sS_k\]
commutes with base change, for a suitable $d>0$. This implies the claim for
$\tilde{C}$, since
\[\tilde{C}=\proj\sS=\proj\sS^{(d)}.\qed\]
\renewcommand{\qed}{}\end{pf}

{\em Claim 2. }The structure map $\tilde{\pi}:\tilde{C}\rightarrow T$ is flat
and projective.
\begin{pf}
The flatness of $\proj\sS^{(d)}$ follows from the fact that $\pi\lst L^{\otimes
k}$ is locally free, for $d\mid k$,
which follows from the fact that its formation commutes with base change. By
passing to a larger $d$ if necessary, we
may assume that for every $k\geq0$ the homomorphism
\[\pi\lst(L^{\otimes d})^{\otimes k}\longrightarrow\pi\lst(L^{\otimes dk}) \]
is surjective. This may be checked on fibers and thus follows from
Lemma~\ref{tfll}(\ref{normgen}). So $\sS^{(d)}$ is
generated by $\sS^{(d)}_1$ and hence $\proj\sS^{(d)}$ is projective by
Proposition~5.5.1 in \cite{ega2}.
\end{pf}

{\em Claim 3. }The canonical morphism from an open subset of $C$ to $\tilde{C}$
is everywhere defined, proper and
surjective.
\begin{pf}
This canonical morphism is defined by $\pi\upst\sS\rightarrow\bigoplus_k
L^{\otimes k}$, or equivalently by
$\sS\rightarrow\bigoplus_k\pi\lst L^{\otimes k}$ (see Section~3.7 in
\cite{ega2}). For it to be everywhere defined, it
suffices to prove that $\pi\upst\pi\lst L^{\otimes k}\rightarrow L^{\otimes k}$
is an epimorphism, for $k$ sufficiently
large. This may be checked on fibers and thus follows from
Lemma~\ref{tfll}(\ref{genglob}). That the canonical morphism
is dominant follows immediately, since $\sS\rightarrow\bigoplus\pi\lst
L^{\otimes k}$ is an isomorphism. That it is
proper, is clear. So it has to be surjective.
\end{pf}

We call this canonical morphism $p:C\rightarrow\tilde{C}$.

{\em Claim 4. }Let $x$ be a geometric point of $\tilde{C}$ and $p^{-1}(x)$ the
fiber of $p$ over $x$. Then either the
cardinality of $p^{-1}(x)$ is one or $p^{-1}(x)$ is a component of
$C_{\tilde{\pi}(x)}$ to be contracted.
\begin{pf}
Without loss of generality we may assume that $T$ is the spectrum of an
algebraically closed field. Then with the
notation of Lemma~\ref{tfll} and by Property~(\ref{mpglob}) of the same lemma,
we have that $p\resto
U:U\rightarrow\tilde{C}$ is an open immersion. If $C'$ is to be contracted,
then $L\resto C'\cong\O_{C'}$, and so $C'$
is mapped to a point in $\tilde{C}$. These facts clearly imply Claim~4.
\end{pf}

{\em Claim 5. }We have $p\lst\O_C=\O_{\tilde{C}}$.
\begin{pf}
With the notation of Claim~4, note that
\[H^1(p^{-1}(x),\O_C\otimes_{\O_{\tilde{C}}}\kappa(x))=0,\]
since $p^{-1}(x)$ is rational if it is one and not zero dimensional. So by
Corollary~1.5 in \cite{knudsen}, the
formation of $p\lst\O_C$ commutes with base change in $T$. So to prove the
claim, we may assume that $T$ is the
spectrum of an algebraically closed field, but then it is clear.
\end{pf}

Now by Lemma~\ref{tloq} the last three claims imply that
$p:C\rightarrow\tilde{C}$ is a universal morphism contracting
the components to be contracted. In particular, we get a unique morphism
$\tilde{f}:C\rightarrow V$ such that
$\tilde{f}\comp p=f$. The fact that
$(\tilde{C},p(x_1),\ldots,p(x_n),\tilde{f})$ is a stable map from an
$n$-pointed
curve of genus $g$ to $V$ of class $\beta$ may now be checked on fibers, which
has already been done.  This finishes the
proof of Proposition~\ref{cae}.
\end{pf}

We now proceed with the proof of Theorem~\ref{mbfc} in Case~V\@. Let $n=\#
F\s(v_0)$. Choose an identification
$\ul{n+1}\rightarrow F\t(v_0)$ mapping $n+1$ to $i_0$, the flag being removed.
Then
$(C_{v_0},x_1,\ldots,x_{n+1},f_{v_0})$ is a stable map to which
Proposition~\ref{cae} applies and we let
$p_{v_0}:C_{v_0}\rightarrow D_{v_0}$ be the universal morphism contracting the
components to be contracted. For
$v\not=v_0$ we let $D_v=C_v$ and $p_v:C_v\rightarrow D_v$ be the identity. It
is then clear how to define $y$ and $h$ to
get a stable map $(D,y,h)$ satisfying the universal mapping property of a
pushforward under the graph morphism
$\tau\rightarrow \sigma$ given by $a:\sigma\rightarrow \tau$. This completes
the proof of Case~V\@. \qed

To complete the proof of Theorem~\ref{mbfc} we need to show that if
$(V,\tau,\alpha)$ is an object of $\VV\GG_s$ and
$p:(C,x,f)\rightarrow(D,y,h)$ is a morphism of stable $(V,\tau,\alpha)$-maps
(covering the identity of
$(V,\tau,\alpha)$), then $p$ is an isomorphism.

This is immediately reduced to the case that $(V,\tau,\alpha)=(V,g,n,\alpha)$
and using Lemma~\ref{tloq} to the case that
$T$ is the spectrum of an algebraically closed field. Then it follows from the
stability condition that $p$ cannot
contract any rational components, so it is injective. To prove that $p$ is
surjective use induction on the number of
nodes of $D$. So let $P$ be a node of $D$ and let $D'$ be the curve obtained
from $D$ by blowing up $P$ and let
$p':C'\rightarrow D'$ by the pullback of $p:C\rightarrow D$ under
$D'\rightarrow D$.

{\em Case 1. } The curve $D'$ is disconnected, $D'=D_1'\amalg D_2'$. Then
$C'=C_1'\amalg C_2'$ with induced maps
$p_i':C_i'\rightarrow D_i'$, for $i=1,2$. Let $g_i=g(D_i')$ and
$\alpha_i=f\lst[D_i']$, for $i=1,2$. Then $g=g_1+g_2$
and $\alpha=\alpha_1+\alpha_2$. Now $g_i(C_i')\leq g_i(D_i')$ and
$f\lst[C_i']\leq f\lst[D_i']$ imply that
$g_i(C_i')=g_i$ and $f\lst[C_i']=\alpha_i$ and thus we may apply the induction
hypothesis to $p_1'$ and $p_2'$, proving
the surjectivity of $p$.

{\em Case 2. } The curve $D'$ is connected. Then $C'$ is connected, since
otherwise we would have two
curves contradicting the induction hypothesis. So me may apply the induction
hypothesis to $p':C'\rightarrow D'$.

This finally completes the proof of Theorem~\ref{mbfc}. \qed

\begin{defn} \label{doosf}
For a given object $(V,\tau)$ of $\VV\GG_s$, we let $\ol{M}(V,\tau)(T)$ be the
fiber of $\ol{M}(T)$ over
$(V,\tau)$ under the cofibration of Theorem~\ref{mbfc}.

Letting $T$ vary we get a stack $\ol{M}(V,\tau)$ on the category of
$k$-schemes with the fppf-topology.

For $(V,\tau)=(V,g,n,\beta)$ we denote $\ol{M}(V,\tau)$ by
$\ol{M}_{g,n}(V,\beta)$.
\end{defn}

If $\chr k\not=0$, let $L$ be a very ample invertible sheaf on $V$. Consider
only stable $V$-graphs for which
$\beta(v)(L)<\chr k$, for all $v\in V\t$. If this condition is satisfied, we
say that $\tau$ or $(V,\tau)$ is {\em
bounded by the characteristic}. If $\chr k=0$, we call {\em every }$(V,\tau)$
bounded by the characteristic, so that we
have uniform terminology.

\begin{them} \label{mgnvas}
For every $(V,\tau)$, bounded by the characteristic, the stack $\ol{M}(V,\tau)$
is a proper algebraic
Deligne-Mumford stack over $k$.
\end{them}
\begin{pf}
The proof will be postponed to a later section (see Corollary~\ref{mgnvass}).
\end{pf}

Every time we refer to $\ol{M}(V,\tau)$ as a Deligne-Mumford stack, we shall
tacitly assume that $(V,\tau)$ is bounded
by the characteristic.

\begin{rmk}
Theorems~\ref{mbfc} and~\ref{mgnvas} give rise to a functor
\begin{eqnarray*}
\ol{M}:\VV\GG_s & \longrightarrow & (\rtext{\normalshape proper algebraic
DM-stacks over $k$})\\
(V,\tau) & \longmapsto & \ol{M}(V,\tau),
\end{eqnarray*}
by choosing for every $k$-scheme $T$ a {\em clivage normalis\'e }(see
Definition~7.1 in \cite[Exp.\ ~VI]{sga1}) of the
cofibered category $\ol{M}(T)$ over $\VV\GG_s$. Of course, this functor
$\ol{M}$ is essentially independent of the
choice of the {\em clivage normalis\'e}.

Another way of stating this would be to construct a fibered category $\ol{M}$
over
$\VV\GG^{\op}_s\times(\text{$k$-schemes})$, such that $\ol{M}(V,\tau)(T)$ is
the fiber of $\ol{M}$ over
$(V,\tau,T)$ and $\ol{M}(T)$ is the fiber of $\ol{M}$ over $T$.
\end{rmk}

\section{Further Study of $\ol{M}$}

\begin{prop} \label{drfu}
Let $(V,\tau)$ be an object of $\VV\GG_s$, bounded by the characteristic of
$k$. Then the diagonal
\[\Delta:\ol{M}(V,\tau)\longrightarrow\ol{M}(V,\tau)\times\ol{M}(V,\tau)\]
is representable, finite and unramified.
\end{prop}
\begin{pf}
The assumption that $(V,\beta)$ is bounded by the characteristic implies that
all stable maps of class $\beta$ are
separable. So by Lemma~\ref{rsmsc} we may reduce the case of stable maps to the
case of stable curves, which is
well-known.
\end{pf}

\begin{lem} \label{rsmsc}
Let $(C,x,f)$ and $(D,y,h)$ be $n$-pointed stable maps to $V$ over the base
$T$, and $t\in T(K)$ a geometric point of
$T$. Assume that $f_t:C_t\rightarrow V$ and $h_t:D_t\rightarrow V$ are
separable morphisms. Then there exists an \'etale
neighborhood $S\rightarrow T$ of $t$, an integer $N$, markings
$x'=(x_1',\ldots,x_N')$ of $C_S$ and
$y'=(y_1',\ldots,y_N')$ of $D_S$ such that $(C_S,x_S,x')$ and $(D_S,y_S,y')$
are stable marked curves over $S$ and a
closed immersion of sheaves on $(\rtext{$S$-schemes})$
\[\Isomu((C,x,f),(D,y,h))_S\longrightarrow\Isomu((C_S,x_S,x'),(D_S,y_S,y')).\]
\end{lem}
\begin{pf}
Without loss of generality assume that $C$ and $D$ have the same genus $g$ and
$f$ and $h$ have the same class $\beta$.
Choose an embedding $\mu:V\hookrightarrow\pp^r$, let $d=\mu\lst\beta$ and
reduce to the case $V=\pp^r$ and $d=\beta$.
Let $N=d(r+1)$. Choose linearly independent hyperplanes $H_0,\ldots,H_r$ in
$\pp^r$ such that for each $i=0,\ldots,r$
\begin{enumerate}
\item no special point of $C_t$ or $D_t$ is mapped into $H_{i,K}$ under $f_t$
or $g_t$,
\item $f_t$ and $g_t$ are transversal to $H_{i,K}$.
\end{enumerate}
Then there exists an \'etale neighborhood $S\rightarrow T$ of $t$ such that
\begin{enumerate}
\item for each $i=0,\ldots,r$
\begin{enumerate}
\item $H_{i,S}\cap C_S$ gives rise to $d$ sections $x_{di+1}',\ldots,x_{di+d}'$
of $C_S$ over $S$,
\item $H_{i,S}\cap D_S$ gives rise to $d$ sections $y_{di+1}',\ldots,y_{di+d}'$
of $D_S$ over $S$,
\end{enumerate}
\item $(C_S,x_S,x')$ and $(D_S,y_S,y')$ are marked prestable curves.
\end{enumerate}
Then  $(C_S,x_S,x')$ and $(D_S,y_S,y')$ are in fact stable and there exists an
obvious morphism
\[\Isomu((C,x,f),(D,y,h))_S\longrightarrow\Isomu((C_S,x_S,x'),(D_S,y_S,y')),\]
which is clearly a closed immersion.
\end{pf}

\begin{lem} \label{lalss}
Let $(C,x_1,\ldots,x_{n+1},f)$ be a stable map and $(D,y_1,\ldots,y_n,h)$ the
stabilization under forgetting $x_{n+1}$.
Let $p:C\rightarrow D$ be the structure morphism. Then any section $y_0$ of $D$
making $(D,y_0,\ldots,y_n)$ a marked
prestable curve lifts uniquely to a section $x_0$ of $C$ making
$(C,x_0,\ldots,x_n)$ a marked prestable curve. If $y_0$
avoids $p(x_{n+1})$, then $(C,x_0,\ldots,x_{n+1})$ is a marked prestable curve.
\end{lem}
\begin{pf}
Let $V\subset D$ be the open subset consisting of smooth points of $D$ which
are not in the image of $y_i$, for any
$i=1,\ldots,n$. Let $U=p^{-1}(V)$. Then $p$ induces an isomorphism $p\resto
U:U\iso V$. Moreover, $U$ is smooth and
$x_{n+1}$ is the only section of $C$ which may meet $U$.
\end{pf}

\begin{prop} \label{uors}
Let $(C,x_1,\ldots,x_{n+1},f)$ and
$(\tilde{C},\tilde{x}_1,\ldots,\tilde{x}_{n+1},\tilde{f})$ be stable maps with
isomorphic stabilizations forgetting the $(n+1)$-st section. Let
$(C,y_1,\ldots,y_n,h)$ be such a stabilization, with
structure maps $p:C\rightarrow D$ and $\tilde{p}:\tilde{C}\rightarrow D$. If
$p(x_{n+1})=\tilde{p}(\tilde{x}_{n+1})$
then there exists a unique isomorphism $q:C\rightarrow\tilde{C}$ of stable maps
such that $\tilde{p}\comp q=p$.
\end{prop}
\begin{pf}
This is local over the base, so we may freely choose sections as necessary. In
fact, choose sections $z_1,\ldots,z_N$ of
$D$ in the smooth locus, avoiding $y_1,\ldots,y_n$ and
$\Delta=p(x_{n+1})=\tilde{p}(\tilde{x}_{n+1})$ and making
$$(D,z_1,\ldots,z_N,y_1,\ldots,y_n)$$ a stable marked curve. By
Lemma~\ref{lalss} these lift uniquely to sections
$w_1,\ldots,w_N$ of $C$ and $\tilde{w}_1,\ldots,\tilde{w}_N$ of $\tilde{C}$
making
$$(C,w_1,\ldots,w_N,x_1,\ldots,x_{n+1})$$ and
%% FOLLOWING LINE CANNOT BE BROKEN BEFORE 80 CHAR
$$(\tilde{C},\tilde{w}_1,\ldots,\tilde{w}_N,\tilde{x}_1,
\ldots,\tilde{x}_{n+1})$$ marked prestable curves. Moreover,
these are clearly marked {\em stable } curves with a common stabilization
$$(D,z_1,\ldots,z_N,y_1,\ldots,y_n)$$ forgetting
the last section, such that $p(x_{n_1})=\tilde{p}(\tilde{x}_{n+1})$. Then they
have to be isomorphic by Knutson's theorem
(see~\cite{knudsen}) that $\ol{M}_{g,N+n+1}$ is the universal curve over
$\ol{M}_{N+n}$.
\end{pf}

\begin{prop} \label{flpuc}
Let $(C,x_1,\ldots,x_n,f)$ be a stable map and $\Delta$ a section of $C$. Then
there exists up to isomorphism a unique
stable map $(\tilde{C},\tilde{x}_1,\ldots,\tilde{x}_{n+1},\tilde{f})$ such that
$(C,x_1,\ldots,x_n,f)$ is the
stabilization of $(\tilde{C},\tilde{x}_1,\ldots,\tilde{x}_{n+1},\tilde{f})$
forgetting the $(n+1)$-st section and
$p(\tilde{x}_{n+1})=\Delta$, where $p:\tilde{C}\rightarrow C$ is the structure
map.
\end{prop}
\begin{pf}
Uniqueness follows from Proposition~\ref{uors}, hence existence is a local
question. Thus we may choose sections
$z_1,\ldots,z_N$ of $C$ such that $$(C,z_1,\ldots,z_N,x_1,\ldots,x_n)$$ is a
stable marked curve. By Knudsen's result
again, there exists a stable curve
$$(C',z_1',\ldots,z_N',x_1',\ldots,x_{n+1}')$$ whose stabilization forgetting
the last
section is $$(C,z_1,\ldots,z_N,x_1,\ldots,x_n)$$ and such that
$q(x_{n+1}')=\Delta$, where $q:C'\rightarrow C$ is the
structure map. Clearly, $$(C',z_1',\ldots,z_N',x_1',\ldots,x_{n+1}',f\comp q)$$
is a stable map. Then let
$(\tilde{C},\tilde{x}_1,\ldots,\tilde{x}_{n+1},\tilde{f})$ be the stabilization
of
$$(C',z_1',\ldots,z_N',x_1',\ldots,x_{n+1}',f\comp q)$$ forgetting the sections
$z_1',\ldots,z_N'$. By its universal
mapping property there exists a morphism $p:\tilde{C}\rightarrow C$ which makes
$(C,x_1,\ldots,x_n,f)$ the stabilization of
$(\tilde{C},\tilde{x}_1,\ldots,\tilde{x}_{n+1},\tilde{f})$ forgetting
$\tilde{x}_{n+1}$.
\end{pf}

\begin{cor} \label{esuc}
Let $C_{g,n}(V,\beta)$ be the universal curve over $\ol{M}_{g,n}(V,\beta)$.
Then the canonical morphism
$\ol{M}_{g,n+1}(V,\beta)\rightarrow C_{g,n}(V,\beta)$ induced by the $(n+1)$-st
section is an isomorphism. \qed
\end{cor}

\begin{prop}
Let $(C,x,f)$ be a stable $(V,g,n,\beta)$-map over $T$. Then the set of $t\in
T$ such that $(C,x)$ is a stable marked
curve is open in $T$.
\end{prop}
\begin{pf}
The set of such $t$ is the set of all $t\in T$ for which $(C,x)$ is isomorphic
to its stabilization. For any morphism of
schemes, the set of elements of its source at which it is an isomorphism is
always open. Finally, use properness of
prestable curves.
\end{pf}

By this proposition we may define
\[U_{g,n}(V,\beta)\subset\ol{M}_{g,n}(V,\beta)\]
to be the open substack of those stable maps, whose underlying marked curve is
stable. The canonical morphism
$U_{g,n}(V,\beta)\rightarrow \ol{M}_{g,n}$ has as fiber over the marked curve
$(C,x)$ the scheme of morphisms form $C$
to $V$ of class $\beta$. By results of Grothendieck in \cite{fgaIV} this is a
quasi-projective scheme. Hence
$U_{g,n}(V,\beta)$ is an algebraic $k$-stack of finite type. Now, for given $n$
there exists an $N>n$ such that
$U_{g,N}(V,\beta)\rightarrow\ol{M}_{g,n}(V,\beta)$  is surjective. Since this
morphism is flat by Corollary~\ref{esuc},
it is a flat epimorphism, hence a presentation of $\ol{M}_{g,n}(V,\beta)$.
Together with Proposition~\ref{drfu} this
implies that $\ol{M}_{g,n}(V,\beta)$ is a finite type separated algebraic
Deligne-Mumford stack over $k$. This is then
true for all objects of $\VV\GG_s$, bounded by the characteristic.

\begin{cor} \label{mgnvass}
Theorem~\ref{mgnvas} is true.
\end{cor}
\begin{pf}
It only remains to show properness. This is easily reduced to the case
$(V,\tau)=(\pp^r,g,n,d)$ and follows from
Proposition~3.3 of \cite{pandh}.
\end{pf}

\section{An Operadic Picture}

\begin{defn}
Let $(\tau,\alpha)$ be an $A$-graph. Let $R\t\subset F\t\times F\t$ be defined
by $(f,\ol{f})\in R\t$ if and only if one
of the conditions
\begin{enumerate}
\item $\ol{f}=j\t(f)$,
\item $\del f=\del\ol{f}$ and for $v=\del f=\del\ol{f}$ we have
$g(v)=\alpha(v)=0$
\end{enumerate}
is satisfied. Let $\sim$ be the equivalence relation on $F\t$ generated by
$R\t$ and let
\[P\t=F\t/\sim.\]
(In fact, $P_{(\tau,\alpha)}$ would be better notation, but we will stick with
the abuse of notation $P_{\tau}$.)
\end{defn}

\begin{prop} \label{cmmgpe}
Let $a:(B,\sigma)\rightarrow(A,\tau)$ be a combinatorial morphism of marked
graphs. Then $a_F:F\s\rightarrow F\t$
preserves equivalence.  \qed
\end{prop}

\begin{rmk}
In fact, Condition~(\ref{commor3}) of Definition~\ref{commor} may be replaced
by requiring $a_F$ to preserve
equivalence.
\end{rmk}

\begin{prop} \label{dafpea}
Let $\phi:\tau\rightarrow\sigma$ be a contraction of $A$-graphs. Then
$\phi^F:F\s\rightarrow F\t$ preserves equivalence.
\qed
\end{prop}

\begin{prop} \label{dpc}
If
\[\begin{array}{ccccc}
B & & \pi & \stackrel{\psi}{\longrightarrow} & \rho \\
\ldiagup{\xi} & \phantom{\longrightarrow} & \ldiag{b} &  & \rdiag{a} \\
A & & \sigma & \stackrel{\phi}{\longrightarrow} & \tau
\end{array}\]
is a stable pullback, then the induced diagram
\[\comdiaback{P_\pi}{\psi^F}{P_\rho}{b}{}{a}{P\s}{\phi^F}{P\t}\]
commutes. \qed
\end{prop}

By Propositions~\ref{cmmgpe},~\ref{dafpea} and~\ref{dpc}, we have a
contravariant functor
\[P:\GG_s\longrightarrow(\rtext{finite sets})\]
given by $P(A,\tau)=P\t$ on objects. Composing with the functor
$\VV\GG_s\rightarrow\GG_s$ we get a contravariant
functor
\begin{eqnarray*}
P:\VV\GG_s & \longrightarrow & (\rtext{finite sets}) \\
(V,\tau)   & \longmapsto     & P\t.
\end{eqnarray*}
There is an obvious functor
\begin{eqnarray*}
\VV\times(\rtext{finite sets}) & \longrightarrow & \VV \\
(V,P)                          & \longmapsto     & V^P,
\end{eqnarray*}
contravariant in the second argument, and composing with $P$ times the natural
functor $\VV\GG_s\rightarrow\VV$ gives
rise to a covariant functor
\begin{eqnarray*}
P:\VV\GG_s & \longrightarrow & \VV \\
(V,\tau)   & \longmapsto     & V^{P\t},
\end{eqnarray*}
still denoted $P$, by abuse of notation. We may consider $\VV$ as a subcategory
of the 2-category of proper algebraic
Deligne-Mumford stacks over $k$ and consider this as a functor
\[P:\VV\GG_s\longrightarrow(\rtext{proper algebraic DM-stacks over $k$}).\]

Now fix an object $(V,\tau)$ of $\VV\GG_s$. Let $(C,x,f)$ be a stable
$(V,\tau)$-map over $T$. Then $x$ and $f$ define a
morphism
\begin{eqnarray*}
f(x):T & \longrightarrow & V^{F\t} \\
t      & \longmapsto     & (f(x_i(t)))_{i\in F\t}.
\end{eqnarray*}
By Corollary~\ref{zc} this morphism $f(x)$ factors through $V^{P\t}\subset
V^{F\t}$, so we consider it as a morphism
\[f(x):T\longrightarrow V^{P\t}.\]
Thus we get a map $\ol{M}(V,\tau)(T)\rightarrow P(V,\tau)(T)$. Since it is
compatible with base change $S\rightarrow
T$, we have a morphism of $k$-stacks
\[\ev(V,\tau):\ol{M}(V,\tau)\longrightarrow P(V,\tau).\]

\begin{prop}
We have defined a natural transformation of functors from $\VV\GG_s$ to
$(\rtext{proper algebraic DM-stacks over $k$})$
\[\ev:\ol{M}\longrightarrow P,\]
called {\em evaluation}.
\end{prop}

In the general framework of $\Gamma$-operads, this allows us to consider
(appropriate subfunctors of) $\ol{M}$ and $P$
as a modular operad and a cyclic endomorphism operad, respectively. The
evaluation functor then induces a structure of
$\ol{M}$-algebra on $V$.

\vfill\eject
\part{Gromov-Witten Invariants}

\section{Isogenies}

\begin{defn}
Let $\tau$ be a stable $A$-graph.
\begin{enumerate}
\item The {\em class }of $\tau$ is
\[\beta(\tau)=\sum_{v\in V\t}\beta(v).\]
\item The {\em Euler characteristic } of $\tau$ is
\[\chi(\tau)=\chi(|\tau|)-\sum_{v\in v\t}g(v).\]
\item If $|\tau|$ is non-empty and connected the {\em genus }of $\tau$ is
\[g(\tau)=1-\chi(\tau).\]
\end{enumerate}
\end{defn}

\begin{defn}
Let $\tau$ be a stable $V$-graph, where $V$ is of pure dimension.
\begin{enumerate}
\item The {\em dimension } of $\tau$ is
\[\dim(V,\tau)=\chi(\tau)(\dim V-3)-\beta(\tau)(\omega_V)+\# S\t-\# E\t,\]
where $\omega_V$ is the canonical line bundle on $V$.
\item The {\em degree } of $\tau$ is
\begin{multline*}
\deg(V,\tau)=\\
\beta(\tau)(\omega_V)+(\dim V-3)(\chi(\tau^s)-\chi(\tau))+(\# S_{\tau^s}-\#
S\t)-(\# E_{\tau^s}-\# E\t),
\end{multline*}
where $\tau^s$ is the absolute stabilization of $\tau$.
\end{enumerate}
\end{defn}

Note that
\[\dim(V,\tau)-\dim(\tau^s)=\chi(\tau^s)\dim V-\deg(V,\tau).\]

\begin{defn}
The stable $A$-graph with one vertex of genus and class zero and three tails
(no edges) shall be called the
{\em $A$-tripod}, or simply a {\em tripod}.
\end{defn}

\begin{defn}
Let $a:\tau'\rightarrow\tau$ be a combinatorial morphism of stable $A$-graphs.
A {\em tail map } for $a$ is a map
$m:S_{\tau'}\rightarrow S\t$. Let $(a,\tau',\phi):\tau\rightarrow\sigma$ be a
morphism of stable $A$-graphs. A {\em tail
map } for $(a,\tau',\phi)$ is a tail map for $a$.
\end{defn}

Since for a contraction $\phi:\tau'\rightarrow\sigma$ the map
$\phi^F:F\s\rightarrow F_{\tau'}$ induces a bijection
$\phi^S:S\s\rightarrow S_{\tau'}$, there is an obvious way to compose morphisms
with tail map of stable $A$-graphs. By
abuse we will also call the composition $m\comp\phi^S$ the tail map of
$(a,\tau',\phi,m)$.

\begin{defn} \label{dsft}
Let $a:\tau'\rightarrow \tau$ be a combinatorial morphism of stable $A$-graphs
with tail map $m:S_{\tau'}\rightarrow
S\t$. We say that $(a,m)$ is of type {\em stably forgetting a tail}, or that
$\tau'$ is obtained form $\tau$ by {\em
stably forgetting a tail}, if there exists a tail $f$ of $\tau$ such that
\begin{enumerate}
\item $\tau'$ is the stabilization of $\tau''$, where $\tau''$ is obtained from
$\tau$ by forgetting the tail $f$,
\item for all $h\in S_{\tau'}$, such that $a(h)$ is a tail, we have
$m(h)=a(h)$,
\item $f\not\in m(S_{\tau'})$.
\end{enumerate}
\end{defn}

Note that the tail $f$ is uniquely determined by $(a,m)$ and $f$ determines
$(a,m)$ uniquely. We say that $(a,m)$ stably
forgets the tail $f$.

\begin{numrmk} \label{sft}
Every combinatorial morphism (with tail map) of type stably forgetting a tail
is of one of the following types (notation
of Definition~\ref{dsft}).

{\em Type I (Incomplete case). } No stabilization is needed, i.e.\
$\tau'=\tau''$.
\[
\beginpicture%---tau
\setcoordinatesystem units <.3cm,.3cm> point at 3 2
\setplotarea x from 2 to 6, y from 0 to 4
\plot 3 2 4 3 /
\plot 2 2 4 2 /
\plot 3 2 4 1 /
\shaderectangleson
\setshadegrid span <1mm>
\putrectangle corners at 4 0 and 6 4
\put {\circle*{4}} [Bl] at 3 2
\axis left invisible label {$\scriptstyle\tau$} /
\axis right invisible label {\phantom{$\scriptstyle\tau$}} /
\endpicture  \stackrel{a}{\longleftarrow}
\beginpicture%---tau'
\setcoordinatesystem units <.3cm,.3cm> point at 3 2
\setplotarea x from 3 to 6, y from 0 to 4
\plot 3 2 4 3 /
\plot 3 2 4 2 /
\plot 3 2 4 1 /
\shaderectangleson
\setshadegrid span <1mm>
\putrectangle corners at 4 0 and 6 4
\put {\circle*{4}} [Bl] at 3 2
\axis left invisible label {\phantom{$\scriptstyle\tau'$}} /
\axis right invisible label {{$\scriptstyle\tau'$}} /
\endpicture  \]

{\em Type II (Removing a tripod from a tail). }Only in this case does the tail
map $m$ contain any information. It
serves to `remember' which of the two tails in the following diagram is being
forgotten by $a$.
\[
\beginpicture%---tau
\setcoordinatesystem units <.3cm,.3cm> point at 3 2
\setplotarea x from 2 to 6, y from 0 to 4
\plot 2 3 3 2 /
\plot 3 2 4 2 /
\plot 2 1 3 2 /
\shaderectangleson
\setshadegrid span <1mm>
\putrectangle corners at 4 0 and 6 4
\put {\circle*{4}} [Bl] at 3 2
\axis left invisible label {$\scriptstyle\tau$} /
\axis right invisible label {\phantom{$\scriptstyle\tau$}} /
\endpicture  \stackrel{a}{\longleftarrow}
\beginpicture%---tau'
\setcoordinatesystem units <.3cm,.3cm> point at 3 2
\setplotarea x from 3 to 6, y from 0 to 4
\plot 3 2 4 2 /
\shaderectangleson
\setshadegrid span <1mm>
\putrectangle corners at 4 0 and 6 4
\axis left invisible label {\phantom{$\scriptstyle\tau'$}} /
\axis right invisible label {{$\scriptstyle\tau'$}} /
\endpicture  \]

{\em Type III (Removing a tripod from an edge). }
\[
\beginpicture%---tau
\setcoordinatesystem units <.3cm,.3cm> point at 3 2
\setplotarea x from 2 to 6, y from 0 to 4
\plot 2 2 3 2 /
\circulararc 180 degrees from 4 3 center at 4 2
\shaderectangleson
\setshadegrid span <1mm>
\putrectangle corners at 4 0 and 6 4
\put {\circle*{4}} [Bl] at 3 2
\axis left invisible label {$\scriptstyle\tau$} /
\axis right invisible label {\phantom{$\scriptstyle\tau$}} /
\endpicture  \stackrel{a}{\longleftarrow}
\beginpicture%---tau'
\setcoordinatesystem units <.3cm,.3cm> point at 3 2
\setplotarea x from 3 to 6, y from 0 to 4
\circulararc 180 degrees from 4 3 center at 4 2
\shaderectangleson
\setshadegrid span <1mm>
\putrectangle corners at 4 0 and 6 4
\axis left invisible label {\phantom{$\scriptstyle\tau'$}} /
\axis right invisible label {{$\scriptstyle\tau'$}} /
\endpicture  \]

{\em Type IV (Forgetting a lonely tripod or a lonely elliptic component.) }
Only in this case does the number of
connected components of the geometric realization change.
\[
\beginpicture%---tau
\setcoordinatesystem units <.3cm,.3cm> point at 3 2
\setplotarea x from 2 to 6, y from 0 to 4
\plot 2 3 3 2 /
\plot 2 2 3 2 /
\plot 2 1 3 2 /
\shaderectangleson
\setshadegrid span <1mm>
\putrectangle corners at 4 0 and 6 4
\put {\circle*{4}} [Bl] at 3 2
\axis left invisible label {$\scriptstyle\tau$} /
\axis right invisible label {\phantom{$\scriptstyle\tau$}} /
\endpicture  \stackrel{a}{\longleftarrow}
\beginpicture%---tau'
\setcoordinatesystem units <.3cm,.3cm> point at 3 2
\setplotarea x from 3 to 6, y from 0 to 4
\shaderectangleson
\setshadegrid span <1mm>
\putrectangle corners at 4 0 and 6 4
\axis left invisible label {\phantom{$\scriptstyle\tau'$}} /
\axis right invisible label {{$\scriptstyle\tau'$}} /
\endpicture  \]
\[
\beginpicture%---tau
\setcoordinatesystem units <.3cm,.3cm> point at 3 2
\setplotarea x from 2 to 6, y from 0 to 4
\plot 2 2 3 2 /
\shaderectangleson
\setshadegrid span <1mm>
\putrectangle corners at 4 0 and 6 4
\put {\circle*{4}} [Bl] at 3 2
\axis left invisible label {$\scriptstyle\tau$} /
\axis right invisible label {\phantom{$\scriptstyle\tau$}} /
\endpicture  \stackrel{a}{\longleftarrow}
\beginpicture%---tau'
\setcoordinatesystem units <.3cm,.3cm> point at 3 2
\setplotarea x from 3 to 6, y from 0 to 4
\shaderectangleson
\setshadegrid span <1mm>
\putrectangle corners at 4 0 and 6 4
\axis left invisible label {\phantom{$\scriptstyle\tau'$}} /
\axis right invisible label {{$\scriptstyle\tau'$}} /
\endpicture  \]
Here, the genus of the vertex displayed in the last diagram is equal to one.
\end{numrmk}

\begin{defn}
Let $(a,\tau',\phi,m):\tau\rightarrow\sigma$ be a morphism of stable $A$-graphs
with tail map. We call $(a,\tau',\phi,m)$
an {\em isogeny}, if
\begin{enumerate}
\item $(a,m)$ is a composition of morphisms of type stably forgetting a tail,
\item $\pi_0|\sigma|\rightarrow\pi_0|\tau|$ is bijective.
\end{enumerate}
We call the isogeny $\Phi:\tau\rightarrow\sigma$ an {\em elementary isogeny},
if it is an elementary contraction, or if
$\sigma$ is obtained from $\tau$ by stably forgetting a tail.
\end{defn}

\begin{note}
If $\Phi:\tau\rightarrow\sigma$ is an isogeny of stable $A$-graphs, then
$\chi(\sigma)=\chi(\tau)$.

An elementary isogeny either contracts a loop, or a non-looping edge or is of
type stably forgetting a tail~I, II,
or~III.

If we write $\Phi:\tau\rightarrow\sigma$ for an isogeny of stable $A$-graphs,
we denote the tail map of $\Phi$ by
$\Phi^S:S\s\rightarrow S\t$.
\end{note}

\begin{prop}
The composition of isogenies is an isogeny.
\end{prop}
\begin{pf}
Let
\[\begin{array}{ccccc}
A & & \pi & \stackrel{\psi}{\longrightarrow} & \rho \\
\ldiagup{\id} & \phantom{\longrightarrow} & \ldiag{b} &  & \rdiag{a} \\
A & & \sigma & \stackrel{\phi}{\longrightarrow} & \tau
\end{array}\]
be a stable pullback, where $a$ stably forgets the tail $f$ of $\tau$,
$\pi_0|\rho|\rightarrow\pi_0|\tau|$ is bijective
and $\phi$ is an elementary contraction of stable $A$-graphs. Then $b$ stably
forgets the tail $\phi^F(f)$ of $\sigma$.
Even if there is a vertex $v_0$ of $\tau$ which does not appear in $\rho$, this
vertex $v_0$ cannot be the vertex onto
which $\phi$ contracts an edge.
\end{pf}

Fix a semi-group with indecomposable zero $A$. We shall define a category
$\tilde{\GG}_s(A)$ from $\GG_s(A)$, retaining
only isogenies and morphisms of type cutting edges, but reversing the direction
of the latter, making them morphisms
{\em gluing tails}.

In fact, define the category $\tilde{\GG}_s(A)$ as follows. Objects of
$\tilde{\GG}_s(A)$ are stable $A$-graphs. A
morphism $\sigma\rightarrow\tau$ is a triple $(a,\sigma',\Phi)$, where
$a:\sigma\rightarrow\sigma'$ is a combinatorial
morphism of $A$-graphs of type cutting edges and $\Phi:\sigma'\rightarrow\tau$
is an isogeny of stable $A$-graphs. To
compose $(a,\sigma',\Phi):\sigma\rightarrow\tau$ and
$(b,\tau',\Psi):\tau\rightarrow\rho$, we need to construct a
diagram
\begin{equation}\label{dcmei}
\begin{array}{ccccc}
\sigma'' & \stackrel{\Xi}{\longrightarrow} & \tau' &
\stackrel{\Psi}{\longrightarrow} & \rho \\
\ldiagup{c} & & \rdiagup{b} &  & \\
\sigma' & \stackrel{\Phi}{\longrightarrow} & \tau & & \\
\ldiagup{a} & & & & \\
\sigma, & & & &
\end{array}
\end{equation}
where $c:\sigma'\rightarrow\sigma''$ is a combinatorial morphism of type
cutting edges and
$\Xi:\sigma''\rightarrow\tau'$ is an isogeny of stable $A$-graphs.

Let $f$ and $\ol{f}$ be two tails of $\tau$ such that $\{b(f),b(\ol{f})\}$ is
an edge of $\tau'$. Then construct
$\sigma''$ from $\sigma'$ by gluing the two tails $\Phi^S(f)$ and
$\Phi^S(\ol{f})$ to an edge. If $b$, cuts more than
one edge, iterate this process to construct $\sigma''$. This defines
composition of morphisms in $\tilde{\GG}_s(A)$,
which is clearly associative.

\begin{note}
In the situation of (\ref{dcmei}), we get a diagram in $\GG_s(A)$
\[\comdia{\sigma''}{\Xi}{\tau'}{\ol{c}}{}{\ol{b}}{\sigma'}{\Phi}{\tau},\]
which is easily seen to commute. Here, $\ol{b}$ and $\ol{c}$ are the morphisms
of stable $A$-graphs induced by $b$ and
$c$, respectively.
\end{note}

\begin{defn} \label{ecisg}
We call $\tilde{\GG}_s(A)$ the {\em extended category of isogenies of stable
$A$-graphs}, or the {\em extended isogeny
category }over $A$.

The morphisms in $\tilde{\GG}_s(A)$ are called {\em extended isogenies}. An
extended isogeny is called {\em elementary},
if it is an elementary isogeny or glues two tails to an edge.
\end{defn}

Now consider the following situation. Fix a smooth projective variety $V$ of
pure dimension. Let
$\Phi:\tau\rightarrow\sigma$ be an elementary extended isogeny of stable
modular graphs. Let $\sigma'$ be a stable
$V$-graph and $b:\sigma\rightarrow\sigma'$ a combinatorial morphism identifying
$\sigma$ as the absolute stabilization
of $\sigma'$. Note that $b$ is injective on vertices and complete, so that
$b:F\s(v)\rightarrow F_{\sigma'}(b(v))$ is
bijective, for all $v\in V\s$. Let $(a_i,\tau_i)_{i\in I}$ be a family of
pairs, where $I$ is a finite set and for each
$i\in I$ we have a combinatorial morphism $a_i:\tau\rightarrow\tau_i$
identifying $\tau$ as the absolute stabilization
of $\tau_i$. Finally, let for every $i\in I$ be given an extended isogeny of
stable $V$-graphs
$\Phi_i:\tau_i\rightarrow\sigma'$. In particular, for each $i\in I$ we have a
diagram of stable marked graphs (but note
that the horizontal and vertical morphisms live in different categories)
%% FOLLOWING LINE CANNOT BE BROKEN BEFORE 80 CHAR
\[\comdia{\tau_i}{\Phi_i}{\sigma'}{\ol{a}_i}{}{\ol{b}}{\tau}{
\Phi}{\phantom{b}\sigma \phantom{b}.}\]
We shall now define what we mean by $(a_i,\tau_i,\Phi_i)_{i\in I}$ to be {\em
cartesian}, or a {\em pullback } of
$\sigma'$ under $\Phi$.  We have to distinguish six cases, according to which
kind of elementary extended isogeny $\Phi$
is.

Let us first consider the case that $\Phi$ is an elementary contraction
$\phi:\tau\rightarrow\sigma$, contracting the
edge $\{f,\ol{f}\}$ of $\tau$. As usual, let $v_1=\del f$, $v_2=\del\ol{f}$ and
$v_0=\phi(v_1)=\phi(v_2)$. Let
$w_0=b(v_0)$.

{\em Case I (Contracting a loop). } In this case $v_1=v_2$. The set $I$ has one
element, say $0$, and
$(a_0,\tau_0,\Phi_0)$ is cartesian, if $\Phi_0$ is a contraction contracting a
single loop $\{a_0(f),a_0(\ol{f})\}$ onto
$w_0$.

{\em Case II (Contracting a non-looping edge). } In this case $v_1\not=v_2$. We
require each $\Phi_i$ to contract
exactly one edge, namely $\{a_i(f),a_i(\ol{f})\}$ onto $w_0$. In particular,
this means that the only way the various
$(a_i,\tau_i,\Phi_i)$ differ is in the classes of $a_i(v_1)$ and $a_i(v_2)$. We
require that
$(\beta(a_i(v_1)),\beta(a_i(v_2)))_{i\in I}$ be a complete and non-repetitive
list of all pairs of elements of
$H_2(V)^+$ adding up to $\beta(w_0)$.

{\em Case III (Forgetting a tail). } Let us now deal with the case that
$\Phi:\tau\rightarrow\sigma$ stably forgets the
tail $f\in S\t$. We require $I$ to have one element, say $0$, and call
$(a_0,\tau_0,\Phi_0)$ cartesian if $\Phi_0$
stably forgets the tail $a_0(f)$ (and does nothing else).

{\em Case IV (Gluing two tails to an edge). } Finally, let us consider the case
that $\Phi$ is given by a combinatorial
morphism $c:\tau\rightarrow\sigma$, gluing the two tails $f$ and $\ol{f}$ of
$\tau$ to an edge $\{c(f),c(\ol{f})\}$ of
$\sigma$. Again, $I$ is required to have one element, say $0$, and
$(a_0,\tau_0,\Phi_0)$ is called cartesian if $\Phi_0$
glues two tails of $\tau_0$ to an edge of $\sigma'$ (and does nothing else).
Moreover, we require that $\Phi_0\comp
a_0=b\comp c$. An example:
\[\begin{array}{ccc}
\beginpicture%---tau_0
\setcoordinatesystem units <.3cm,.3cm> point at 3 2
\setplotarea x from 2 to 6, y from 0 to 4
\plot 2 3 4 3 /
\plot 4 1 3 1 /
\shaderectangleson
\setshadegrid span <1mm>
\putrectangle corners at 4 0 and 6 4
\put {\circle*{4}} [Bl] at 3 3
\axis left invisible label {$\scriptstyle\tau_0$} /
\axis right invisible label {\phantom{$\scriptstyle\tau_0$}} /
\axis bottom invisible label {\phantom{.}} /
\endpicture  & \stackrel{\Phi_0}{\longrightarrow} &
\beginpicture%---sigma'
\setcoordinatesystem units <.3cm,.3cm> point at 3 2
\setplotarea x from 3 to 6, y from 0 to 4
\circulararc 180 degrees from 4 3 center at 4 2
\shaderectangleson
\setshadegrid span <1mm>
\putrectangle corners at 4 0 and 6 4
\put {\circle*{4}} [Bl] at 3 2
\axis left invisible label {\phantom{$\scriptstyle\sigma'$}} /
\axis right invisible label {{$\scriptstyle\sigma'$}} /
\axis bottom invisible label {\phantom{.}} /
\endpicture  \\
\ldiagup{a_0} & & \rdiagup{b} \\
\beginpicture%---tau
\setcoordinatesystem units <.3cm,.3cm> point at 3 2
\setplotarea x from 2 to 6, y from 0 to 4
\plot 4 1 3 1 /
\plot 3 3 4 3 /
\shaderectangleson
\setshadegrid span <1mm>
\putrectangle corners at 4 0 and 6 4
\axis left invisible label {{$\scriptstyle\tau$}} /
\axis right invisible label {\phantom{$\scriptstyle\tau$}} /
\axis top invisible label {\phantom{.}} /
\endpicture & \stackrel{c}{\longrightarrow} &
\beginpicture%---sigma
\setcoordinatesystem units <.3cm,.3cm> point at 3 2
\setplotarea x from 3 to 6, y from 0 to 4
\circulararc 180 degrees from 4 3 center at 4 2
\shaderectangleson
\setshadegrid span <1mm>
\putrectangle corners at 4 0 and 6 4
\axis left invisible label {\phantom{$\scriptstyle\sigma$}} /
\axis right invisible label {{$\scriptstyle\sigma$}} /
\axis top invisible label {\phantom{.}} /
\endpicture
\end{array} \]

Note that in each case pullbacks exist, even though they are not necessarily
unique, even up to isomorphism, in the last
two cases. Note also, that for each $i\in I$ we have
$\deg(\tau_i)=\deg(\sigma')$.

We shall now define still another category, denoted $\tilde{\GG}_s(V)_{\cart}$,
called the {\em cartesian extended
isogeny category }over $V$.

\begin{defn} \label{ceicv}
Objects of $\tilde{\GG}_s(V)_\cart$ are pairs $(\tau,(a_i,\tau_i)_{i\in I})$,
where $\tau$ is a stable modular graph, $I$
is a finite set and for each $i\in I$ the pair $(a_i,\tau_i)$ is a stable
$V$-graph $\tau_i$, together with a
combinatorial morphism $a_i:\tau\rightarrow\tau_i$, identifying $\tau$ as the
absolute stabilization of $\tau_i$.

A {\em premorphism } from $(\tau,(a_i,\tau_i)_{i\in I})$ to
$(\sigma,(b_j,\sigma_j)_{j\in J})$ is a triple
$(\Phi,\lambda,(\Phi_i)_{i\in I})$, where $\Phi:\tau\rightarrow\sigma$ is an
extended isogeny of stable modular graphs,
$\lambda:I\rightarrow J$ is a map and for each $i\in I$ we have an extended
isogeny of stable $V$-graphs
$\Phi_i:\tau_i\rightarrow\sigma_{\lambda(i)}$. It is clear how to compose such
premorphisms and that composition is
associative.

A {\em morphism }in $\tilde{\GG}_s(V)_\cart$ is defined to be a premorphism
which can be factored into a composition of
elementary morphisms and isomorphisms. It is clear what an isomorphism is. An
{\em elementary morphism }is a premorphism
$(\Phi,\lambda,(\Phi_i)_{i\in I})$, as above, such that
\begin{enumerate}
\item $\Phi$ is an elementary extended isogeny,
\item for each $j\in J$ we have that
$(a_i,\tau_i,\Phi_i)_{i\in\lambda^{-1}(j)}$ is cartesian in the sense defined
in
Cases~I through~IV, above.
\end{enumerate}
\end{defn}

\begin{numrmk} \label{cfnc}
Projecting onto the first component defines a functor
\[\tilde{\GG}_s(V)_\cart\longrightarrow\tilde{\GG}_s(0).\]
Despite the notation, this is not a fibration of categories.
\end{numrmk}

We shall, in what follows, often shorten the notation $(\tau,(a_i,\tau_i)_{i\in
I})$ to $(\tau,(\tau_i)_{i\in I})$ or
even $(\tau_i)_{i\in I}$.

Call an object $(\tau_i)_{i\in I}$ of $\tilde{\GG}_s(V)_\cart$ {\em homogeneous
} of degree $n\in\zz$, if for all $i\in
I$ we have $\deg(V,\tau_i)=n$.

For a stable modular graph $\tau$, we may consider the fiber
$\tilde{\GG}_s(V)_{\cart/\tau}$ of the functor
$\tilde{\GG}_s(V)_\cart\rightarrow\tilde{\GG}_s(0)$ over $\tau$. In every such
fiber $\tilde{\GG}_s(V)_{\cart/\tau}$ we
have a functor
%% FOLLOWING LINE CANNOT BE BROKEN BEFORE 80 CHAR
\[\oplus:\tilde{\GG}_s(V)_{\cart/\tau}\times\tilde{\GG}_s(V)_{\cart/\tau}
\longrightarrow\tilde{\GG}_s(V)_{\cart/\tau},\]
given by
\[(\tau_i)_{i\in I}\oplus(\sigma_j)_{j\in J}=((\tau_i)_{i\in
I},(\sigma_j)_{j\in J}),\]
where we think of the object on the right hand side as a family parametrized by
$I\amalg J$. The functor $\oplus$
satisfies some obvious properties, which we shall not list.

Every object $X=(\tau_i)_{i\in I}$ of $\tilde{\GG}_s(V)_\cart$ has a unique
decomposition $X=\bigoplus_{n\in\zz}X_n$
into homogeneous components. Every morphism in $\tilde{\GG}_s(V)_\cart$
respects this decomposition.

Finally, $\tilde{\GG}_s(V)_\cart$ is a tensor category (in the sense of
\cite{delmil}) with tensor product given by
%% FOLLOWING LINE CANNOT BE BROKEN BEFORE 80 CHAR
\[\otimes:\tilde{\GG}_s(V)_{\cart}\times\tilde{\GG}_s(V)_{\cart}
\longrightarrow \tilde{\GG}_s(V)_{\cart},\]
which is defined by the formula
\[(\tau,(\tau_i)_{i\in I})\otimes(\sigma,(\sigma_j)_{j\in
J})=(\tau\times\sigma,(\tau_i\times\sigma_j)_{(i,j)\in I\times
J}).\]
For two graphs $\sigma$ and $\tau$ we denote by $\sigma\times\tau$ the graph
whose geometric realization is the
disjoint union of $|\sigma|$ and $|\tau|$. This notion extends in an obvious
way to marked graphs. The identity object
for $\otimes$ is the one element family with value the empty graph.

There are obvious compatibilities between these various structures on
$\tilde{\GG}_s(V)_\cart$. For example, if
$X=\bigoplus_nX_n$ and $Y=\bigoplus_mY_m$ are objects of
$\tilde{\GG}_s(V)_\cart$, then the decomposition of $X\otimes
Y$ into homogeneous components is given by
\[X\otimes Y=\bigoplus_r\left(\bigoplus_{n+m=r}X_n\otimes Y_m\right).\]
We summarize these properties by saying that $\tilde{\GG}_s(V)_\cart$ has
$\oplus$, $\otimes$ and $\deg$ structures.

A formally similar situation arises, for example, if we consider the category
of morphisms of an additive tensor
category $\CC$ in which all homomorphism groups are graded.  If we denote this
morphism category by $\MM\CC$, there is a
functor $\MM\CC\rightarrow\CC\times\CC$, given by source and target, whose
fibers have a graded $\oplus$-structure as
above. Also, $\MM\CC$ becomes a tensor category compatible with $\deg$ and
$\oplus$. So $\MM\CC$ has $\oplus$, $\otimes$
and $\deg$ structures. In fact, Gromov-Witten invariants may be thought of as a
functor from $\tilde{\GG}_s(V)_\cart$ to
$\MM\CC$ respecting the $\oplus$, $\otimes$ and $\deg$ structures. In this case
$\CC$ will be a category of motives.

\begin{defn} \label{doasvg}
A full subcategory $\tilde{\TT}_s(A)\subset\tilde{\GG}_s(A)$ is called {\em
admissible}, if it satisfies the following
axioms.
\begin{enumerate}
\item If $\Phi:\sigma\rightarrow\tau$ is an extended isogeny in
$\tilde{\GG}_s(A)$ and $\tau\in \ob\tilde{\TT}_s(A)$,
then $\sigma\in\ob\tilde{\TT}_s(A)$.
\item If $\sigma$ and $\tau$ are in $\tilde{\TT}_s(A)$, then so is
$\sigma\times\tau$.
\end{enumerate}
\end{defn}

For an admissible subcategory $\tilde{\TT}_s(A)\subset\tilde{\GG}_s(A)$ and a
homomorphism $\xi:A\rightarrow B$, the
full subcategory $\tilde{\TT}_s(B)\subset\tilde{\GG}_s(B)$ of graphs which are
stabilizations of objects of
$\tilde{\TT}_s(A)$ is admissible.

For a smooth projective variety $V$ of pure dimension, we may construct the
full subcategory
$\tilde{\TT}_s(V)_\cart\subset\tilde{\GG}_s(V)_\cart$, called the {\em
associated cartesian category}, which may be
characterized as the subcategory of $\tilde{\GG}_s(V)_\cart$ such that for each
object $(\tau,(a_i,\tau_i)_{i\in I})$ we
have that $\tau\in\ob\tilde{\TT}_s(0)$ and for all $i\in I$ that
$\tau_i\in\ob\tilde{\TT}_s(V)$. Note that
$\tilde{\TT}_s(V)_\cart$ inherits the $\oplus$, $\otimes$ and $\deg$ structures
from $\tilde{\GG}_s(V)_\cart$.

\begin{examples}
{\em I}. Call a marked graph $\tau$ a {\em forest}, if
\begin{enumerate}
\item $H^1(|\tau|)=0$,
\item $g(v)=0$, for all $v\in V\t$.
\end{enumerate}
Let $\tilde{\TT}_s(A)\subset\tilde{\GG}_s(A)$ be the full subcategory whose
objects are forests. Then $\tilde{\TT}_s(A)$
is an admissible subcategory, called the {\em tree level }subcategory of
$\tilde{\GG}_s(A)$.

{\em II}. Let $\tilde{\TT}_s(A)\subset\tilde{\GG}_s(A)$ be an admissible
subcategory. Let $d:A\rightarrow\zz_{\geq0}$ be
an additive map and $N>0$ an integer. Then let $\tilde{\TT}_s(A)_{d<N}$ be the
full subcategory of $\tilde{\TT}_s(A)$
given by the condition
%% FOLLOWING LINE CANNOT BE BROKEN BEFORE 80 CHAR
\[\tau\in\ob\tilde{\TT}_s(A)_{d<N}\quad \Longleftrightarrow\quad
\mbox{$d(\beta(v))<N$,  for all $v\in V\t$}.\]
The subcategory $\tilde{\TT}_s(A)_{d<N}\subset\tilde{\GG}_s(A)$ is admissible.

If $A=H_2(V)^+$, then a very ample invertible sheaf $L$ on $V$ gives rise to
$d:H_2(V)^+\rightarrow\zz_{\geq0}$, by
setting $d(\beta)=\beta(L)$. If $\chr k\not=0$, we shall always pass to
$\tilde{\TT}_s(V)_{d<\chr k}$, in other words
assume that
\[\tilde{\TT}_s(V)=\tilde{\TT}_s(V)_{d<\chr k}.\]
But for emphasis, we may say that $\tilde{\TT}_s(V)$ is {\em bounded by the
characteristic}.
\end{examples}

\section{Orientations}

Fix a smooth projective variety $V$ of pure dimension.  Recall the following
five basic properties of $\ol{M}$.

{\em Property I (Mapping to a point). } Let $\tau$ be a stable $V$-graph of
class zero. Then $\tau$ is absolutely
stable. The evaluation morphism factors through $V^{\pi_0|\tau|}\subset
V^{P\t}$ and the canonical morphism
\[\ol{M}(V,\tau)\longrightarrow V^{\pi_0|\tau|}\times\ol{M}(\tau)\]
is an isomorphism. This follows immediately from Corollary~\ref{zc}. In
particular, $\ol{M}(V,\tau)$ is smooth.

Assume that $|\tau|$ is non-empty and connected. Let $(C,x)$ be the universal
family of stable marked curves over
$\ol{M}(\tau)$. Glue the $(C_v)_{v\in F\t}$ according to the edges of $\tau$ to
obtain a stable marked curve
$\pi:\tilde{C}\rightarrow\ol{M}(\tau)$ over $\ol{M}(\tau)$. Denote the vector
bundle of rank $g(\tau)\dim V$ on
$\ol{M}(V,\tau)$ given by $T_V\boxtimes R^1\pi\lst\O_{\tilde{C}}$ by
$\tT^{(1)}$.

{\em Property II (Products). } Let $\sigma$ and $\tau$ be stable $V$-graphs and
$\sigma\times\tau$ the obvious stable
$V$-graph whose geometric realization is the disjoint union of $|\sigma|$ and
$|\tau|$. There are obvious
combinatorial morphisms $\sigma\rightarrow\sigma\times\tau$ and
$\tau\rightarrow\sigma\times\tau$ giving rise to morphisms
of stable $V$-graphs $\sigma\times\tau\rightarrow\sigma$ and
$\sigma\times\tau\rightarrow\tau$ called the {\em
projections}. The induced morphism
%% FOLLOWING LINE CANNOT BE BROKEN BEFORE 80 CHAR
\[\ol{M}(V,\sigma\times\tau)\longrightarrow\ol{M}(V,\sigma)
\times\ol{M}(V,\tau)\]
is an isomorphism. This follows directly from the definitions.

{\em Property III (Cutting edges). } Let $\Phi:\sigma\rightarrow\tau$ be a
morphism of stable $V$-graphs of type cutting
an edge. So $\Phi$ is induced by a combinatorial morphism
$a:\tau\rightarrow\sigma$. Let $f$ and $\ol{f}$ be the tails
of $\tau$ that come from the edge of $\sigma$ which is being cut by $\Phi$. So
this edge is $\{a(f),a(\ol{f})\}$.
The diagram of algebraic $k$-stacks
\begin{equation} \label{ceedd}
%% FOLLOWING LINE CANNOT BE BROKEN BEFORE 80 CHAR
\comdia{\ol{M}(V,\sigma)}{\ev_{\{a(f),a(\ol{f})\}}}{V}{\ol{M}(\Phi)}{
}{\Delta}{\ol{M}(V, \tau)}{\ev_f\times
\ev_{\ol{f}}}{V\times V,}
\end{equation}
where the horizontal maps are evaluations at the indicated flags, is cartesian.
In particular, $\ol{M}(\Phi)$ is a
closed immersion. Again, this follows directly from the definitions.

{\em Property IV (Forgetting tails). } Let $\Phi:\sigma\rightarrow\tau$ be a
morphism of stable $V$-graphs stably
forgetting a tail. Denote the combinatorial morphism giving rise to $\Phi$ by
$a:\tau\rightarrow\sigma$ and the
forgotten tail by $f\in F\s$.

If $\Phi$ is of Type~I (i.e.\ incomplete), let $v=\del\s(f)$. Let
$\pi':C'\rightarrow\ol{M}(V,\sigma)$ be the universal
curve indexed by $v$ and $x:\ol{M}(V,\sigma)\rightarrow C'$ the universal
section given by $f$. Let
$\pi:C\rightarrow\ol{M}(V,\tau)$ be the universal curve indexed by the unique
vertex $w$ of $\tau$ such that
$a(w)=v$. Then by definition there is a commutative diagram
%% FOLLOWING LINE CANNOT BE BROKEN BEFORE 80 CHAR
\[\comdia{C'}{}{C}{\pi'}{}{\pi}{\ol{M}(V,\sigma)}{\ol{M}(\Phi)}{
\ol{M}(V,\tau),}\]
and the section $x$ induces an $\ol{M}(V,\tau)$-morphism
\[\ol{M}(V,\sigma)\rightarrow C.\]
This is an isomorphism. In particular, $\ol{M}(\Phi)$ is proper and flat of
relative dimension one. This follows from
Corollary~\ref{esuc}.

If $\Phi:\sigma\rightarrow\tau$ removes a tripod, then
\[\ol{M}(\Phi):\ol{M}(V,\sigma)\rightarrow\ol{M}(V,\tau)\]
is an isomorphism. This is because
$\ol{M}(\text{$0$-tripod})=\ol{M}_{0,3}=\spec k$.

{\em Property V (Isogenies). } Let
\[(\Phi,\lambda,(\Phi_i)_{i\in I}):(\tau,(a_i,\tau_i)_{i\in
I})\longrightarrow(\sigma,(b_j,\sigma_j)_{j\in J})\]
be a morphism in $\tilde{\GG}_s(V)_\cart$, where $\Phi$ (and hence all
$\Phi_i$) is an isogeny, i.e.\ free of any tail
gluing factors.  For each $j\in J$ we have a commutative diagram
\[\comdia{\displaystyle\coprod_{i\in
%% FOLLOWING LINE CANNOT BE BROKEN BEFORE 80 CHAR
I\atop\lambda(i)=j}\ol{M}(V,\tau_i)}{\amalg\ol{M}(\Phi_i)}{
\ol{M}(V,\sigma_j)}{\amalg
%% FOLLOWING LINE CANNOT BE BROKEN BEFORE 80 CHAR
\ol{M}(\ol{a}_i)}{}{\ol{M}(\ol{b})}{\ol{M}(\tau)}{\ol{M}(\Phi)}{
\ol{M}(\sigma).}\]
This diagram should be considered close to being cartesian. See
Definition~\ref{domb} for a more precise statement. For
the moment let us note that the induced morphism
\[\coprod_{i\in
%% FOLLOWING LINE CANNOT BE BROKEN BEFORE 80 CHAR
I\atop\lambda(i)=j}\ol{M}(V,\tau_i)\longrightarrow\ol{M}(\tau)
\times_{\ol{M}(\sigma)} \ol{M}(V,\sigma_j)\]
is surjective.

If $X$ is a separated algebraic Deligne-Mumford stack, by $A\lst(X)$ we shall
mean the rational Chow group of $X$ (see
\cite{vistoli}). If $X\rightarrow Y$ is a morphism of separated algebraic
Deligne-Mumford stacks, $A\upst(X\rightarrow
Y)$ will denote the rational bivariant intersection theory defined in
\cite{vistoli}.

\begin{defn} \label{domb}
Let $\tilde{\TT}_s(V)\subset\tilde{\GG}_s(V)$ be an admissible subcategory
(bounded by the characteristic). Let for each
$\tau\in\ob\tilde{\TT}_s(V)$ be given a cycle class
\[J(V,\tau)\in A_{\dim(V,\tau)}(\ol{M}(V,\tau)).\]
This collection of cycle classes is called an {\em orientation } of $\ol{M}$
over $\tilde{\TT}_s(V)$, if the following
axioms are satisfied.
\begin{enumerate}
\item \label{domb1} {\em (Mapping to a point). }We have
\[J(V,\tau)=c_{g(\tau)\dim V}(\tT^{(1)})\cdot[\ol{M}(V,\tau)],\]
for every stable $\tau\in\ob\tilde{\TT}_s(V)$ of class zero such that $|\tau|$
is non-empty and connected.
\item \label{domb2} {\em (Products). }In the situation of Property~II we have
\[J(V,\sigma\times\tau)=J(V,\sigma)\times J(V,\tau).\]
\item \label{domb3} {\em (Cutting edges). }In the situation of Property~III the
following is true. Let $[\ol{M}(\Phi)]\in
A^{\dim V}(\ol{M}(V,\sigma)\rightarrow\ol{M}(V,\tau))$ be the orientation class
of $\ol{M}(\Phi)$ obtained by pullback
(using Diagram~(\ref{ceedd})) from the canonical orientation $[\Delta]\in
A^{\dim V}(V\rightarrow V\times V)$. Then we
have
\[J(V,\sigma)=[\ol{M}(\Phi)]\cdot J(V,\tau).\]
In other words,
\[J(V,\sigma)=\Delta\upsh J(V,\tau),\]
where $\Delta\upsh$ is the Gysin homomorphism given by the complete
intersection morphism $\Delta$.
\item \label{domb4} {\em (Forgetting tails). }In the situation of Property~IV
the morphism $\ol{M}(\Phi)$ has a canonical
orientation $[\ol{M}(\Phi)]\in
A\upst(\ol{M}(V,\sigma)\rightarrow\ol{M}(V,\tau))$. We require that
\[J(V,\sigma)=[\ol{M}(\Phi)]\cdot J(V,\tau).\]
In other words,
\[J(V,\sigma)=\ol{M}(\Phi)\upst J(V,\tau),\]
where $\ol{M}(\Phi)\upst$ is given by flat pullback.
\item \label{domb5} {\em (Isogenies). }In the situation of Property~V, we have
for every $j\in J$ a class
\[\ol{M}(\Phi)\upsh J(V,\sigma_j)\in
A_{\dim(V,\sigma_j)}(\ol{M}(\tau)\times_{\ol{M}(\sigma)}\ol{M}(V,\sigma_j)),\]
since $\ol{M}(\Phi)$ has a canonical orientation, $\ol{M}(\tau)$ and
$\ol{M}(\sigma)$ being smooth of pure dimension. We
also have a morphism
%% FOLLOWING LINE CANNOT BE BROKEN BEFORE 80 CHAR
\[h:\coprod_{\lambda(i)=j}\ol{M}(V,\tau_i)\longrightarrow\ol{M}(
\tau)\times_{\ol{M}(
\sigma)}\ol{M}(V,\sigma_j)),\]
which is proper. The requirement is that
\[h\lst(\sum_{\lambda(i)=j}J(V,\tau_i))=\ol{M}(\Phi)\upsh J(V,\tau).\]
\end{enumerate}
\end{defn}

\begin{numrmk} \label{afmsl}
To check Axiom~(\ref{domb5}), it suffices to do so for $\Phi$ an elementary
isogeny, $\#J=1$ and
$(a_i,\tau_i,\Phi_i)_{i\in I}$ a pullback. This follows from the projection
formula.
\end{numrmk}

\begin{example}
If $\tau$ is a stable $V$-graph such that $|\tau|$ is non-empty and connected,
define
\[J_0(V,\tau)=
\begin{cases}
c_{g(\tau)\dim V}(\tT^{(1)})\cdot[\ol{M}(V,\tau)] & \text{if $\beta(\tau)=0$,}
\\
0                                                 & \text{otherwise.}
\end{cases}\]
For an arbitrary stable $V$-graph $\tau$, let
$\tau=\tau_1\times\ldots\times\tau_n$, for stable $V$-graphs
$\tau_1,\ldots,\tau_n$, such that $|\tau|=|\tau_1|\amalg\ldots\amalg|\tau_n|$
is the decomposition of $|\tau|$ into
connected components. Then set
\[J_0(V,\tau)=J_0(V,\tau_1)\times\ldots\times J_0(V,\tau_n).\]
We claim that $J_0$ is an orientation of $\ol{M}$ over $\tilde{\GG}_s(V)$,
called the {\em trivial orientation}.
\end{example}

\begin{defn}
Call a smooth projective variety $V$ {\em convex}, if for every morphism
$f:\pp^1\rightarrow V$ (defined over an
extension $K$ of $k$) we have $H^1(\pp^1,f\upst T_V)=0$.
\end{defn}

\noprint{
\begin{prop}
Let $\pi:C\rightarrow T$ be a prestable curve of genus zero and $f:C\rightarrow
V$ a morphism to a convex variety. Then
$R^1\pi\lst f\upst T_V=0$.
\end{prop}
\begin{pf}
\end{pf}
}

\begin{prop}
Let $V$ be convex and $\tau$ a stable $V$-forest. Then $\ol{M}(V,\tau)$ is
smooth of dimension $\dim(V,\tau)$. Moreover,
the morphism
\[\ol{M}(V,\tau)\longrightarrow\ol{M}(\tau^s)\]
is flat of relative dimension $\chi(\tau^s)\dim V-\deg(V,\tau)$.
\end{prop}
\begin{pf}
Let us start with some general remarks. Let $\tau$ be an absolutely stable
$V$-graph. Then we define
\[U(V,\tau)\subset\ol{M}(V,\tau)\]
to be the open substack of those stable maps $(C,x,f)$, such that
$(C_v,(x_i)_{i\in F\t(v)})$ is  a stable marked curve,
for all $v\in V\t$.  Let $(C,x):T\rightarrow\ol{M}(\tau)$ be  a $T$-valued
point of $\ol{M}(\tau)$, i.e.\
$(C_v,(x_i)_{i\in F\t(v)})_{v\in V\t}$ is a family of stable marked curves
parametrized by $T$. Let
$(\tilde{C},\tilde{x})$ be the stable marked curve over $T$ obtained by gluing
the $C_v$ according to the edges of
$\tau$. The diagram
\[\comdia{\Mor_T(\tilde{C},V_T)}{}{T}{}{}{}{U(V,\tau)}{}{\ol{M}(\tau)}\]
is cartesian. In particular, by Grothendieck \cite{fgaIV}, the morphism
$U(V,\tau)\rightarrow\ol{M}(\tau)$ is
representable, separated and of finite type. Moreover, let $(C,x,f)$ be a
$K$-valued point of $U(V,\tau)$. Let
$(\tilde{C},\tilde{x})$ be the marked curve obtained by gluing the $C_v$ and
$\tilde{f}:\tilde{C}\rightarrow V$
the morphism induced by the $f_v$. If $H^1(\tilde{C},\tilde{f}\upst T_V)=0$,
then $(C,x,f)$ is a smooth point of
$U(V,\tau)\rightarrow\ol{M}(\tau)$ and we have
\[T_{U(V,\tau)/\ol{M}(\tau)}(C,x,f)=H^0(\tilde{C},\tilde{f}\upst T_V)\]
for the relative tangent space. (This is the case, if $\tau$ is a $V$-forest
and $V$ is convex.)

In this smooth case we may calculate the relative dimension of $U(V,\tau)$ over
$\ol{M}(\tau)$ at $(C,x,f)$ as
\begin{eqnarray*}
\dim_K H^0(\tilde{C},\tilde{f}\upst T_V)  & = & \chi(\tilde{f}\upst T_V) \\
& = & \deg\tilde{f}\upst T_V + \rk(\tilde{f}\upst T_V)\chi(\tilde{C}) \\
& = & -\beta(\tau)(\omega_V) + \dim V \chi(\tau) \\
& = & \dim(V,\tau)-\dim(\tau).
\end{eqnarray*}
Since $\ol{M}(\tau)$ is smooth of dimension $\dim(\tau)$, we get that
$U(V,\tau)$ is smooth of dimension $\dim(V,\tau)$
at $(C,x,f)$.

Now let $\tau$ be an arbitrary stable $V$-graph. Then there exists an
absolutely stable $V$-graph $\tau'$, together with
a morphism $\tau'\rightarrow\tau$ of type forgetting tails, such that the
morphism
\[U(V,\tau')\longrightarrow\ol{M}(V,\tau)\]
is surjective, hence a flat epimorphism of relative dimension
$\#S_{\tau'}-\#S\t$. So if $U(V,\tau')$ is smooth of
dimension $\dim(V,\tau')$, then $\ol{M}(V,\tau)$ is smooth of dimension
\[\dim(V,\tau')-\#S_{\tau'}+\#S\t=\dim(V,\tau).\]
Finally, by considering the commutative diagram
%% FOLLOWING LINE CANNOT BE BROKEN BEFORE 80 CHAR
\[\comdia{U(V,\tau')}{}{\ol{M}(V,\tau)}{}{}{
}{\ol{M}(\tau')}{}{\ol{M}(\tau^s),}\]
we see that in this case $\ol{M}(V,\tau)\rightarrow\ol{M}(\tau^s)$ is flat of
relative dimension $\chi(\tau^s)\dim
V-\deg(V,\tau)$.
\end{pf}

\begin{them} \label{uosc}
Let $V$ be a convex variety and $\tilde{\TT}_s(V)\subset\tilde{\GG}_s(V)$ the
admissible subcategory of
$V$-forests bounded by the characteristic. Then the collection
\[J(V,\tau)=[\ol{M}(V,\tau)]\]
is an orientation of $\ol{M}$ over $\tilde{\TT}_s(V)$.
\end{them}
\begin{pf}
Let us check the axioms.

(1) {\em Mapping to a point. }  This follows from the fact that $g(\tau)=0$ and
hence
\[c_{g(\tau)\dim V}(\tT^{(1)})=c_0(0)=1.\]

(2) {\em Products. } In complete generality we have for smooth proper
Deligne-Mumford stacks $X$ and $Y$ that
\[[X\times Y]=[X]\times[Y]\]
in $A\lst(X\times Y)$.

(3) {\em Cutting edges. } Again we have a general fact to the following effect.
Consider the cartesian diagram of
separated Deligne-Mumford stacks
\[\comdia{X}{f}{V}{j}{}{i}{Y}{}{W,}\]
where $i$ and $j$ are regular embeddings such that for the normal bundles we
have
\[f\upst N_{V/W}=N_{X/Y}.\]
Then $i\upsh[Y]=[X]$. If all four participating stacks are smooth and $i$ and
$j$ are closed immersions of the same
codimension, then these conditions are automatically satisfied (see for example
Proposition~17.13.2 in \cite{ega4}).
Thus we may apply this fact in our case.

More generally, we have that $i\upsh[Y]=[X]$ if all participating stacks are
smooth and
\[\dim X+\dim W=\dim Y +\dim V.\]

(4) {\em Forgetting tails. } Again, there is a general fact that
$f\upsh[Y]=[X]$ if $f:X\rightarrow Y$ is a flat
morphism of smooth and proper Deligne-Mumford stacks.

(5) {\em Isogenies. } In accordance with Remark~\ref{afmsl} we assume that
$\Phi$ is an elementary isogeny, $\#J=1$ and
that $(a_i,\tau_i,\Phi_i)_{i\in I}$ is a pullback. There are five cases to
consider, according to what type of
elementary isogeny $\Phi$ is. We use notation as in the definition of pullback.

{\em Case I (Contracting a loop). } This case does not occur, since $\sigma$
and $\tau$ are forests.

{\em Case II (Contracting an edge). } We will start with some general remarks.
Let $\tau$ be a stable $V$-graph, and
$v_1,\ldots,v_n$ absolutely stable vertices of $\tau$, i.e.\ vertices $v$ such
that $2g(v)+|v|\geq3$. (To avoid
ill-defined notation we assume that $n\geq1$.) Let
\[U_{v_1,\ldots,v_n}(V,\tau)\subset\ol{M}(V,\tau)\]
be the open substack of all those stable maps $(C,x,f)\in\ol{M}(V,\tau)$ such
that $$(C_{v_{\nu}},(x_i)_{i\in
F\t(v_{\nu})})$$ is a stable marked curve, for all $\nu=1,\ldots,n$.

With this notation the diagram
\[\comdia{\displaystyle\coprod_{i\in I}
%% FOLLOWING LINE CANNOT BE BROKEN BEFORE 80 CHAR
U_{a_i(v_1),a_i(v_2)}(V,\tau_i)}{}{U_{b(v_0)}(V,\sigma')}{}{
}{}{\ol{M}(\tau)}{}{\ol{M}(\sigma)}\]
is cartesian. Consider for a fixed $i\in I$ the open immersion
\[U_{a_i(v_1),a_i(v_2)}(V,\tau_i)\subset\ol{M}(V,\tau_i).\]
Let
\[Z_{a_i(v_1),a_i(v_2)}(V,\tau_i)\subset\ol{M}(V,\tau_i)\]
be the closed complement. We have
\[\dim Z_{a_i(v_1),a_i(v_2)}(V,\tau_i)<\dim \ol{M}(V,\tau_i).\]
Thus, to prove the equality of two cycles of degree $\dim(V,\tau_i)$ in
$\ol{M}(\tau)\times_{\ol{M}(\sigma)}\ol{M}(V,\sigma')$, it suffices to prove
the equality of the cycles restricted to
$\coprod_i U_{a_i(v_1),a_i(v_2)}(V,\tau_i)$. This reduces us to proving that
%% FOLLOWING LINE CANNOT BE BROKEN BEFORE 80 CHAR
\[\ol{M}(\Phi)\upsh[U_{b(v_0)}(V,\sigma')]=\sum_i[U_{a_i(v_1),
a_i(v_2)}(V,\tau_i)].\]
This claim finally follows from the general fact already mentioned in the proof
of Axiom~(3).

{\em Case III (Forgetting a tail, incompletely). }
Let $f\in F\t$ be the forgotten flag, $v=\del_{\tau_0}(a_0(f))$ and $w\in
V_{\sigma'}$ the vertex of $\sigma'$
corresponding to $v$ via $\Phi_0$. We have an open immersion
\[U_v(V,\tau_0)\subset \ol{M}(V,\tau_0)\]
with closed complement
\[Z_v(V,\tau_0)\subset \ol{M}(V,\tau_0)\]
of strictly smaller dimension. Thus, as in the previous case, we may reduce to
proving that
\[\ol{M}(\Phi)\upsh[U_w(V,\sigma')]=[U_v(V,\tau_0)].\]
This follows from the fact that the diagram
%% FOLLOWING LINE CANNOT BE BROKEN BEFORE 80 CHAR
\[\comdia{U_v(V,\tau_0)}{}{U_w(V,\sigma')}{}{}{
}{\ol{M}(\tau)}{\ol{M}(\Phi)}{\ol{M}(\sigma)}\]
is cartesian.

{\em Cases IV and V (Removing a tripod). } These cases are trivial, since
$\ol{M}(\Phi_0)$ and $\ol{M}(\Phi)$ are
isomorphisms.
\end{pf}

\noprint{
\begin{conjecture} \label{gmc}
Let $V$ be a Grassmannian variety. Then the collection
\[J(V,\tau)=[\ol{M}(V,\tau)]^0\]
is an orientation of $\ol{M}$ over $\tilde{\GG}_s(V)$.
\end{conjecture}
\begin{pf}
\end{pf}
}

\section{Deligne-Mumford-Chow Motives}

We shall imitate the usual construction of the category of Chow motives, as
described for example in \cite{scholl}.

Fix a ground field $k$. Let $\WW$ be the category of smooth and proper
algebraic Deligne-Mumford stacks over $k$. For an
object $X$ of $\WW$, let $A\upst(X)$ be the rational Chow ring of $X$ defined
by Vistoli \cite{vistoli}. Then $A\upst$
is a generalized cohomology theory with coefficient field $\qq$ in the sense of
\cite{kleiman}. Moreover, it is a graded
global intersection theory with Poincar\'e duality and cycle map in the
terminology of \cite{kleiman}.

If $X$ and $Y$ are objects of $\WW$ we define $S^d(Y,X)$, the group of {\em
correspondences }from $Y$ to $X$ of degree
$d$, to be
\[S^d(Y,X)=A^{n+d}(Y\times X),\]
if $Y$ is purely $n$-dimensional and
\[S^d(Y,X)=\bigoplus_i S^d(Y_i,X),\]
if $Y=\coprod_i Y_i$ is the decomposition of $Y$ into irreducible components.
Note that $S^d(Y,X)\subset A\upst(Y\times
X)$. The isomorphism $Y\times X\cong X\times Y$ exchanging components induces
an isomorphism
\[S^d(Y,X)\cong S^{d+n-m}(X,Y),\]
if $\dim Y=n$ and $\dim X=m$. We call this isomorphism {\em transpose }of
correspondences.
For objects $Z$, $Y$ and $X$ of $\WW$ we define composition of correspondences
by the usual formula
\[g\comp f={p_{13}}\lst(p_{12}\upst f\cdot p_{23}\upst g),\]
for $f\in S^d(Z,Y)$ and $g\in S^e(Y,X)$. Then $g\comp f\in S^{d+e}(Z,X)$.

The category $\ol{\WW}$ of {\em Deligne-Mumford-Chow motives } (or DMC-motives)
is now defined to be the category of
triples $(X,p,n)$, where $X\in \ob\WW$, $p\in S^0(X,X)$ such that $p^2=p$ and
$n\in\zz$. Homomorphisms are defined by
\[\Hom_{\ol{\WW}}((Y,q,m),(X,p,n))=p S^{n-m}(Y,X) q.\]
Note that $\Hom_{\ol{\WW}}((Y,q,m),(X,p,n))\subset S^{n-m}(Y,X)$. Composition
of homomorphisms in $\ol{\WW}$ is defined
as composition of correspondences.

There is a contravariant involution $\ol{\WW}\rightarrow\ol{\WW}$, denoted
$M\mapsto M^{\vee}$, defined by
$(X,p,n)^{\vee}=(X,^tp,\dim X-n)$, where $^tp$ is the transpose of $p$, on
objects and by transpose of correspondences on
homomorphisms.

\begin{prop}
The category $\ol{\WW}$ is a $\qq$-linear pseudo-abelian category. \qed
\end{prop}

Every morphism $f:X\rightarrow Y$ in $\WW$ defines a correspondence of degree
zero $\ol{f}\in S^0(Y,X)$ by
\[\ol{f}={\Gamma_f}\lst[X]\in A\upst(Y\times X),\]
where $\Gamma_f:X\rightarrow Y\times X$ is the graph of $f$.
We define the contravariant functor $h:\WW\rightarrow\ol{\WW}$ by
$h(X)=(X,\ol{\id}_X,0)$ and $h(f)=\ol{f}$. We usually
write $f\upst$ for $h(f)$ and $f\lst$ for $h(f)^{\vee}$.

Let $\ll=(\spec k,\ol{\id},-1)$ be the {\em Lefschetz motive}.  We shall use
the notation
\[M(n)=M\otimes\ll^{-n}.\]
We set
\[\Hom_{\ol{\WW}}^i(M,N)=\Hom_{\ol{\WW}}(M\otimes\ll^i,N)\]
and
\[\Hom_{\ol{\WW}}\upst(M,N)=\bigoplus_{i\in\zz}\Hom_{\ol{\WW}}^i(M,N).\]
The category with the same objects as $\ol{\WW}$, but with homomorphism groups
given by $\Hom_{\ol{\WW}}\upst(M,N)$ will be called the
category of {\em graded }DMC-motives.

For a DMC-motive $M$, define
\[A^i(M)=\Hom(\ll^i,M)\]
and
\[A\upst(M)=\bigoplus_iA^i(M).\]

\begin{prop}[Identity principle]
If $f,g:M\rightarrow N$ are two homomorphisms of DMC-motives, such that the
induced homomorphisms
\[A\upst(M\otimes h(X))\longrightarrow A\upst(N\otimes h(X))\]
agree, for all $X\in\ob\WW$, then $f=g$. \qed
\end{prop}

Let $\ol{\VV}$ be the category of Chow motives (which is defined as $\ol{\WW}$
is above, but starting with $\VV$ instead
of $\WW$). There is a natural fully faithful functor
$\ol{\VV}\rightarrow\ol{\WW}$.

\begin{question}
Is the functor $\ol{\VV}\rightarrow\ol{\WW}$ an equivalence of categories?
\end{question}

Let $H$ be a graded generalized cohomology theory on $\WW$ with a coefficient
field $\Lambda$ of characteristic zero,
possessing a cycle map such that $\pp^1$ satisfies epu (see \cite{kleiman}).
Then $H$ induces a covariant functor
(called a {\em realization functor\/})
\[\ol{H}:(\rtext{graded DMC-motives})\longrightarrow(\rtext{graded
$\Lambda$-algebras}),\]
such that for $X\in\ob\WW$ we have $\ol{H}(h(X))=H(X)$ and for a correspondence
$\xi\in S^d(Y,X)$ we have an induced
homomorphism
\begin{eqnarray*}
\ol{H}(\xi):H(Y) & \longrightarrow & H(X) \\
\alpha           & \longmapsto     &
{p_X}\lst({p_Y}\upst(\alpha)\cup\cl_{Y\times X}(\xi)).
\end{eqnarray*}
The functor $\ol{H}$ doubles the degree of a homomorphism.

The following are examples of such a cohomology theory $H$.
\begin{enumerate}
\item If $k=\cc$, consider to $X$ the associated topological stack $X^{\topo}$.
This is a stack on the category of
topological spaces with the \'etale topology. It has an associated \'etale
topos $X^{\topo}_{\et}$. Set
\[H_B(X)=H\upst(X^{\topo}_{\et},\qq)\]
and call it the {\em Betti cohomology }of $X$. Here $\Lambda=\qq$.
\item If $\ell\not=\chr k$ set
\[H_{\ell}(X)=H\upst(\ol{X}_{\et},\qql)={\displaystyle\projlim\limits_n}
H\upst(\ol{X}_{\et},\zz/\ell^n),\]
where $\ol{X}=X\times_{\spec k}\spec \ol{k}$ is the lift of $X$ to an algebraic
closure of $k$ and $\ol{X}_{\et}$
denotes the \'etale topos of $\ol{X}$. We call $H_{\ell}(X)$ the {\em
$\ell$-adic cohomology } of $X$. In this case
$\Lambda=\qql$.
\item If $\chr k=0$, let $\Omega_X\com$ be the algebraic deRham complex of $X$
and set
\[H_{dR}(X)=\hh\upst(X,\Omega_X\com).\]
We call $H_{dR}(X)$ the {\em algebraic deRham cohomology }of $X$. Here
$\Lambda=k$.
\end{enumerate}

\section{Motivic Gromov-Witten Classes}

Define the contravariant tensor functor
\[h(\ol{M}):\tilde{\GG}_s(0)\longrightarrow(\rtext{DMC-motives})\]
by $h(\ol{M})(\tau)=h(\ol{M}(\tau))$ on objects. For a morphism
$(a,\sigma',\Phi):\sigma\rightarrow\tau$ we have
$\ol{M}(\ol{a}):\ol{M}(\sigma')\rightarrow\ol{M}(\sigma)$ and
$\ol{M}(\Phi):\ol{M}(\sigma')\rightarrow\ol{M}(\tau)$.
Then let
\[h(\ol{M})(a,\sigma',\Phi)=\ol{M}(\ol{a})\lst\comp\ol{M}(\Phi)\upst.\]
This makes sense, because $\ol{M}(\ol{a})\lst$ is of degree zero,
$\ol{M}(\ol{a})$ being an isomorphism. This is also
why $h(\ol{M})$ is functorial.

Now fix a smooth projective variety $V$ of pure dimension and consider the
contravariant tensor functor
\[h(V)^{\otimes S}{(\chi\dim
V)}:\tilde{\GG}_s(0)\longrightarrow(\rtext{DMC-motives})\]
defined on objects by
\[\tau\longmapsto h(V)^{\otimes S\t}({\chi(\tau)\dim V}).\]
For a morphism $(a,\sigma',\Phi):\sigma\rightarrow\tau$ let $E$ be the set of
edges of $\sigma'$ which are cut by
$a:\sigma\rightarrow\sigma'$. Then we have
$V^{S\s}=V^{S_{\sigma'}}\times(V\times V)^E$. Let $p:V^{S_{\sigma'}}\times
V^E\rightarrow V^{S_{\sigma'}}$ be the projection,
$\Delta:V^{S_{\sigma'}}\times V^E\rightarrow V^{S_{\sigma'}}\times
(V\times V)^E=V^{S\s}$ the identity times the $E$-fold power of the diagonal.
Finally, we have an injection
$\Phi^S:S\t\rightarrow S_{\sigma'}$ giving rise to
$\Phi^S:V^{S_{\sigma'}}\rightarrow V^{S\t}$. We define the
homomorphism
\[h(V)^{\otimes S\t}{(\chi(\tau)\dim V)}\longrightarrow h(V)^{\otimes
S\s}{(\chi(\sigma)\dim V)}\]
as the composition of the three homomorphisms
\[(\Phi^S)\upst:h(V)^{\otimes S\t}{(\chi(\tau)\dim V)}\longrightarrow
h(V)^{\otimes
S_{\sigma'}}{(\chi(\sigma')\dim V)},\]
\[p\upst:h(V)^{\otimes S_{\sigma'}}{(\chi(\sigma')\dim V)}\longrightarrow
h(V)^{\otimes
S_{\sigma'}\cup E}{(\chi(\sigma')\dim V)}\]
and
\[\Delta\lst:h(V)^{\otimes S_{\sigma'}\cup E}{(\chi(\sigma')\dim
V)}\longrightarrow h(V)^{\otimes
S\s}{(\chi(\sigma)\dim V)},\] noting that $\chi(\tau)=\chi(\sigma')$ and
$\chi(\sigma')=\chi(\sigma)-\#
E$. Functoriality is a straightforward check using the identity principle.

Pulling back $h(\ol{M})$ and $h(V)^{\otimes S}{(\chi\dim V)}$ to the cartesian
extended isogeny category over $V$ via the
functor of Remark~\ref{cfnc}, we get two contravariant tensor functors
\[\tilde{\GG}_s(V)_\cart\longrightarrow(\text{graded DMC-motives}).\]

Now let $\tilde{\TT}_s(V)\subset\tilde{\GG}_s(V)$ be an admissible subcategory
(bounded by the characteristic) and $J$
an orientation of $\ol{M}$ over $\tilde{\TT}_s(V)$. For every object $\tau$ of
$\tilde{\TT}_s(V)$ we have a morphism
\[\phi_{(V,\tau)}:\ol{M}(V,\tau)\longrightarrow V^{S_{\tau^s}}\times
\ol{M}(\tau^s).\]
The first component is given by evaluation, noting that we have a map
$F_{\tau^s}\rightarrow F\t$.
Then
\begin{eqnarray*}
{\phi_{(V,\tau)}}\lst J(V,\tau) & \in &
S^{\dim(\tau^s)-\dim(V,\tau)}(V^{S_{\tau^s}},\ol{M}(\tau^s)) \\
 & = & \Hom_{\ol{\WW}}^{\deg(V,\tau)}(h(V^{S_{\tau^s}}){(\chi(\tau^s)\dim
V)},h(\ol{M}(\tau^s))).
\end{eqnarray*}

\begin{defn}
Define
\[I(V,\tau)={\phi_{(V,\tau)}}\lst J(V,\tau),\]
so that we have a homomorphism
\[I(V,\tau):h(V)^{\otimes S_{\tau^s}}{(\chi(\tau^s)\dim V)}\longrightarrow
h(\ol{M}(\tau^s)){(\deg(V,\tau))}\]
of DMC-motives over $k$. We call $I$ the system of {\em Gromov-Witten classes
}associated to the orientation $J$.
\end{defn}

Restricting the two functors $h(\ol{M})$ and $h(V)^{\otimes S}{(\chi\dim V)}$
to $\tilde{\TT}_s(V)_\cart$, we
get two contravariant tensor functors
\[\tilde{\TT}_s(V)_\cart\longrightarrow(\text{graded DMC-motives}).\]
We shall now define a natural transformation
\[I:h(V)^{\otimes S}{(\chi\dim V)}\longrightarrow h(\ol{M}).\]
So let $(\tau,(\tau_i)_{i\in I})$ be an object of $\tilde{\TT}_s(V)_\cart$, and
define
\[I(\tau,(\tau_i)_{i\in I})=\sum_{i\in I}I(V,\tau_i):h(V)^{\otimes
S_{\tau}}{(\chi(\tau)\dim V)}\longrightarrow
h(\ol{M}(\tau)).\]

\begin{them} \label{vagwc}
The Gromov-Witten transformation $I$ is a natural transformation compatible
with the $\oplus$, $\otimes$ and $\deg$
structures. Moreover,
\begin{enumerate}
\item \label{vag1} {\em (Mapping to a point). } The triangle
\[\comtri{h(V)^{\otimes S\t}{(\chi(\tau)\dim
V)}}{\rtext{mult}}{h(V){(\chi(\tau)\dim V)}}{I(V,\tau)}{c_{g(\tau)\dim
V}(\tT^{(1)})}{h(\ol{M}(\tau))}\]
commutes, for any stable $V$-graph $\tau$  of class zero in
$\tilde{\TT}_s(V)$, such that $|\tau|$ is non-empty and
connected.
\item \label{vag2} {\em (Divisor). }  Let $\lL\in\Pic(V)$ be a line bundle, so
its Chern class induces a homomorphism
$c_1(\lL):\ll\rightarrow h(V)$. Let $\Phi:\sigma\rightarrow\tau$ be a morphism
in $\tilde{\TT}_s(V)$ of type forgetting a tail,
such that the corresponding vertex of $\tau$ is absolutely stable. Then the
square
\[\begin{array}{ccc}
h(V)^{\otimes S_{\sigma^s}}{(\chi(\sigma^s)\dim V)} &
\stackrel{I(V,\sigma)}{\longrightarrow} &
h(\ol{M}(\sigma^s)){(\deg(V,\sigma))}\\
\ldiagup{c_1(\lL)} & & \rdiag{\ol{M}(\Phi)\lst} \\
h(V)^{\otimes S_{\tau^s}}{(\chi(\tau^s)\dim V)} \otimes\ll &
\stackrel{\beta(\lL)I(V,\tau)}{\longrightarrow} &
h(\ol{M}(\tau^s)){(\deg(V,\tau))}\otimes\ll
\end{array}\]
commutes.
\end{enumerate}
\end{them}
\begin{rmk}
To make this statement more precise, consider to $(\text{graded DMC-motives})$
the associated category of morphisms.
Then the natural transformation $I$ may be considered as a functor
\[I:\tilde{\TT}_s(V)_\cart\longrightarrow(\text{graded morphisms of
DMC-motives}).\]
Both categories have $\oplus$, $\otimes$ and $\deg$ structures and $I$
preserves them.  This essentially means that
\begin{enumerate}
\item $I((\tau_i)\oplus(\sigma_j))=I((\tau_i))+I((\sigma_j))$,
\item $\deg I((\tau_i))=\deg(\tau_i)$, if $(\tau_i)$ is homogeneous,
\item $I((\tau,\tau_i)\otimes(\sigma,\sigma_j))=I(\tau,\tau_i)\otimes
I(\sigma,\sigma_j)$.
\end{enumerate}
\end{rmk}
\begin{pf}
All this follows formally from Definition~\ref{domb} using the identity
principle and the bivariant formalism (as
explained for example in \cite{fulmacp}).
\end{pf}

\begin{rmks}
\begin{enumerate}
\item
Applying Theorem~\ref{uosc} we get the tree level system of Gromov-Witten
invariants for convex varieties.
\item
By applying a realization functor, we get Betti, $\ell$-adic or deRham
Gromov-Witten classes.
\item
Theorem~\ref{vagwc} implies all the axioms for Gromov-Witten classes listed in
\cite{KM}. Perhaps only Formula~(2.7) is
not quite evident. In view of its importance (it implies that the fundamental
class remains the identity with respect to
quantum multiplication), we will show that it follows from the rest of the
axioms.  In fact, assume that
\begin{equation} \label{leman}
\langle I_{0,3,\beta}\rangle (\gamma_1\otimes\gamma_2\otimes e^0) \ne 0.
\end{equation}

Choose a divisorial class $\delta$ with nonvanishing intersection with $\beta$.
In view of the Divisor Axiom, we must
then have

$$\langle I_{0,4,\beta}\rangle (\gamma_1\otimes\gamma_2\otimes\delta\otimes
e^0) \ne 0.$$

In view of (2.6), the last class is the lift of

$$\langle I_{0,3,\beta}\rangle (\gamma_1\otimes\gamma_2\otimes\delta ).$$

But this cannot be non-vanishing simultaneously with~(\ref{leman}) because the
Grading Axiom does not allow this.

More generally, this argument shows that whenever $e^0$ is among the arguments,
then $\langle I\rangle = 0$ for $\beta
\ne 0$, any genus, any $n.$ Geometrically: `if one of the points on $C$ is
unconstrained, the problem cannot have
finitely many (and non-zero) solutions.'

\noprint{The composition
\[h(V)^{\otimes2}{-\dim V}\stackrel{p\upst}{\longrightarrow}
h(V)^{\otimes3}{-\dim
V}\stackrel{I_{0,3}(V,\beta)}{\longrightarrow}\ll^{\beta(\omega_V)}\]
is zero, if $\beta\not=0$. Here $p:V^3\rightarrow V^2$ is a projection.}

\end{enumerate}
\end{rmks}

\begin{thebibliography}{10}

\bibitem{delmil}
P.~Deligne and J.~Milne.
\newblock Tannakian categories.
\newblock In {\em Hodge Cycles, Motives and Shimura Varieties}, Lecture Notes
  in Mathematics Vol.\ 900, pages 101--228. Springer, Berlin, Heidelberg, New
  York, 1982.

\bibitem{fulmacp}
W.~Fulton and R.~MacPherson.
\newblock {\em Categorical framework for the study of singular spaces}.
\newblock Memoirs of the American Mathematical Society No.\ 243. American
  Mathematical Society, Providence, 1981.

\bibitem{fgaIV}
A.~Grothendieck.
\newblock Techniques de construction et th\'eor\`emes d'existence en
  g\'eom\'etrie alg\'ebrique {IV}: Les sch\'emas de {H}ilbert.
\newblock {\em S\'eminaire Bourbaki}, 13e ann\'ee(221), 1960--61.

\bibitem{ega2}
A.~Grothendieck.
\newblock {\em Elements de G\'eom\'etrie Alg\'ebrique II: Etude globale
  \'el\'ementaire de quelques classes de morphismes, {\em\bf EGA2}}.
\newblock Publications Math\'ematiques No.\ 8. Institut des Hautes \'Etudes
  Scientifiques, Bois-Marie Bures-sur-Yvette, 1961.

\bibitem{ega4}
A.~Grothendieck.
\newblock {\em Elements de G\'eom\'etrie Alg\'ebrique IV: Etude Locale des
  Sch\'emas et des Morphismes de Sch\'emas, {\em\bf EGA4}}.
\newblock Publications Math\'ematiques Nos.\ 20, 24, 28, 32. Institut des
  Hautes \'Etudes Scientifiques, Bois-Marie Bures-sur-Yvette, 1964, 65, 66, 67.

\bibitem{sga1}
A.~Grothendieck et~al.
\newblock {\em Rev\^etements \'Etales et Groupe Fondamental {\em\bf SGA1}}.
\newblock Lecture Notes in Mathematics No.\ 224. Springer, Berlin, Heidelberg,
  New York, 1971.

\bibitem{KapMan}
M.~Kapranov and Yu. Manin.
\newblock Paper in preparation.

\bibitem{kleiman}
S.~Kleiman.
\newblock Motives.
\newblock In F.~Oort, editor, {\em Algebraic Geometry, Oslo 1970}, pages
  53--96. Wolters-Noordhoff, 1972.

\bibitem{knudsen}
F.~Knudsen.
\newblock The projectivity of the moduli space of stable curves, {I}{I}: The
  stacks ${M}_{g,n}$.
\newblock {\em Math.\ Scand.}, 52:161--199, 1983.

\bibitem{K}
M.~Kontsevich.
\newblock Enumeration of rational curves via torus actions.
\newblock Preprint, Max-Planck-Institut f\"ur Mathematik, Bonn, 1994.

\bibitem{KM}
M.~Kontsevich and Yu. Manin.
\newblock {G}romov-{W}itten classes, quantum cohomology, and enumerative
  geometry.
\newblock {\em Commuications in Mathematical Physics}, 164:525--562, 1994.

\bibitem{KM2}
M.~Kontsevich and Yu. Manin.
\newblock Quantum cohomology of a product.
\newblock Preprint, Max-Planck-Institut f\"ur Mathematik, Bonn, 1995.

\bibitem{pandh}
R.~Pandharipande.
\newblock Notes on {K}ontsevich's compactification of the space of maps.
\newblock Preprint, preliminary version, March 7, 1995.

\bibitem{scholl}
A.~J. Scholl.
\newblock Classical motives.
\newblock In {\em Motives}, Proceedings of Symposia in Pure Mathematics Vol.\
  55, Part~I, pages 163--187. American Mathematical Society, Providence, 1994.

\bibitem{vistoli}
A.~Vistoli.
\newblock Intersection theory on algebraic stacks and on their moduli spaces.
\newblock {\em Inventiones mathematicae}, 97:613--670, 1989.

\end{thebibliography}

\end{document}
